% Équations, inéquations et systèmes — Mathématiques 1ère A — Cours signé OnBuch+
% Programme officiel MINESEC. Compiler avec : tectonic NA01-equations-inequations-systemes-a.tex
\newcommand{\DOCMATIERE}{Mathématiques}
\newcommand{\DOCNIVEAU}{1ère A}
\newcommand{\DOCMODULE}{Relations et opérations fondamentales dans l'ensemble des nombres réels}
\newcommand{\DOCLECON}{NA01}
\newcommand{\DOCTITRE}{Équations, inéquations et systèmes}
% OnBuch+ — préambule « cours signés » ENRICHI (compilé avec tectonic / XeLaTeX)
% Système visuel « Pop » d'OnBuch+ : fond crème (jamais blanc), bordures encre
% épaisses, ombres « dures » décalées, titres Archivo Black en capitales.
% Version MATHÉMATIQUES : graphes (pgfplots/tikz), boîtes pédagogiques
% (définition, propriété, méthode, attention, le savais-tu, exercice résolu,
% exercice, TP, bilan), page de garde et signature OnBuch+ ancrée en bas.
%
% Macros attendues dans main.tex AVANT \input{preamble} :
%   \DOCMATIERE  \DOCNIVEAU  \DOCMODULE  \DOCLECON (numéro)  \DOCTITRE
\documentclass[11pt]{article}

\usepackage{fontspec}
\usepackage[french]{babel}
\usepackage[a4paper,margin=2cm,top=3.4cm,bottom=2.8cm,headheight=1.4cm,footskip=1.3cm]{geometry}
\usepackage{amsmath,amssymb,amsfonts}
\usepackage{mathtools} % \xmapsto, \coloneqq... souvent utilisés par les modèles
\usepackage{cancel} % simplifications barrées
% \abs{z} (module, valeur absolue) et \norm{u} (norme), utilisés par les modèles
\providecommand{\abs}{}\let\abs\relax\DeclarePairedDelimiter{\abs}{\lvert}{\rvert}
\providecommand{\norm}{}\let\norm\relax\DeclarePairedDelimiter{\norm}{\lVert}{\rVert}
\usepackage[table]{xcolor} % [table] : fournit aussi \rowcolors (alternance de lignes)
\usepackage{pagecolor}
\usepackage{tikz}
\usetikzlibrary{arrows.meta,positioning,calc,shapes.geometric,shapes.misc,shapes.arrows,decorations.pathmorphing,decorations.pathreplacing,decorations.markings,patterns,shadows,mindmap,trees,fit,backgrounds,matrix,intersections,angles,quotes}
\usepackage{pgfplots}
\usepackage{pgf-pie} % diagrammes circulaires \pie (statistiques)
\pgfplotsset{compat=1.18}
\usepgfplotslibrary{fillbetween,groupplots} % aires entre courbes (calcul intégral)
\usepackage{enumitem}
\usepackage{fancyhdr}
% [most] charge toutes les bibliothèques tcolorbox (listings, minted, raster,
% documentation...) et gonfle inutilement la mémoire TeX sur un document long
% (~300 boîtes) : on ne charge que ce qui est réellement utilisé.
\usepackage[skins,breakable]{tcolorbox}
\usepackage{graphicx}
\usepackage{colortbl}
\usepackage{booktabs}
\usepackage{tabularx}
\usepackage{diagbox} % case d'en-tête diagonale des tableaux à double entrée
% Colonnes extensibles centrée / gauche / droite (C, L, R), souvent utilisées par les modèles
\newcolumntype{C}{>{\centering\arraybackslash}X}
\newcolumntype{L}{>{\raggedright\arraybackslash}X}
\newcolumntype{R}{>{\raggedleft\arraybackslash}X}
\usepackage{array}
\usepackage{multicol}
\usepackage{siunitx}
% parse-numbers=false : accepte aussi \SI{2,5 \times 10^{-3}}{...}
\sisetup{locale=FR,output-decimal-marker={,},group-minimum-digits=4,parse-numbers=false}
\DeclareSIUnit\ohm{\ensuremath{\symup{\Omega}}} % U+2126 absent de la police
\usepackage{qrcode}
% \degree : commande fréquemment utilisée par les modèles pour un angle ou une
% température en dehors de \SI{}{} (non fournie par les paquets chargés).
\providecommand{\degree}{^{\circ}}
% \qed (fin de démonstration), utilisé par les modèles sans amsthm
\providecommand{\qed}{\hfill$\square$}
% \barre : double barre des valeurs interdites dans un tableau de variations (tkz-tab)
\providecommand{\barre}{\ensuremath{\|}}
% environnement proof (amsthm non chargé) utilisé par les modèles
\ifdefined\proof\else\newenvironment{proof}[1][Démonstration]{\par\noindent\textit{#1.}\ }{\qed\par}\fi
\usepackage{lastpage}
\usepackage{hyperref}
\hypersetup{hidelinks}

% ============================================================
% Palette « Pop » OnBuch+ — mêmes hex que la classe P (plus_ui.dart)
% ============================================================
\definecolor{poporange}{HTML}{F59321}
\definecolor{popdark}{HTML}{E07A0C}
\definecolor{popcream}{HTML}{FFF6E8}
\definecolor{popink}{HTML}{1C1714}
\definecolor{poppurple}{HTML}{7A5AE0}
\definecolor{popgreen}{HTML}{2E9E6A}
\definecolor{poppink}{HTML}{E0457B}
\definecolor{popblue}{HTML}{2F7DE1}
\definecolor{popgold}{HTML}{C98A12}
\definecolor{popmuted}{HTML}{8A7F76}
\definecolor{popline}{HTML}{EADFD2}
\definecolor{popgray}{HTML}{8A7F76}
\colorlet{popgrey}{popgray}
% Alias tolérants (le modèle nomme parfois les couleurs différemment)
\colorlet{popwhite}{white}
\colorlet{popblack}{popink}
\colorlet{popred}{poppink}
\colorlet{popyellow}{popgold}
% Teintes claires pour fonds de tableaux / zones
\colorlet{popblueL}{popblue!10!white}
\colorlet{poporangeL}{poporange!14!white}
\colorlet{popgreenL}{popgreen!12!white}
\colorlet{poppurpleL}{poppurple!12!white}
\colorlet{poppinkL}{poppink!12!white}
\colorlet{popgoldL}{popgold!14!white}

\pagecolor{popcream}

% ============================================================
% Typographie
% ============================================================
\defaultfontfeatures{Ligatures=TeX}
\setmainfont{PlusJakartaSans-Medium.ttf}[Path=../../../fonts/,BoldFont=PlusJakartaSans-Bold.ttf]
% Maths en sans-serif (Fira Math) avec chiffres et lettres droites de Plus
% Jakarta, pour que les formules se fondent dans le texte.
\usepackage[math-style=ISO,bold-style=ISO]{unicode-math}
\setmathfont{FiraMath-Regular.otf}[Path=../../../fonts/]
\setmathfont{PlusJakartaSans-Medium.ttf}[Path=../../../fonts/,range={up/num,up/{Latin,latin},\mathup}]
\usepackage{icomma} % virgule décimale sans espace en mode math
% Caractères mathématiques Unicode absents des polices de texte (Plus Jakarta,
% Archivo Black) : rendus par leur équivalent mathématique, en texte comme en
% formule — sinon ils s'affichent en carré vide (ex. « Arithmétique dans ℤ »).
\usepackage{newunicodechar}
\newunicodechar{ℝ}{\ensuremath{\mathbb{R}}}
\newunicodechar{ℤ}{\ensuremath{\mathbb{Z}}}
\newunicodechar{ℂ}{\ensuremath{\mathbb{C}}}
\newunicodechar{ℕ}{\ensuremath{\mathbb{N}}}
\newunicodechar{ℚ}{\ensuremath{\mathbb{Q}}}
\newunicodechar{𝔻}{\ensuremath{\mathbb{D}}}
\newunicodechar{α}{\ensuremath{\alpha}}
\newunicodechar{β}{\ensuremath{\beta}}
\newunicodechar{γ}{\ensuremath{\gamma}}
\newunicodechar{δ}{\ensuremath{\delta}}
\newunicodechar{ε}{\ensuremath{\varepsilon}}
\newunicodechar{θ}{\ensuremath{\theta}}
\newunicodechar{λ}{\ensuremath{\lambda}}
\newunicodechar{μ}{\ensuremath{\mu}}
\newunicodechar{σ}{\ensuremath{\sigma}}
\newunicodechar{τ}{\ensuremath{\tau}}
\newunicodechar{φ}{\ensuremath{\varphi}}
\newunicodechar{ψ}{\ensuremath{\psi}}
\newunicodechar{ω}{\ensuremath{\omega}}
\newunicodechar{Σ}{\ensuremath{\Sigma}}
% majuscules produites par \MakeUppercase dans les titres (α-aminés -> Α)
\newunicodechar{Α}{\ensuremath{\alpha}}
\newunicodechar{Β}{\ensuremath{\beta}}
\newunicodechar{Γ}{\ensuremath{\Gamma}}
\newunicodechar{Δ}{\ensuremath{\Delta}}
\newunicodechar{Ε}{\ensuremath{\varepsilon}}
\newunicodechar{Θ}{\ensuremath{\Theta}}
\newunicodechar{Λ}{\ensuremath{\Lambda}}
\newunicodechar{Μ}{\ensuremath{\mu}}
\newunicodechar{Τ}{\ensuremath{\tau}}
\newunicodechar{Φ}{\ensuremath{\Phi}}
\newunicodechar{Ψ}{\ensuremath{\Psi}}
\newunicodechar{Ω}{\ensuremath{\Omega}}
\newunicodechar{↦}{\ensuremath{\mapsto}}
\newunicodechar{∈}{\ensuremath{\in}}
\newunicodechar{∉}{\ensuremath{\notin}}
\newunicodechar{∧}{\ensuremath{\wedge}}
\newunicodechar{⇔}{\ensuremath{\Leftrightarrow}}
\newunicodechar{⟺}{\ensuremath{\Longleftrightarrow}}
\newunicodechar{⇒}{\ensuremath{\Rightarrow}}
\newunicodechar{⟹}{\ensuremath{\Longrightarrow}}
\newunicodechar{∘}{\ensuremath{\circ}}
\newunicodechar{∪}{\ensuremath{\cup}}
\newunicodechar{∩}{\ensuremath{\cap}}
\newunicodechar{≡}{\ensuremath{\equiv}}
\newunicodechar{⊂}{\ensuremath{\subset}}
\newunicodechar{⊥}{\ensuremath{\perp}}
\newunicodechar{∃}{\ensuremath{\exists}}
\newunicodechar{∀}{\ensuremath{\forall}}
\newunicodechar{∛}{\ensuremath{\sqrt[3]{\,}}}
\newunicodechar{∜}{\ensuremath{\sqrt[4]{\,}}}
\newunicodechar{⃗}{\ensuremath{{}^{\to}}}
\newunicodechar{ⁿ}{\ifmmode^{n}\else\textsuperscript{\ensuremath{n}}\fi}
\newunicodechar{ˣ}{\ifmmode^{x}\else\textsuperscript{\ensuremath{x}}\fi}
\newunicodechar{ᵇ}{\ifmmode^{b}\else\textsuperscript{\ensuremath{b}}\fi}
\newunicodechar{ʳ}{\ifmmode^{r}\else\textsuperscript{\ensuremath{r}}\fi}
\newunicodechar{ᵘ}{\ifmmode^{u}\else\textsuperscript{\ensuremath{u}}\fi}
\newunicodechar{ʸ}{\ifmmode^{y}\else\textsuperscript{\ensuremath{y}}\fi}
\newunicodechar{ᵃ}{\ifmmode^{a}\else\textsuperscript{\ensuremath{a}}\fi}
\newunicodechar{ᵉ}{\ifmmode^{e}\else\textsuperscript{\ensuremath{e}}\fi}
\newunicodechar{ᵏ}{\ifmmode^{k}\else\textsuperscript{\ensuremath{k}}\fi}
\newunicodechar{ᵖ}{\ifmmode^{p}\else\textsuperscript{\ensuremath{p}}\fi}
\newunicodechar{ᴺ}{\ifmmode^{N}\else\textsuperscript{\ensuremath{N}}\fi}
\newunicodechar{ᶜ}{\ifmmode^{c}\else\textsuperscript{\ensuremath{c}}\fi}
\newunicodechar{ʰ}{\ifmmode^{h}\else\textsuperscript{\ensuremath{h}}\fi}
\newunicodechar{ᵠ}{\ifmmode^{q}\else\textsuperscript{\ensuremath{q}}\fi}
\newunicodechar{ₙ}{\ifmmode_{n}\else\textsubscript{\ensuremath{n}}\fi}
\newunicodechar{ₐ}{\ifmmode_{a}\else\textsubscript{\ensuremath{a}}\fi}
\newunicodechar{ᵢ}{\ifmmode_{i}\else\textsubscript{\ensuremath{i}}\fi}
\newunicodechar{ᵣ}{\ifmmode_{r}\else\textsubscript{\ensuremath{r}}\fi}
\newunicodechar{ₖ}{\ifmmode_{k}\else\textsubscript{\ensuremath{k}}\fi}
\newunicodechar{ᵦ}{\ifmmode_{\beta}\else\textsubscript{\ensuremath{\beta}}\fi}
\newunicodechar{ₓ}{\ifmmode_{x}\else\textsubscript{\ensuremath{x}}\fi}
\newfontfamily\popdisplay{ArchivoBlack-Regular.ttf}[Path=../../../fonts/]
\newfontfamily\poplogo{SpaceGrotesk-Bold.ttf}[Path=../../../fonts/]
\newfontfamily\popsemi{PlusJakartaSans-SemiBold.ttf}[Path=../../../fonts/]
\newfontfamily\popxbold{PlusJakartaSans-ExtraBold.ttf}[Path=../../../fonts/]
\newfontfamily\popxboldls{PlusJakartaSans-ExtraBold.ttf}[Path=../../../fonts/,LetterSpace=3.0]

\setlist[enumerate]{leftmargin=1.7em,itemsep=5pt,topsep=4pt}
\setlist[itemize]{leftmargin=1.7em,itemsep=4pt,topsep=4pt,label={\color{poporange}\textbullet}}
\setlength{\parindent}{0pt}
\setlength{\parskip}{5pt}
\renewcommand{\arraystretch}{1.25}

% pgfplots : style Pop par défaut
\pgfplotsset{
  every axis/.append style={axis line style={popink,line width=1pt},tick style={popink},
    grid style={popline,line width=0.5pt},label style={font=\small},tick label style={font=\footnotesize}},
  popcurve/.style={poporange,line width=1.8pt,no markers,smooth},
}
% Même style disponible dans un \draw TikZ simple (hors pgfplots)
\tikzset{popcurve/.style={poporange,line width=1.8pt}}

% ============================================================
% Pill / squiggle
% ============================================================
\newcommand{\poppill}[3]{%
  \tikz[baseline=-0.55ex]\node[draw=popink,line width=1.2pt,rounded corners=2.6pt,%
    fill=#2,inner xsep=6.5pt,inner ysep=3pt]{%
    \popxboldls\fontsize{7}{7}\selectfont\color{#3}\MakeUppercase{#1}};%
}
\newcommand{\popsquiggle}[1]{%
  \tikz[baseline=-0.4ex]{
    \draw[#1,line width=2.6pt,line cap=round]
      plot [smooth] coordinates {(0,0.09) (0.34,0.01) (0.68,0.21) (1.02,0.01) (1.36,0.10)};
  }%
}
% Mot-clé surligné (marqueur orange sous le texte)
\newcommand{\cle}[1]{{\popxbold\color{popdark}#1}}

% ============================================================
% En-tête / pied
% ============================================================
\pagestyle{fancy}
\fancyhf{}
\renewcommand{\headrulewidth}{0pt}
\renewcommand{\footrulewidth}{0pt}
\lhead{\raisebox{-0.15ex}{\poplogo\fontsize{13}{13}\selectfont\color{popink}OnBuch\color{poporange}+}}
\rhead{\poppill{\DOCMATIERE\ \textbullet\ \DOCNIVEAU\ \textbullet\ Leçon \DOCLECON}{white}{popink}}
\lfoot{\popxbold\fontsize{7.6}{7.6}\selectfont\color{popmuted}\MakeUppercase{Cours signé OnBuch+}\ \textbullet\ {\popsemi\DOCTITRE}}
\rfoot{\tikz[baseline=-0.6ex]\node[rounded corners=8.5pt,fill=popink,minimum height=17pt,minimum width=17pt,inner xsep=5pt,inner sep=0]{\popxbold\fontsize{8}{8}\selectfont\color{popcream}\thepage\,/\,\pageref*{LastPage}};}

% ============================================================
% Titres
% ============================================================
\newcommand{\chaptitre}[2]{%
  \begin{tcolorbox}[enhanced,colback=poporange,colframe=popink,boxrule=2.6pt,arc=11pt,
      drop shadow=popink,left=15pt,right=15pt,top=12pt,bottom=11pt,before skip=3pt,after skip=18pt]
    \poppill{#2}{popink}{popcream}\par\vspace{9pt}
    {\raggedright\hyphenpenalty=10000\exhyphenpenalty=10000\popdisplay\fontsize{22}{24}\selectfont\color{popcream}\MakeUppercase{#1}\par}
  \end{tcolorbox}
}
% Section numérotée : pastille orange avec le numéro + titre Archivo
\newcounter{popsec}
\newcommand{\coursec}[1]{%
  \stepcounter{popsec}\setcounter{popsub}{0}%
  \par\vspace{16pt}\noindent
  \begin{minipage}{\linewidth}
  \tikz[baseline=-0.7ex]\node[draw=popink,line width=1.6pt,fill=poporange,rounded corners=4pt,minimum size=22pt,inner sep=0]{\popdisplay\fontsize{12}{12}\selectfont\color{popcream}\thepopsec};%
  \hspace{8pt}{\popdisplay\fontsize{15}{17}\selectfont\color{popink}\MakeUppercase{#1}}\par
  \vspace{4pt}\noindent{\color{poporange}\rule{36pt}{3.6pt}}
  \end{minipage}\par\nopagebreak\vspace{8pt}
}
\newcounter{popsub}
\newcommand{\courssub}[1]{%
  \stepcounter{popsub}%
  \par\vspace{9pt}\noindent{\popxbold\fontsize{11.5}{13}\selectfont\color{popink}{\color{poporange}\thepopsec.\thepopsub}\hspace{6pt}#1}\par\nopagebreak\vspace{4pt}
}

% ============================================================
% Boîtes pédagogiques — même recette : fond blanc, bordure épaisse colorée,
% ombre dure encre, badge plein. Titre optionnel [#1].
% ============================================================
\tcbset{popbase/.style={breakable,enhanced,colback=white,boxrule=2.3pt,arc=9pt,
  shadow={3pt}{-3pt}{0pt}{popink},left=14pt,right=14pt,top=11pt,bottom=11pt,before skip=11pt,after skip=15pt,
  fontupper=\popsemi\fontsize{10.6}{14.5}\selectfont\color{popink},
  fontlower=\popsemi\fontsize{10.6}{14.5}\selectfont\color{popink}}}
% Coche dessinée (le glyphe ✓ n'existe pas dans les polices du document)
\newcommand{\popcheck}[1]{\tikz[baseline=-0.5ex]\draw[#1,line width=1.6pt,line cap=round,line join=round](-0.13,0.0)--(-0.04,-0.09)--(0.13,0.1);}
\newcommand{\popboxhead}[3]{\poppill{#1}{#2}{white}\ifx\relax#3\relax\else\hspace{6pt}{\popxbold\fontsize{10.6}{12}\selectfont\color{popink}#3}\fi\par\vspace{7pt}}

\newtcolorbox{aretenir}[1][]{popbase,colframe=popgreen,before upper={\popboxhead{À retenir}{popgreen}{#1}}}
\newtcolorbox{exemplebox}[1][]{popbase,colframe=poporange,before upper={\popboxhead{Exemple}{poporange}{#1}}}
\newtcolorbox{definition}[1][]{popbase,colframe=popblue,before upper={\popboxhead{Définition}{popblue}{#1}}}
\newtcolorbox{propriete}[1][]{popbase,colframe=poppurple,before upper={\popboxhead{Propriété}{poppurple}{#1}}}
\newtcolorbox{methode}[1][]{popbase,colframe=popgold,before upper={\popboxhead{Méthode}{popgold}{#1}}}
\newtcolorbox{attention}[1][]{popbase,colframe=poppink,before upper={\popboxhead{Attention}{poppink}{#1}}}
\newtcolorbox{experience}[1][]{popbase,colframe=popblue,colback=popblueL,before upper={\popboxhead{Expérience}{popblue}{#1}}}
% « Le savais-tu ? » : carte violette pleine (contexte, culture, Cameroun)
\newtcolorbox{savaistu}[1][]{popbase,colback=poppurple,colframe=popink,
  fontupper=\popsemi\fontsize{10.4}{14.2}\selectfont\color{white},
  before upper={\poppill{Le savais-tu\,?}{white}{poppurple}\ifx\relax#1\relax\else\hspace{6pt}{\popxbold\color{white}#1}\fi\par\vspace{7pt}}}
% Exercice résolu : énoncé en haut, \tcblower puis la solution sur fond crème
\newcounter{popexo}\newcounter{popexr}
\newtcolorbox{exoresolu}[1][]{popbase,colframe=poporange,colbacklower=popcream,
  segmentation style={popink,line width=1.2pt,dashed},
  before upper={\stepcounter{popexr}\popboxhead{Exercice résolu \thepopexr}{poporange}{#1}},
  before lower={\poppill{Solution}{popink}{popcream}\par\vspace{6pt}}}
% Exercice (non résolu en place) — la correction est dans la partie Corrigés
\newtcolorbox{exercice}[1][]{popbase,colframe=popink,
  before upper={\stepcounter{popexo}\popboxhead{Exercice \thepopexo}{popink}{#1}}}
% Corrigé
\newtcolorbox{corrige}[1][]{popbase,colframe=popgreen,colback=popgreenL,
  before upper={\popboxhead{Corrigé}{popgreen}{#1}}}
% Objectifs / prérequis / bilan
\newtcolorbox{objectifs}{popbase,colframe=popink,colback=poporangeL,
  before upper={\popboxhead{Objectifs de la leçon}{popink}{}}}
\newtcolorbox{prerequis}{popbase,colframe=popink,colback=popblueL,
  before upper={\popboxhead{Prérequis}{popblue}{}}}
\newtcolorbox{bilan}{popbase,colframe=popink,colback=white,boxrule=2.8pt,
  before upper={\popboxhead{Fiche bilan}{poporange}{L'essentiel à connaître pour l'examen}}}
% Figure encadrée avec légende
\newtcolorbox{popfigure}[1][]{enhanced,breakable=false,colback=white,colframe=popline,boxrule=1.4pt,arc=8pt,
  left=10pt,right=10pt,top=10pt,bottom=8pt,before skip=10pt,after skip=12pt,halign=center,
  lower separated=false,halign lower=center,
  fontlower=\popsemi\fontsize{9}{11.5}\selectfont\color{popmuted},
  #1}
\newcommand{\legende}[1]{\par\vspace{6pt}{\popsemi\fontsize{9}{11.5}\selectfont\color{popmuted}#1}}

% ============================================================
% Signature OnBuch+
% ============================================================
\newcommand{\poplogoinline}[1]{{\poplogo\fontsize{#1}{#1}\selectfont\color{popink}OnBuch\color{poporange}+}}
\newcommand{\signatureonbuch}{%
  \begin{tcolorbox}[enhanced,colback=popink,colframe=popink,boxrule=2.2pt,arc=10pt,
      drop shadow=poporange,left=14pt,right=14pt,top=10pt,bottom=10pt,before skip=0pt,after skip=0pt]
    \tikz[baseline=-0.7ex]\node[circle,fill=popgreen,draw=popcream,line width=1.3pt,minimum size=22pt,inner sep=0]{\popcheck{white}};%
    \hspace{9pt}\begin{minipage}[c]{0.62\linewidth}
      {\popxbold\fontsize{12}{13}\selectfont\color{popcream}Signé\ \textbullet\ {\poplogo OnBuch\color{poporange}+}}\par\vspace{2pt}
      {\raggedright\popsemi\fontsize{8}{10.5}\selectfont\color{popline}\MakeUppercase{Équipe pédagogique OnBuch+}\\\MakeUppercase{Conforme au programme officiel MINESEC}\par}
    \end{minipage}\hfill
    {\poplogo\fontsize{22}{22}\selectfont\color{popcream}OnBuch\color{poporange}+}
  \end{tcolorbox}%
}

% ============================================================
% Page de garde
% #1 titre · #2 matière · #3 niveau · #4 module · #5 n° leçon · #6 durée ·
% #7 description / objectifs (texte ou itemize)
% ============================================================
\newcommand{\pagegarde}[7]{%
  \thispagestyle{empty}
  \newgeometry{a4paper,margin=1.7cm,top=1.6cm,bottom=1.4cm}%
  \begin{tikzpicture}[remember picture,overlay]
    \fill[poporange!18] ([xshift=-3.2cm,yshift=-3.6cm]current page.north east) circle (4.6cm);
    \fill[poppurple!14] ([xshift=2.4cm,yshift=7.6cm]current page.south west) circle (3.8cm);
  \end{tikzpicture}%
  {\poplogo\fontsize{30}{30}\selectfont\color{popink}OnBuch\color{poporange}+}\hfill
  \raisebox{0.6ex}{\poppill{Cours signé \textbullet\ Premium}{popink}{popcream}}\par
  \vspace{1pt}\popsquiggle{poppurple}\par
  \vspace{8pt}
  \begin{tcolorbox}[enhanced,colback=poporange,colframe=popink,boxrule=2.8pt,arc=13pt,
      drop shadow=popink,left=18pt,right=18pt,top=12pt,bottom=13pt,after skip=10pt]
    \poppill{#2 \textbullet\ #3}{popink}{popcream}\hspace{5pt}\poppill{#4}{white}{popink}\par\vspace{8pt}
    {\popdisplay\fontsize{14}{14}\selectfont\color{popink}LEÇON #5}\par\vspace{4pt}
    {\raggedright\hyphenpenalty=10000\exhyphenpenalty=10000\popdisplay\fontsize{25}{27}\selectfont\color{popcream}\MakeUppercase{#1}\par}
    \vspace{8pt}
    {\popxbold\fontsize{9.5}{9.5}\selectfont\color{popink}\MakeUppercase{Durée conseillée : #6}}
  \end{tcolorbox}
  \begin{tcolorbox}[enhanced,colback=white,colframe=popink,boxrule=2.2pt,arc=10pt,
      drop shadow=popink,left=16pt,right=16pt,top=10pt,bottom=10pt,after skip=8pt]
    \poppill{Dans cette leçon}{poppurple}{white}\par\vspace{5pt}
    \popsemi\fontsize{9.3}{12.6}\selectfont\color{popink}#7
  \end{tcolorbox}
  \noindent\begin{minipage}[t]{0.487\linewidth}
    \begin{tcolorbox}[enhanced,colback=popgreen,colframe=popink,boxrule=2.2pt,arc=10pt,
        drop shadow=popink,left=11pt,right=11pt,top=10pt,bottom=10pt,height=3.1cm,valign=center]
      \tcbox[colback=white,colframe=popink,boxrule=1.3pt,arc=5pt,boxsep=0pt,nobeforeafter,
        left=3pt,right=3pt,top=3pt,bottom=3pt,valign=center]{\qrcode[height=1.9cm]{https://play.google.com/store/apps/details?id=com.luvvix.onbuch&hl=fr}}%
      \hspace{8pt}\begin{minipage}[c]{0.5\linewidth}
        {\popdisplay\fontsize{11}{12.5}\selectfont\color{white}\MakeUppercase{Télécharge OnBuch}}\par\vspace{4pt}
        {\raggedright\popsemi\fontsize{8.4}{11}\selectfont\color{white}Gratuit \textbullet\ Google Play\\Tuteur IA, annales, concours, résultats\par}
      \end{minipage}
    \end{tcolorbox}
  \end{minipage}\hfill
  \begin{minipage}[t]{0.487\linewidth}
    \begin{tcolorbox}[enhanced,colback=poppurple,colframe=popink,boxrule=2.2pt,arc=10pt,
        drop shadow=popink,left=11pt,right=11pt,top=10pt,bottom=10pt,height=3.1cm,valign=center]
      \tikz[baseline=-0.7ex]\node[circle,fill=white,draw=popink,line width=1.3pt,minimum size=1.6cm,inner sep=0]{\popdisplay\fontsize{24}{24}\selectfont\color{poppurple}+};%
      \hspace{8pt}\begin{minipage}[c]{0.6\linewidth}
        {\popdisplay\fontsize{11}{12.5}\selectfont\color{white}\MakeUppercase{Passe à OnBuch+}}\par\vspace{4pt}
        {\raggedright\popsemi\fontsize{8.4}{11}\selectfont\color{white}Tous les cours signés, Tuteur IA illimité, fiches PDF — onglet OnBuch+ de l'app\par}
      \end{minipage}
    \end{tcolorbox}
  \end{minipage}
  \vspace{8pt}
  \popcontenu
  \vfill
  \signatureonbuch
  \restoregeometry
  \newpage
}

% Rangée « Ce cours contient » (page de garde)
\newcommand{\poptile}[3]{% #1 couleur #2 gros texte #3 légende
  \begin{tcolorbox}[enhanced,colback=#1,colframe=popink,boxrule=2pt,arc=9pt,shadow={2.5pt}{-2.5pt}{0pt}{popink},
      left=6pt,right=6pt,top=9pt,bottom=9pt,halign=center,valign=center,height=1.75cm,width=\linewidth,nobeforeafter]
    {\popdisplay\fontsize{12.5}{13}\selectfont\color{white}\MakeUppercase{#2}}\par\vspace{3pt}
    {\centering\popsemi\fontsize{7.8}{9.5}\selectfont\color{white}#3\par}
  \end{tcolorbox}}
\newcommand{\popcontenu}{%
  \par\noindent{\popxboldls\fontsize{7.5}{7.5}\selectfont\color{popmuted}CE COURS SIGNÉ CONTIENT}\par\vspace{7pt}
  \noindent\begin{minipage}[t]{0.235\linewidth}\poptile{popblue}{Cours}{complet, illustré, pas à pas}\end{minipage}\hfill
  \begin{minipage}[t]{0.235\linewidth}\poptile{poporange}{Méthodes}{et exercices résolus}\end{minipage}\hfill
  \begin{minipage}[t]{0.235\linewidth}\poptile{poppink}{Exercices}{type examen + corrigés}\end{minipage}\hfill
  \begin{minipage}[t]{0.235\linewidth}\poptile{popgreen}{Fiche bilan}{l'essentiel à retenir}\end{minipage}\par}

% Sommaire
\newtcolorbox{sommaire}{enhanced,colback=white,colframe=popink,boxrule=2.2pt,arc=10pt,
  drop shadow=popink,left=18pt,right=18pt,top=13pt,bottom=13pt,before skip=6pt,after skip=20pt,
  before upper={\poppill{Sommaire}{poppurple}{white}\par\vspace{9pt}\popsemi\fontsize{10.6}{14.5}\selectfont\color{popink}}}

% Signature de fin de cours — poussée tout en bas de la dernière page
\newcommand{\findecours}{%
  \par\vfill\nopagebreak
  \begin{center}\popsquiggle{poporange}\end{center}\vspace{4pt}
  \signatureonbuch
}
\AtBeginDocument{\def\llbracket{\mathopen{[\mkern-3.5mu[}}\def\rrbracket{\mathclose{]\mkern-3.5mu]}}} % intervalles d'entiers (glyphe ⟦ absent de la police)
% Glyphes absents de Fira Math : redéfinis à partir de glyphes disponibles
\AtBeginDocument{%
  \def\setminus{\mathbin{\backslash}}\let\smallsetminus\setminus
  \def\vdots{\mathord{\vbox{\baselineskip3.2pt\lineskiplimit0pt\kern4pt\hbox{.}\hbox{.}\hbox{.}}}}%
  \def\ddots{\mathinner{\mkern1mu\raise6.5pt\hbox{.}\mkern2mu\raise3.6pt\hbox{.}\mkern2mu\raise0.7pt\hbox{.}\mkern1mu}}%
  \def\triangle{\mathord{\scalebox{0.9}{$\symup{\Delta}$}}}%
  \def\nabla{\mathord{\raisebox{\depth}{\scalebox{1}[-1]{$\symup{\Delta}$}}}}%
  \def\checkmark{\popcheck{popdark}}%
  \def\star{\mathbin{\ast}}\def\bigstar{\mathord{\ast}}%
  \def\diamond{\mathbin{\scalebox{0.6}{$\Diamond$}}}%
  \def\bigcup{\mathop{\mathchoice{\raisebox{-0.3em}{\scalebox{1.7}{$\cup$}}}{\scalebox{1.25}{$\cup$}}{\cup}{\cup}}\displaylimits}%
  \def\bigcap{\mathop{\mathchoice{\raisebox{-0.3em}{\scalebox{1.7}{$\cap$}}}{\scalebox{1.25}{$\cap$}}{\cap}{\cap}}\displaylimits}%
  \def\lesseqqgtr{\mathrel{\lessgtr}}\def\bigoplus{\mathop{\oplus}\displaylimits}%
}
\providecommand{\ssi}{si et seulement si} % abréviation fréquente
\AtBeginDocument{\providecommand{\wideparen}{\overparen}} % arc de cercle

%% FIGFIX-BEGIN v4 — correctifs globaux des figures (docs/onbuch-generation/tools/figfix)
\usepackage{adjustbox}\usepackage{etoolbox}
% toute figure TikZ/pgfplots plus large que la ligne est réduite à la largeur disponible
\BeforeBeginEnvironment{tikzpicture}{\begin{adjustbox}{max width=\linewidth}}
\AfterEndEnvironment{tikzpicture}{\end{adjustbox}}
% légendes pgfplots lisibles par-dessus les courbes
\pgfplotsset{every axis/.append style={legend style={fill=white,fill opacity=0.92,text opacity=1,draw=black!15,inner sep=3pt}}}
% étiquettes d'arêtes (syntaxe edge["..."]) avec fond blanc
\tikzset{every edge quotes/.append style={fill=white,inner sep=1.2pt}}
%% FIGFIX-END
\begin{document}

\pagegarde{Équations, inéquations et systèmes}{Mathématiques}{1ère A}{Relations et opérations fondamentales dans l'ensemble des nombres réels}{NA01}{12 h}{Cette leçon outille l'élève pour résoudre dans ℝ les équations et inéquations du premier degré à quotient, du second degré, ainsi que les systèmes de deux équations à deux inconnues. Elle développe les techniques fondamentales (forme canonique, discriminant, factorisation, tableaux de signes) indispensables pour la suite du parcours scientifique et pour modéliser des situations concrètes de la vie quotidienne.
\vspace{4pt}
\begin{multicols}{2}\small
\begin{itemize}
\item À la fin de cette leçon, je sais résoudre dans ℝ une équation ou une inéquation se ramenant à (ax+b)/(cx+d) = 0 ou (ax+b)/(cx+d) < 0 (≤, >, ≥)
\item À la fin de cette leçon, je sais dresser et exploiter le tableau de signes d'un quotient (ax+b)/(cx+d)
\item À la fin de cette leçon, je sais écrire un polynôme du second degré sous sa forme canonique
\item À la fin de cette leçon, je sais calculer le discriminant d'un polynôme du second degré et résoudre une équation du second degré dans ℝ
\item À la fin de cette leçon, je sais factoriser un polynôme du second degré et dresser son tableau de signes
\item À la fin de cette leçon, je sais résoudre dans ℝ une inéquation du second degré
\end{itemize}\end{multicols}}

\begin{sommaire}
\begin{multicols}{2}
\begin{enumerate}
\item Équations et inéquations se ramenant à un quotient (ax+b)/(cx+d)
\item Forme canonique d'un polynôme du second degré
\item Discriminant et résolution des équations du second degré
\item Factorisation et signe d'un polynôme du second degré
\item Systèmes d'équations se ramenant au second degré dans ℝ²
\item Problèmes de synthèse et modélisation
\item Activité d'intégration
\item Méthodes et exercices résolus
\item Exercices
\item Corrigés des exercices
\item Fiche bilan
\end{enumerate}
\end{multicols}
\end{sommaire}

\begin{objectifs}
\begin{itemize}
\item À la fin de cette leçon, je sais résoudre dans ℝ une équation ou une inéquation se ramenant à (ax+b)/(cx+d) = 0 ou (ax+b)/(cx+d) < 0 (≤, >, ≥)
\item À la fin de cette leçon, je sais dresser et exploiter le tableau de signes d'un quotient (ax+b)/(cx+d)
\item À la fin de cette leçon, je sais écrire un polynôme du second degré sous sa forme canonique
\item À la fin de cette leçon, je sais calculer le discriminant d'un polynôme du second degré et résoudre une équation du second degré dans ℝ
\item À la fin de cette leçon, je sais factoriser un polynôme du second degré et dresser son tableau de signes
\item À la fin de cette leçon, je sais résoudre dans ℝ une inéquation du second degré
\item À la fin de cette leçon, je sais vérifier qu'un couple de réels est solution d'un système d'équations dans ℝ²
\item À la fin de cette leçon, je sais résoudre dans ℝ² un système d'équations se ramenant à une équation du second degré
\item À la fin de cette leçon, je sais modéliser et résoudre un problème concret se ramenant à une équation, une inéquation ou un système étudié
\end{itemize}
\end{objectifs}

\begin{prerequis}
\begin{itemize}
\item Résolution des équations et inéquations du premier degré à une inconnue (classe de Seconde)
\item Signe d'un binôme ax + b et tableau de signes d'un produit (classe de Seconde)
\item Développement et factorisation d'expressions algébriques, identités remarquables
\item Résolution des systèmes de deux équations linéaires à deux inconnues (substitution, combinaison)
\item Notion d'ensemble de définition d'une expression littérale (valeurs interdites)
\end{itemize}
\end{prerequis}

\newpage

\coursec{Équations et inéquations se ramenant à un quotient (ax+b)/(cx+d)}

\textbf{Situation d'accroche.} Au marché Mokolo, Mama Ngo Bell est une vendeuse de beignets réputée. Chaque matin, elle se pose la même question : à quel prix vendre ses beignets pour que la journée soit rentable ? Son fils, qui a étudié les mathématiques, lui a expliqué que si elle fixe le prix unitaire à $x$ francs, sa recette journalière (en milliers de francs) est donnée par la formule
\[ R(x) = -2x^2 + 300x. \]
Mama Ngo Bell veut savoir : \emph{quel prix fixer pour obtenir une recette d'au moins \num{10000} F}, c'est-à-dire $R(x) \geq 10$ ? Pendant ce temps, son fils, qui gère aussi un atelier de couture à Yaoundé, doit partager \num{2400000} F entre deux associés dont les parts vérifient des conditions précises reliant leur somme et leur produit.

Ces deux problèmes, en apparence très différents, se résolvent avec les mêmes outils : les \cle{équations}, les \cle{inéquations} et les \cle{systèmes}. Dans cette première section, nous commençons par le cas le plus délicat à manipuler : les expressions en \emph{quotient}, de la forme $\dfrac{ax+b}{cx+d}$, qui apparaissent naturellement dès qu'on compare des grandeurs liées de manière inversement proportionnelle (prix moyen, vitesse moyenne, concentration d'une solution...).

\textbf{Problématique.} Comment déterminer rigoureusement pour quelles valeurs de $x$ un quotient $\dfrac{ax+b}{cx+d}$ est nul, positif ou négatif, sans jamais tomber dans les pièges du dénominateur ?

\courssub{Rappel : signe d'un binôme $ax + b$}

Avant d'étudier un quotient, il faut maîtriser le signe de chacun de ses morceaux. Rappelons le résultat fondamental vu en Seconde.

\begin{propriete}[Signe du binôme $ax+b$ ($a \neq 0$)]
Le binôme $ax+b$ s'annule en $x_0 = -\dfrac{b}{a}$. De plus :
\begin{itemize}
\item si $a > 0$ : $ax+b$ est \textbf{négatif} avant $x_0$ et \textbf{positif} après $x_0$ ;
\item si $a < 0$ : $ax+b$ est \textbf{positif} avant $x_0$ et \textbf{négatif} après $x_0$.
\end{itemize}
Autrement dit, le binôme $ax+b$ est \emph{du signe de $a$} après sa racine.
\end{propriete}

\begin{exemplebox}[Signe de $2x-4$ et de $x+1$]
$2x-4$ s'annule en $x = 2$ ; comme $a = 2 > 0$, il est négatif sur $]-\infty ; 2]$ et positif sur $[2 ; +\infty[$.

$x+1$ s'annule en $x = -1$ ; comme $a = 1 > 0$, il est négatif sur $]-\infty ; -1]$ et positif sur $[-1 ; +\infty[$.
\end{exemplebox}

\begin{attention}[Le signe dépend de celui de $a$]
L'erreur classique consiste à dire « le binôme est négatif avant sa racine » sans regarder le signe de $a$. Pour $-3x+6$ (racine $2$, $a = -3 < 0$), c'est l'inverse : le binôme est \textbf{positif} avant $2$ et négatif après. Vérifie toujours le signe du coefficient de $x$ avant de remplir ton tableau.
\end{attention}

\courssub{Ensemble de définition et valeur interdite}

Voici l'intuition essentielle : \emph{on ne peut jamais diviser par zéro}. Si tu partages équitablement $0$ F entre $5$ amis, chacun reçoit $0$ F : pas de problème. Mais partager $5000$ F entre $0$ personne n'a aucun sens. En mathématiques, une expression quotient n'existe que si son dénominateur est non nul.

\begin{definition}[Valeur interdite, ensemble de définition]
Soit l'expression quotient $Q(x) = \dfrac{ax+b}{cx+d}$ avec $c \neq 0$.
\begin{itemize}
\item La \cle{valeur interdite} est le nombre $x_0 = -\dfrac{d}{c}$ qui annule le dénominateur $cx+d$.
\item L'\cle{ensemble de définition} de $Q$ est l'ensemble des réels pour lesquels le quotient existe :
\[ \mathcal{D}_Q = \mathbb{R} \setminus \left\{ -\frac{d}{c} \right\} = \left] -\infty ; -\frac{d}{c} \right[ \cup \left] -\frac{d}{c} ; +\infty \right[. \]
\end{itemize}
\end{definition}

\begin{exemplebox}[Ensemble de définition de $\dfrac{3x-6}{x+2}$]
Le dénominateur $x+2$ s'annule en $x = -2$. La valeur interdite est donc $-2$, et
\[ \mathcal{D} = \mathbb{R} \setminus \{-2\} = ]-\infty ; -2[ \cup ]-2 ; +\infty[. \]
Toute résolution d'équation ou d'inéquation faisant intervenir ce quotient devra \textbf{commencer} par cette remarque et \textbf{se terminer} en excluant $-2$ des solutions éventuelles.
\end{exemplebox}

\courssub{Résolution de l'équation $\dfrac{ax+b}{cx+d} = 0$}

Quand une fraction vaut-elle zéro ? Pense à une part de gâteau : si le gâteau (le numérateur) n'existe pas, chacun reçoit zéro ; mais si le nombre de convives (le dénominateur) est nul, la question elle-même n'a pas de sens. D'où le théorème suivant.

\begin{propriete}[Quotient nul]
Pour tout $x \in \mathcal{D}_Q$ :
\[ \frac{ax+b}{cx+d} = 0 \iff \begin{cases} ax+b = 0 \\ cx+d \neq 0 \end{cases} \]
Autrement dit : \emph{un quotient est nul si et seulement si son numérateur est nul et son dénominateur non nul}.
\end{propriete}

\begin{experience}[Démonstration guidée]
Supposons $cx+d \neq 0$ (sinon le quotient n'existe pas).
\begin{enumerate}
\item \textbf{Sens direct.} Si $\dfrac{ax+b}{cx+d} = 0$, multiplions les deux membres par le nombre non nul $cx+d$ : on obtient $ax+b = 0 \times (cx+d) = 0$. Donc le numérateur est nul.
\item \textbf{Sens réciproque.} Si $ax+b = 0$ et $cx+d \neq 0$, alors $\dfrac{ax+b}{cx+d} = \dfrac{0}{cx+d} = 0$, car zéro divisé par un nombre non nul donne zéro.
\end{enumerate}
Les deux sens étant établis, l'équivalence est démontrée. $\blacksquare$
\end{experience}

\begin{exoresolu}[Équation quotient nulle]
Résoudre dans $\mathbb{R}$ l'équation $\dfrac{2x-4}{x+1} = 0$.
\tcblower
\textbf{Étape 1 — Ensemble de définition.} $x+1 = 0 \iff x = -1$. Donc $\mathcal{D} = \mathbb{R} \setminus \{-1\}$.

\textbf{Étape 2 — Numérateur nul.} $2x-4 = 0 \iff x = 2$.

\textbf{Étape 3 — Vérification.} $2 \neq -1$ : la valeur trouvée n'est pas interdite, elle est bien dans $\mathcal{D}$.

\textbf{Conclusion.} $S = \{2\}$.
\end{exoresolu}

\begin{methode}[Équations se ramenant à un quotient nul]
Pour résoudre une équation comportant plusieurs quotients :
\begin{enumerate}
\item Déterminer \textbf{toutes} les valeurs interdites (chaque dénominateur non nul).
\item \textbf{Se ramener à zéro} : tout passer dans un seul membre.
\item \textbf{Réduire au même dénominateur} pour obtenir un unique quotient $\dfrac{N(x)}{D(x)} = 0$.
\item Résoudre $N(x) = 0$, puis \textbf{écarter} toute solution qui serait une valeur interdite.
\end{enumerate}
\end{methode}

\begin{exoresolu}[Mise au même dénominateur]
Résoudre dans $\mathbb{R}$ : $\dfrac{3}{x-1} = \dfrac{2}{x+3}$.
\tcblower
\textbf{Valeurs interdites :} $x \neq 1$ et $x \neq -3$, donc $\mathcal{D} = \mathbb{R} \setminus \{-3 ; 1\}$.

On se ramène à zéro et on réduit au même dénominateur $(x-1)(x+3)$ :
\[ \frac{3}{x-1} - \frac{2}{x+3} = 0 \iff \frac{3(x+3) - 2(x-1)}{(x-1)(x+3)} = 0 \iff \frac{x + 11}{(x-1)(x+3)} = 0. \]
En effet, $3(x+3) - 2(x-1) = 3x + 9 - 2x + 2 = x + 11$. Le quotient est nul si et seulement si $x + 11 = 0$, soit $x = -11$. Or $-11 \notin \{-3 ; 1\}$ : la solution est acceptée.

\textbf{Conclusion :} $S = \{-11\}$.

\emph{Remarque :} sur $\mathcal{D}$, le produit en croix $3(x+3) = 2(x-1)$ donne le même résultat, car les dénominateurs y sont non nuls. Mais cette méthode n'est légitime qu'\textbf{après} avoir exclu les valeurs interdites.
\end{exoresolu}

\courssub{Tableau de signes du quotient $\dfrac{ax+b}{cx+d}$}

Le signe d'un quotient se déduit des signes du numérateur et du dénominateur par la \emph{règle des signes} : positif si les deux ont le même signe, négatif sinon. On rassemble tout dans un tableau à trois lignes, en plaçant une \textbf{double barre} sous la valeur interdite : elle matérialise le « mur » infranchissable où le quotient n'existe pas.

\begin{exemplebox}[Tableau de signes de $\dfrac{2x-4}{x+1}$]
Le numérateur s'annule en $2$, le dénominateur en $-1$ (valeur interdite). On range les valeurs dans l'ordre croissant : $-1$ puis $2$.
\begin{center}
\begin{tabular}{|c|cc||ccccc|}
\hline
$x$ & $-\infty$ & & $-1$ & & $2$ & & $+\infty$ \\
\hline
$2x-4$ & & $-$ & & $-$ & $0$ & $+$ & \\
\hline
$x+1$ & & $-$ & $0$ & $+$ & & $+$ & \\
\hline
$\dfrac{2x-4}{x+1}$ & & $+$ & & $-$ & $0$ & $+$ & \\
\hline
\end{tabular}
\end{center}
Lecture : le quotient est strictement positif sur $]-\infty ; -1[ \cup ]2 ; +\infty[$, strictement négatif sur $]-1 ; 2[$, nul en $2$, et \textbf{inexistant} en $-1$ (la double barre du tableau le rappelle).
\end{exemplebox}

\begin{attention}[Double barre obligatoire]
Dans la ligne du quotient, la valeur interdite est marquée par une double barre verticale dans le tableau, jamais par un $0$. Un $0$ signifierait que le quotient vaut zéro en ce point, ce qui est faux : il n'y est même pas défini. Cette confusion coûte cher au Probatoire !
\end{attention}

\courssub{Résolution des inéquations quotient}

Pour résoudre $\dfrac{ax+b}{cx+d} < 0$ (ou $\leq 0$, $> 0$, $\geq 0$), on lit simplement le tableau de signes, en appliquant les règles suivantes :
\begin{itemize}
\item la \textbf{valeur interdite} est toujours exclue (crochet ouvert) ;
\item la racine du numérateur est \textbf{incluse} (crochet fermé) si l'inégalité est large ($\leq$ ou $\geq$), exclue si elle est stricte ($<$ ou $>$).
\end{itemize}

\begin{methode}[Résoudre une inéquation quotient en 4 étapes]
\begin{enumerate}
\item \textbf{Ensemble de définition} : trouver la ou les valeurs interdites.
\item \textbf{Se ramener à une comparaison à 0} : tout passer dans un membre et réduire au même dénominateur pour obtenir $\dfrac{N(x)}{D(x)} \leq 0$ (par exemple).
\item \textbf{Tableau de signes} : lignes de $N(x)$, de $D(x)$, puis du quotient, avec double barre à la valeur interdite.
\item \textbf{Conclusion} : lire les intervalles convenables en excluant toujours la valeur interdite.
\end{enumerate}
\end{methode}

\begin{exoresolu}[L'exemple de référence]
Résoudre dans $\mathbb{R}$ l'inéquation $\dfrac{3x-6}{x+2} \leq 0$.
\tcblower
\textbf{Étape 1.} $x+2 = 0 \iff x = -2$ : valeur interdite. $\mathcal{D} = \mathbb{R} \setminus \{-2\}$.

\textbf{Étape 2.} L'inéquation est déjà comparée à $0$.

\textbf{Étape 3.} $3x-6 = 0 \iff x = 2$. Tableau de signes :
\begin{center}
\begin{tabular}{|c|cc||ccccc|}
\hline
$x$ & $-\infty$ & & $-2$ & & $2$ & & $+\infty$ \\
\hline
$3x-6$ & & $-$ & & $-$ & $0$ & $+$ & \\
\hline
$x+2$ & & $-$ & $0$ & $+$ & & $+$ & \\
\hline
$\dfrac{3x-6}{x+2}$ & & $+$ & & $-$ & $0$ & $+$ & \\
\hline
\end{tabular}
\end{center}

\textbf{Étape 4.} On cherche où le quotient est négatif ou nul : c'est l'intervalle $]-2 ; 2[$, auquel on ajoute $x = 2$ (inégalité large, le quotient y vaut $0$) mais pas $-2$ (valeur interdite).

\textbf{Conclusion :} $S = ]-2 ; 2]$.
\end{exoresolu}

\begin{exoresolu}[Inéquation se ramenant à un quotient]
Résoudre dans $\mathbb{R}$ : $\dfrac{2x+1}{x-3} \geq 1$.
\tcblower
\textbf{Étape 1.} Valeur interdite : $x = 3$. $\mathcal{D} = \mathbb{R} \setminus \{3\}$.

\textbf{Étape 2.} On se ramène à $0$ (surtout pas de produit en croix !) :
\[ \frac{2x+1}{x-3} - 1 \geq 0 \iff \frac{2x+1 - (x-3)}{x-3} \geq 0 \iff \frac{x+4}{x-3} \geq 0. \]
En effet, $2x+1-(x-3) = 2x+1-x+3 = x+4$.

\textbf{Étape 3.} Racines : $-4$ (numérateur) et $3$ (dénominateur).
\begin{center}
\begin{tabular}{|c|ccccc||cc|}
\hline
$x$ & $-\infty$ & & $-4$ & & $3$ & & $+\infty$ \\
\hline
$x+4$ & & $-$ & $0$ & $+$ & & $+$ & \\
\hline
$x-3$ & & $-$ & & $-$ & $0$ & $+$ & \\
\hline
$\dfrac{x+4}{x-3}$ & & $+$ & $0$ & $-$ & & $+$ & \\
\hline
\end{tabular}
\end{center}

\textbf{Étape 4.} Quotient positif ou nul sur $]-\infty ; -4]$ (on inclut $-4$) et sur $]3 ; +\infty[$ (on exclut $3$).

\textbf{Conclusion :} $S = ]-\infty ; -4] \cup ]3 ; +\infty[$.
\end{exoresolu}

\begin{attention}[Le piège du produit en croix dans une inéquation]
Face à $\dfrac{2x+1}{x-3} \geq 1$, beaucoup d'élèves écrivent $2x+1 \geq x-3$, c'est-à-dire multiplient les deux membres par $x-3$. \textbf{C'est faux en général !} Multiplier une inégalité par un nombre \emph{négatif} en change le sens. Or le signe de $x-3$ dépend de $x$ : pour $x < 3$, il faudrait écrire $2x+1 \leq x-3$, et pour $x > 3$, $2x+1 \geq x-3$. Plutôt que de discuter deux cas, on applique systématiquement la méthode en 4 étapes : \emph{on ne multiplie jamais une inéquation par une expression dont on ne connaît pas le signe}.
\end{attention}

\courssub{Lien avec les fonctions homographiques}

\begin{definition}[Fonction homographique]
On appelle \cle{fonction homographique} toute fonction de la forme
\[ f : x \longmapsto \frac{ax+b}{cx+d} \quad (c \neq 0,\ ad - bc \neq 0), \]
définie sur $\mathbb{R} \setminus \left\{-\dfrac{d}{c}\right\}$.
\end{definition}

La condition $ad - bc \neq 0$ garantit que $f$ n'est pas une fonction constante déguisée : si $ad - bc = 0$, le numérateur serait proportionnel au dénominateur et le quotient se simplifierait en une constante.

Sa courbe représentative est une \cle{hyperbole} : elle possède une \textbf{asymptote verticale} d'équation $x = -\dfrac{d}{c}$ (la valeur interdite !) et une \textbf{asymptote horizontale} d'équation $y = \dfrac{a}{c}$. Le signe du quotient se lit alors directement sur le graphique : le quotient est positif là où la courbe est \emph{au-dessus} de l'axe des abscisses, négatif là où elle est \emph{en dessous}.

\begin{popfigure}
\begin{center}
\begin{tikzpicture}
\begin{axis}[
  width=0.9\linewidth, height=7cm,
  axis lines=middle, xlabel=$x$, ylabel=$y$,
  xmin=-6, xmax=5, ymin=-6, ymax=8,
  xtick={-4,-2,-1,2,4}, ytick={-4,-2,2,4,6},
  grid=both, grid style={popline!40},
  clip=true]
  \addplot[popcurve, domain=-6:-1.15, samples=150] {(2*x-4)/(x+1)};
  \addplot[popcurve, domain=-0.85:5, samples=150] {(2*x-4)/(x+1)};
  \addplot[poporange, dashed, thick] coordinates {(-1,-6) (-1,8)};
  \addplot[popgreen, dashed, thick, domain=-6:5] {2};
  \node[poporange, anchor=west] at (axis cs:-0.9,-4.5) {\small $x=-1$};
  \node[popgreen, anchor=south west] at (axis cs:2.2,2.1) {\small $y=2$};
  \addplot[only marks, mark=*, popdark] coordinates {(2,0)};
  \node[popdark, anchor=south] at (axis cs:2,0.2) {\small $2$};
\end{axis}
\end{tikzpicture}
\end{center}
\legende{Courbe de $f(x)=\dfrac{2x-4}{x+1}$ : hyperbole d'asymptotes $x=-1$ et $y=2$. La courbe est au-dessus de l'axe des abscisses sur $]-\infty;-1[\,\cup\,]2;+\infty[$ (quotient positif) et en dessous sur $]-1;2[$ (quotient négatif) : on retrouve exactement le tableau de signes.}
\end{popfigure}

\begin{savaistu}[Les homographiques autour de nous]
Les fonctions homographiques ne sont pas qu'un exercice de style ! En \textbf{optique}, la relation de conjugaison des lentilles minces, $\dfrac{1}{f'} = \dfrac{1}{p} + \dfrac{1}{p'}$, fait intervenir des fonctions homographiques : c'est grâce à elles qu'on calcule la vergence des lunettes correctrices. En \textbf{économie}, les partages inversement proportionnels (par exemple répartir un coût fixe entre des consommations variables, comme le partage d'une facture CAMTEL entre abonnés) conduisent à des quotients du type $\dfrac{k}{x}$ ou $\dfrac{ax+b}{cx+d}$. Les Grecs connaissaient déjà l'hyperbole comme section d'un cône, mais c'est au XVII\textsuperscript{e} siècle, avec Descartes et Fermat, que son équation a été reliée à ces quotients.
\end{savaistu}

\begin{exercice}[Pour t'entraîner]
\begin{enumerate}
\item Déterminer l'ensemble de définition de $g(x) = \dfrac{5x+10}{2x-6}$, puis résoudre $g(x) = 0$.
\item Résoudre dans $\mathbb{R}$ : $\dfrac{4-x}{x+3} > 0$.
\item Résoudre dans $\mathbb{R}$ : $\dfrac{x}{x-1} \leq 2$.
\end{enumerate}
\end{exercice}

\begin{corrige}[Exercice 1]
\textbf{1.} $2x-6 = 0 \iff x = 3$ : $\mathcal{D} = \mathbb{R} \setminus \{3\}$. $g(x) = 0 \iff 5x+10 = 0 \iff x = -2$, et $-2 \neq 3$. Donc $S = \{-2\}$.

\textbf{2.} Racines : $4$ (numérateur) et $-3$ (dénominateur, valeur interdite). Attention : $4-x$ a un coefficient $a = -1 < 0$, il est donc positif avant $4$ et négatif après. Le tableau de signes donne un quotient strictement positif sur $]-3 ; 4[$. $S = ]-3 ; 4[$.

\textbf{3.} On se ramène à $0$ : $\dfrac{x}{x-1} - 2 \leq 0 \iff \dfrac{x - 2(x-1)}{x-1} \leq 0 \iff \dfrac{-x+2}{x-1} \leq 0$. Racines : $2$ et $1$ (interdite). Le binôme $-x+2$ est positif avant $2$, négatif après ; le tableau donne un quotient négatif ou nul sur $]-\infty ; 1[ \cup [2 ; +\infty[$. $S = ]-\infty ; 1[ \cup [2 ; +\infty[$.
\end{corrige}

\begin{aretenir}[L'essentiel de la section]
\begin{itemize}
\item Une expression quotient $\dfrac{ax+b}{cx+d}$ n'existe que si $cx+d \neq 0$ : la valeur $x = -\dfrac{d}{c}$ est \textbf{interdite} et sera toujours exclue des solutions.
\item Un quotient est \textbf{nul} si et seulement si son numérateur est nul (et son dénominateur non nul).
\item Pour une \textbf{inéquation} : jamais de produit en croix ! On se ramène à une comparaison à $0$, on dresse le tableau de signes (double barre à la valeur interdite), puis on conclut en respectant les crochets.
\item La courbe d'une fonction homographique est une hyperbole ; son asymptote verticale correspond exactement à la valeur interdite, et le signe du quotient se lit par la position de la courbe par rapport à l'axe des abscisses.
\end{itemize}
\end{aretenir}

\coursec{Forme canonique d'un polynôme du second degré}

Après avoir maîtrisé la résolution des équations et inéquations faisant intervenir des quotients de fonctions affines $\frac{ax+b}{cx+d}$, nous disposons maintenant des outils pour attaquer les polynômes du second degré. La \cle{forme canonique} est la clé de voûte de ce chapitre : elle révèle instantanément le sommet de la parabole, son axe de symétrie et permet de résoudre des équations ou des problèmes d'optimisation sans systématiquement calculer le discriminant. C'est l'écriture « intelligente » du trinôme.

\courssub{Rappel et définition : polynôme du second degré}

\begin{definition}[Polynôme du second degré et forme canonique]
Soient $a,b,c \in \mathbb{R}$ avec $a \neq 0$. La fonction $P$ définie sur $\mathbb{R}$ par 
\[
P(x) = ax^2 + bx + c
\]
est appelée \cle{polynôme du second degré} (ou \cle{trinôme du second degré}). Le coefficient $a$ est le \cle{coefficient directeur}, $b$ le coefficient du terme linéaire et $c$ le \cle{terme constant}.

On appelle \cle{forme canonique} de $P$ l'écriture :
\[
P(x) = a(x - \alpha)^2 + \beta
\]
où 
\[
\alpha = -\frac{b}{2a} \qquad \text{et} \qquad \beta = P(\alpha) = \frac{4ac - b^2}{4a} = -\frac{\Delta}{4a},
\]
$\Delta = b^2 - 4ac$ étant le \cle{discriminant} de $P$ (nous l'étudierons en détail à la section 3).
\end{definition}

\begin{savaistu}[Pourquoi « canonique » ?]
Le terme « canonique » vient du grec \emph{kanon} (règle, modèle). C'est l'écriture \emph{de référence}, celle qui met en évidence la structure géométrique de la parabole : son sommet $S(\alpha\,;\,\beta)$ et son axe de symétrie $x = \alpha$. Au Cameroun, on la rencontre dès la Première pour modéliser la trajectoire d'un projectile (arrosage agricole, tir de pénalité au stade d'Olembé) ou optimiser un coût de production (usine Alucam, brasseries).
\end{savaistu}

\courssub{De la forme développée à la forme canonique : la méthode du carré parfait}

L'idée directrice est l'\cle{identité remarquable} du carré parfait : $(u+v)^2 = u^2 + 2uv + v^2$. Si nous avons $x^2 + bx$, il nous manque le carré de la moitié du coefficient de $x$ pour former un carré parfait : $\left(\frac{b}{2}\right)^2$. On « complète le carré » en ajoutant et en soustrayant ce terme.

\begin{methode}[Compléter le carré — Cas $a=1$ (exemple numérique : $x^2 - 4x + 3$)]
\emph{Objectif :} Écrire $P(x) = x^2 - 4x + 3$ sous la forme $(x-\alpha)^2 + \beta$.

\begin{enumerate}
\item \textbf{Repérer le coefficient de $x$ :} ici $-4$.
\item \textbf{Calculer la moitié de ce coefficient :} $\frac{-4}{2} = -2$.
\item \textbf{Écrire le carré correspondant :} $(x - 2)^2 = x^2 - 4x + 4$.
\item \textbf{Comparer avec $P(x)$ :} $P(x) = x^2 - 4x + 3 = (x^2 - 4x + 4) - 1$.
\item \textbf{Conclure :} $P(x) = (x - 2)^2 - 1$. Donc $\alpha = 2$, $\beta = -1$.
\end{enumerate}
\end{methode}

\begin{attention}[Piège majeur : le coefficient $a$ devant $x^2$]
Si $a \neq 1$, \textbf{il est impératif de factoriser $a$ sur les deux premiers termes} avant de compléter le carré.
\emph{Erreur fréquente :} Écrire $2x^2 - 8x + 5 = 2(x^2 - 4x) + 5$ puis compléter le carré \emph{à l'intérieur} de la parenthèse sans reporter le facteur $2$ correctement sur le terme ajouté/soustrait.
\emph{Bonne méthode :} $2x^2 - 8x + 5 = 2(x^2 - 4x) + 5 = 2\bigl[(x-2)^2 - 4\bigr] + 5 = 2(x-2)^2 - 8 + 5 = 2(x-2)^2 - 3$.
\end{attention}

\begin{propriete}[Théorème : Forme canonique d'un trinôme — Démonstration exigible au Probatoire]
Tout polynôme du second degré $P(x) = ax^2 + bx + c$ ($a \neq 0$) admet une unique forme canonique :
\[
P(x) = a\left(x + \frac{b}{2a}\right)^2 - \frac{b^2 - 4ac}{4a} = a\left(x + \frac{b}{2a}\right)^2 - \frac{\Delta}{4a}.
\]

\textbf{Démonstration.}
\begin{enumerate}
\item \textbf{Factoriser $a$ sur les deux premiers termes :}
$P(x) = a\left(x^2 + \frac{b}{a}x\right) + c$.
\item \textbf{Compléter le carré à l'intérieur de la parenthèse :}
Le coefficient de $x$ est $\frac{b}{a}$. Sa moitié est $\frac{b}{2a}$. Son carré est $\frac{b^2}{4a^2}$.
On ajoute et on soustrait ce terme \emph{à l'intérieur} :
$P(x) = a\left(x^2 + \frac{b}{a}x + \frac{b^2}{4a^2} - \frac{b^2}{4a^2}\right) + c$.
\item \textbf{Factoriser le trinôme parfait :}
$x^2 + \frac{b}{a}x + \frac{b^2}{4a^2} = \left(x + \frac{b}{2a}\right)^2$.
\item \textbf{Regrouper les constantes :}
$P(x) = a\left[\left(x + \frac{b}{2a}\right)^2 - \frac{b^2}{4a^2}\right] + c
= a\left(x + \frac{b}{2a}\right)^2 - \frac{b^2}{4a} + c$.
\item \textbf{Mettre au même dénominateur $4a$ :}
$-\frac{b^2}{4a} + c = \frac{-b^2 + 4ac}{4a} = -\frac{b^2 - 4ac}{4a} = -\frac{\Delta}{4a}$.
\item \textbf{Forme finale :}
$P(x) = a\left(x + \frac{b}{2a}\right)^2 - \frac{\Delta}{4a}$.
En posant $\alpha = -\frac{b}{2a}$ et $\beta = -\frac{\Delta}{4a}$, on obtient $P(x) = a(x-\alpha)^2 + \beta$. \qed
\end{enumerate}
\end{propriete}

\begin{exemplebox}[Mise en forme canonique : $P(x) = 2x^2 - 8x + 5$]
Appliquons la méthode rigoureuse (factorisation de $a$ d'abord).

\begin{align*}
P(x) &= 2x^2 - 8x + 5 \\
&= 2(x^2 - 4x) + 5 \quad \text{(factorisation de 2 sur les deux premiers termes)} \\
&= 2\left[ x^2 - 4x + \left(\frac{-4}{2}\right)^2 - \left(\frac{-4}{2}\right)^2 \right] + 5 \quad \text{(complétion du carré : moitié de -4 est -2, carré = 4)} \\
&= 2\left[ (x - 2)^2 - 4 \right] + 5 \quad \text{(factorisation du trinôme parfait)} \\
&= 2(x - 2)^2 - 8 + 5 \quad \text{(distributivité du 2)} \\
&= 2(x - 2)^2 - 3.
\end{align*}

\textbf{Résultat :} Forme canonique $P(x) = 2(x-2)^2 - 3$.
\cle{Sommet} $S(2\,;\,-3)$, \cle{axe de symétrie} $x = 2$.
Comme $a=2>0$, $P$ admet un \cle{minimum} égal à $-3$ atteint en $x=2$.
\end{exemplebox}

\courssub{Lecture géométrique de la forme canonique}

La forme canonique $P(x) = a(x-\alpha)^2 + \beta$ est une mine d'informations géométriques sur la courbe représentative $\mathcal{C}_P$ (parabole).

\begin{aretenir}[Lecture immédiate sur la forme canonique]
Pour $P(x) = a(x-\alpha)^2 + \beta$ ($a \neq 0$) :
\begin{itemize}
\item Le \cle{sommet} de la parabole $\mathcal{C}_P$ est le point $S(\alpha\,;\,\beta)$.
\item L'\cle{axe de symétrie} de $\mathcal{C}_P$ est la droite d'équation $x = \alpha$.
\item \textbf{Sens de la parabole :}
\begin{itemize}
\item Si $a > 0$ : la parabole est tournée vers le haut ($\cup$). $P$ admet un \cle{minimum} $\beta$ atteint en $x=\alpha$.
\item Si $a < 0$ : la parabole est tournée vers le bas ($\cap$). $P$ admet un \cle{maximum} $\beta$ atteint en $x=\alpha$.
\end{itemize}
\item \textbf{Variations :} $P$ est décroissante sur $]-\infty\,;\,\alpha]$ et croissante sur $[\alpha\,;\,+\infty[$ si $a>0$ (inverse si $a<0$).
\end{itemize}
\end{aretenir}

\begin{popfigure}
\begin{tikzpicture}
\begin{axis}[
    width=\linewidth,
    height=8cm,
    axis lines=middle,
    xlabel=$x$, ylabel=$y$,
    xmin=-1, xmax=5,
    ymin=-2.5, ymax=4,
    xtick={-1,0,1,2,3,4,5},
    ytick={-2,-1,0,1,2,3,4},
    tick label style={font=\small},
    label style={font=\small},
    grid=major,
    grid style={popline, opacity=0.5},
    clip=false
]
% Parabole P(x) = x^2 - 4x + 3 = (x-2)^2 - 1
\addplot[popcurve, domain=-0.5:4.5, samples=200, thick] {x^2 - 4*x + 3};
% Sommet S(2,-1)
\addplot[only marks, mark=*, mark size=2.5, poporange] coordinates {(2,-1)};
\node[above right, poporange, font=\small\bfseries] at (axis cs:2,-1) {$S(2\,;\,-1)$};
% Axe de symétrie x=2
\addplot[poppurple, dashed, thick, domain=-2:4] ({2}, x);
\node[poppurple, font=\small\bfseries, rotate=90, anchor=south] at (axis cs:2,3.5) {Axe de symétrie $x=2$};
% Points d'intersection avec Ox : 1 et 3
\addplot[only marks, mark=*, mark size=2, popblue] coordinates {(1,0) (3,0)};
\node[below, popblue, font=\small] at (axis cs:1,0) {$1$};
\node[below, popblue, font=\small] at (axis cs:3,0) {$3$};
% Point d'intersection avec Oy : 3
\addplot[only marks, mark=*, mark size=2, popblue] coordinates {(0,3)};
\node[left, popblue, font=\small] at (axis cs:0,3) {$3$};
% Flèches de variation
\draw[popgreen, thick, -{Stealth[length=3mm]}] (axis cs:0,3) .. controls (axis cs:1,1) .. (axis cs:1.5,-0.5);
\draw[popgreen, thick, -{Stealth[length=3mm]}] (axis cs:2.5,-0.5) .. controls (axis cs:3,1) .. (axis cs:4,3);
\node[popgreen, font=\small\bfseries] at (axis cs:0.5,3.5) {Décroît};
\node[popgreen, font=\small\bfseries] at (axis cs:3.5,3.5) {Croît};
\end{axis}
\end{tikzpicture}
\legende{Parabole $\mathcal{C}_P$ de $P(x)=x^2-4x+3 = (x-2)^2-1$. Sommet $S(2\,;\,-1)$, axe $x=2$, racines $1$ et $3$.}
\end{popfigure}

\begin{popfigure}
\begin{tikzpicture}[scale=0.9, every node/.style={font=\small}]
% Repère de base y = x^2
\begin{scope}[shift={(0,0)}]
\draw[popline, ->] (-2.5,0) -- (2.5,0) node[right] {$x$};
\draw[popline, ->] (0,-0.5) -- (0,3.5) node[above] {$y$};
\draw[popcurve, thick, domain=-1.5:1.5, samples=100] plot (\x, {\x*\x});
\node[popdark, anchor=north east] at (1.5,2.2) {$y = x^2$};
\node[popdark, anchor=south] at (0,0) {$O$};
\draw[dashed, popmuted] (-1.5,0) -- (-1.5,2.25) -- (0,2.25);
\draw[dashed, popmuted] (1.5,0) -- (1.5,2.25);
\end{scope}

% Flèche de translation
\draw[poporange, thick, -{Stealth[length=4mm]}] (3.5,1.5) -- (5.5,1.5);
\node[poporange, above, font=\bfseries] at (4.5,1.8) {Translation du graphe $\vec{v}(2\,;\,-1)$};

% Repère final y = (x-2)^2 - 1
\begin{scope}[shift={(8,0)}]
\draw[popline, ->] (-0.5,0) -- (4.5,0) node[right] {$x$};
\draw[popline, ->] (2,-1.5) -- (2,2.5) node[above] {$y$};
\draw[popcurve, thick, domain=0:4, samples=100] plot (\x, {(\x-2)*(\x-2)-1});
\node[popdark, anchor=north west] at (0.5,1.5) {$y = (x-2)^2 - 1$};
% Sommet S
\fill[poporange] (2,-1) circle (2.5pt);
\node[poporange, above right] at (2,-1) {$S(2\,;\,-1)$};
% Axe de symétrie
\draw[poppurple, dashed, thick] (2,-1.5) -- (2,2.5);
\node[poppurple, right, font=\bfseries] at (2,2) {$x=2$};
% Repère original fantôme
\draw[popmuted, very thin, dashed] (-0.5,0) -- (4.5,0);
\draw[popmuted, very thin, dashed] (2,-1.5) -- (2,2.5);
\end{scope}
\end{tikzpicture}
\legende{Interprétation géométrique : la forme canonique $P(x) = (x-2)^2-1$ montre que la parabole s'obtient par translation du vecteur $\vec{v}(2\,;\,-1)$ de la parabole de référence $y=x^2$. Le sommet $O(0,0)$ devient $S(2,-1)$.}
\end{popfigure}

\courssub{Applications : résolution d'équations et problèmes d'optimisation}

La forme canonique permet de résoudre l'équation $P(x)=0$ (ou $P(x)=k$) par une simple manipulation algébrique, sans passer par la formule du discriminant (surtout quand $a=1$). Elle est aussi l'outil roi pour les problèmes de \cle{maximum} ou \cle{minimum} (optimisation).

\begin{methode}[Résoudre $x^2 + bx + c = 0$ par la forme canonique ($a=1$)]
\begin{enumerate}
\item Écrire $x^2 + bx + c$ sous la forme $(x + \frac{b}{2})^2 - \frac{b^2}{4} + c = (x + \frac{b}{2})^2 - \frac{\Delta}{4}$.
\item L'équation devient $(x + \frac{b}{2})^2 = \frac{\Delta}{4}$.
\item \begin{itemize}
\item Si $\Delta < 0$ : aucune solution (carré ne peut être négatif).
\item Si $\Delta = 0$ : solution unique $x = -\frac{b}{2}$.
\item Si $\Delta > 0$ : $x + \frac{b}{2} = \pm \frac{\sqrt{\Delta}}{2}$, soit $x = \frac{-b \pm \sqrt{\Delta}}{2}$ (on retrouve la formule classique).
\end{itemize}
\end{enumerate}
\end{methode}

\begin{exoresolu}[Optimisation : Recette maximale d'un vendeur de « Poulet DG »]
Un vendeur de rue à Douala constate que s'il vend son plat à $1000$ FCFA, il en écoule $50$ par jour. Pour chaque augmentation de $50$ FCFA, il perd $2$ clients. Quel prix fixer pour maximiser sa recette journalière ?

\textbf{Solution détaillée.}

\begin{enumerate}
\item \textbf{Modélisation.} Soit $x$ le nombre d'augmentations de $50$ FCFA.
\begin{itemize}
\item Prix de vente : $p(x) = 1000 + 50x$ (FCFA).
\item Quantité vendue : $q(x) = 50 - 2x$ (plats). Condition : $q(x) \ge 0 \Rightarrow x \le 25$.
\end{itemize}
\item \textbf{Fonction recette.}
$R(x) = p(x) \times q(x) = (1000 + 50x)(50 - 2x)$.
Développons :
$R(x) = 50000 - 2000x + 2500x - 100x^2 = -100x^2 + 500x + 50000$.
\item \textbf{Forme canonique (factorisation de $-100$).}
\begin{align*}
R(x) &= -100(x^2 - 5x) + 50000 \\
&= -100\left[ x^2 - 5x + \left(\frac{5}{2}\right)^2 - \left(\frac{5}{2}\right)^2 \right] + 50000 \\
&= -100\left[ \left(x - \frac{5}{2}\right)^2 - \frac{25}{4} \right] + 50000 \\
&= -100\left(x - \frac{5}{2}\right)^2 + 625 + 50000 \\
&= -100\left(x - \frac{5}{2}\right)^2 + 50625.
\end{align*}
\item \textbf{Lecture du maximum.}
$a = -100 < 0$, donc $R$ admet un maximum au sommet.
Sommet : $\alpha = \frac{5}{2}$, $\beta = 50625$.
Le maximum de recette est \textbf{50\,625 FCFA}, atteint pour $x = \frac{5}{2}$ augmentations.
\item \textbf{Interprétation réelle.}
$x=\frac{5}{2}$ n'est pas un entier (on ne peut pas faire une demi-augmentation). On teste $x=2$ et $x=3$ :
$R(2) = -100(2-\frac{5}{2})^2 + 50625 = -25 + 50625 = 50600$ FCFA.
$R(3) = -100(3-\frac{5}{2})^2 + 50625 = -25 + 50625 = 50600$ FCFA.
Les deux prix donnent la même recette maximale entière :
\begin{itemize}
\item $x=2$ : Prix $= 1000 + 100 = 1100$ FCFA, clients $= 46$, Recette $= 50600$ FCFA.
\item $x=3$ : Prix $= 1000 + 150 = 1150$ FCFA, clients $= 44$, Recette $= 50600$ FCFA.
\end{itemize}
\textbf{Conclusion :} Le vendeur fixera son prix à \textbf{1100 FCFA} ou \textbf{1150 FCFA} pour gagner \textbf{50\,600 FCFA} par jour.
\end{enumerate}
\end{exoresolu}

\begin{exemplebox}[Résolution d'équation par forme canonique : $x^2 - 6x + 7 = 0$]
\begin{align*}
x^2 - 6x + 7 &= 0 \\
(x - 3)^2 - 9 + 7 &= 0 \quad \text{(moitié de -6 est -3, carré 9)} \\
(x - 3)^2 - 2 &= 0 \\
(x - 3)^2 &= 2 \\
x - 3 &= \pm \sqrt{2} \\
x &= 3 \pm \sqrt{2}.
\end{align*}
Les solutions sont $S = \{3-\sqrt{2}\,;\, 3+\sqrt{2}\}$.
\emph{Remarque :} $\Delta = 36 - 28 = 8$, $\sqrt{\Delta} = 2\sqrt{2}$, formules classiques : $\frac{6 \pm 2\sqrt{2}}{2} = 3 \pm \sqrt{2}$. Même résultat, mais la forme canonique évite la formule « magique » et donne le sommet $S(3\,;\,-2)$ gratuitement.
\end{exemplebox}

\courssub{Exercices d'entraînement}

\begin{exercice}[Mise en forme canonique et lecture]
Écrire les polynômes suivants sous forme canonique. En déduire le sommet de leur parabole, leur axe de symétrie, leurs variations et leur extremum.
\begin{enumerate}
\item $P_1(x) = x^2 + 6x - 2$
\item $P_2(x) = -2x^2 + 8x + 1$
\item $P_3(x) = 3x^2 - 5x + 4$
\end{enumerate}
\end{exercice}

\begin{exercice}[Résolution d'équations et d'inéquations]
Résoudre dans $\mathbb{R}$ en utilisant la forme canonique :
\begin{enumerate}
\item $x^2 - 4x - 1 = 0$
\item $2x^2 + 4x - 1 \le 0$
\item $-x^2 + 2x + 3 > 0$
\end{enumerate}
\end{exercice}

\begin{exercice}[Problème concret : Clôture d'un champ]
Un agriculteur dispose de $200$ mètres de grillage pour clôturer un champ rectangulaire le long d'une rivière (donc seulement trois côtés à clôturer : deux largeurs et une longueur). Quelle doit être la longueur du côté parallèle à la rivière pour maximiser la surface du champ ? Quelle est cette surface maximale ?
\end{exercice}

\begin{corrige}[Exercice 1]
\begin{enumerate}
\item $P_1(x) = x^2 + 6x - 2 = (x+3)^2 - 9 - 2 = (x+3)^2 - 11$.
Sommet $S_1(-3\,;\,-11)$, axe $x=-3$, $a=1>0$ : min $-11$ en $-3$. Décroît sur $]-\infty,-3]$, croît sur $[-3,+\infty[$.
\item $P_2(x) = -2x^2 + 8x + 1 = -2(x^2 - 4x) + 1 = -2[(x-2)^2 - 4] + 1 = -2(x-2)^2 + 9$.
Sommet $S_2(2\,;\,9)$, axe $x=2$, $a=-2<0$ : max $9$ en $2$. Croît sur $]-\infty,2]$, décroît sur $[2,+\infty[$.
\item $P_3(x) = 3x^2 - 5x + 4 = 3(x^2 - \frac{5}{3}x) + 4 = 3[(x-\frac{5}{6})^2 - \frac{25}{36}] + 4 = 3(x-\frac{5}{6})^2 - \frac{25}{12} + \frac{48}{12} = 3(x-\frac{5}{6})^2 + \frac{23}{12}$.
Sommet $S_3(\frac{5}{6}\,;\,\frac{23}{12})$, axe $x=\frac{5}{6}$, $a=3>0$ : min $\frac{23}{12}$ en $\frac{5}{6}$. Décroît sur $]-\infty,\frac{5}{6}]$, croît sur $[\frac{5}{6},+\infty[$.
\end{enumerate}
\end{corrige}

\begin{corrige}[Exercice 2]
\begin{enumerate}
\item $x^2 - 4x - 1 = 0 \Leftrightarrow (x-2)^2 - 5 = 0 \Leftrightarrow (x-2)^2 = 5 \Leftrightarrow x-2 = \pm\sqrt{5} \Leftrightarrow x = 2\pm\sqrt{5}$. $S = \{2-\sqrt{5}\,;\,2+\sqrt{5}\}$.
\item $2x^2 + 4x - 1 \le 0 \Leftrightarrow 2(x^2+2x) - 1 \le 0 \Leftrightarrow 2[(x+1)^2 - 1] - 1 \le 0 \Leftrightarrow 2(x+1)^2 - 3 \le 0 \Leftrightarrow (x+1)^2 \le \frac{3}{2} \Leftrightarrow |x+1| \le \sqrt{\frac{3}{2}} = \frac{\sqrt{6}}{2}$.
$S = \left[-1-\frac{\sqrt{6}}{2}\,;\,-1+\frac{\sqrt{6}}{2}\right]$.
\item $-x^2 + 2x + 3 > 0 \Leftrightarrow -(x^2-2x) + 3 > 0 \Leftrightarrow -[(x-1)^2 - 1] + 3 > 0 \Leftrightarrow -(x-1)^2 + 4 > 0 \Leftrightarrow (x-1)^2 < 4 \Leftrightarrow |x-1| < 2 \Leftrightarrow -2 < x-1 < 2 \Leftrightarrow -1 < x < 3$.
$S = ]-1\,;\,3[$.
\end{enumerate}
\end{corrige}

\begin{corrige}[Exercice 3]
Soit $x$ la longueur du côté parallèle à la rivière (en mètres). Les deux largeurs valent $\frac{200-x}{2}$ chacune.
Surface $A(x) = x \times \frac{200-x}{2} = -\frac{1}{2}x^2 + 100x$.
Forme canonique : $A(x) = -\frac{1}{2}(x^2 - 200x) = -\frac{1}{2}[(x-100)^2 - 10000] = -\frac{1}{2}(x-100)^2 + 5000$.
$a = -\frac{1}{2} < 0$, maximum au sommet $x=100$ m. Surface maximale $= 5000$ m$^2$.
Les largeurs valent $\frac{200-100}{2} = 50$ m.
\textbf{Réponse :} Longueur $100$ m (côté rivière), largeurs $50$ m. Surface max $5000$ m$^2$.
\end{corrige}

\coursec{Discriminant et résolution des équations du second degré}

Dans la section précédente, tu as appris à écrire tout polynôme du second degré $P(x)=ax^2+bx+c$ (avec $a\neq 0$) sous sa \cle{forme canonique} :
\[
P(x)=a\left[\left(x+\dfrac{b}{2a}\right)^2-\dfrac{b^2-4ac}{4a^2}\right].
\]
Cette écriture est un véritable trésor : regarde-la attentivement. Le signe de $P(x)$, l'existence de solutions de l'équation $P(x)=0$, tout semble dépendre d'une seule quantité, celle qui apparaît au numérateur du dernier terme : $b^2-4ac$. Cette quantité porte un nom : le \cle{discriminant}. Toute cette section va montrer qu'elle « discrimine » effectivement, c'est-à-dire qu'elle permet de distinguer, sans aucun calcul de racine carrée inutile, les trois situations possibles : deux solutions, une seule, ou aucune.

\courssub{Le discriminant : définition et intuition}

\begin{definition}[Discriminant d'un polynôme du second degré]
Soit $P(x)=ax^2+bx+c$ un polynôme du second degré, avec $a$, $b$, $c$ réels et $a\neq 0$. On appelle \cle{discriminant} de $P$ le réel noté $\Delta$ (lettre grecque « delta ») défini par :
\[
\Delta=b^2-4ac.
\]
\end{definition}

\begin{attention}[Bien identifier $a$, $b$ et $c$]
Avant de calculer $\Delta$, l'équation doit être écrite sous la forme $ax^2+bx+c=0$, avec \textbf{tous les termes du même côté}. Par exemple, pour l'équation $2x^2-5x=-3$, on ne peut pas lire directement les coefficients : on écrit d'abord $2x^2-5x+3=0$, d'où $a=2$, $b=-5$, $c=3$. N'oublie jamais les signes : $b=-5$ et non $5$.
\end{attention}

\begin{exemplebox}[Calculs de discriminants]
\begin{itemize}
\item Pour $2x^2-5x+3$ : $\Delta=(-5)^2-4\times 2\times 3=25-24=1$.
\item Pour $x^2+x+1$ : $\Delta=1^2-4\times 1\times 1=1-4=-3$.
\item Pour $x^2-6x+9$ : $\Delta=(-6)^2-4\times 1\times 9=36-36=0$.
\end{itemize}
\end{exemplebox}

\courssub{Le théorème fondamental de résolution}

\begin{propriete}[Résolution de l'équation $ax^2+bx+c=0$ dans $\mathbb{R}$ — démonstration exigible]
Soit $a$, $b$, $c$ trois réels avec $a\neq 0$, et soit $\Delta=b^2-4ac$. L'équation $(E)$ : $ax^2+bx+c=0$ admet dans $\mathbb{R}$ :
\begin{itemize}
\item si $\Delta>0$ : \textbf{deux solutions distinctes}
\[
x_1=\dfrac{-b-\sqrt{\Delta}}{2a}\qquad\text{et}\qquad x_2=\dfrac{-b+\sqrt{\Delta}}{2a}\,;
\]
\item si $\Delta=0$ : \textbf{une unique solution} (dite \cle{racine double}) $x_0=-\dfrac{b}{2a}$ ;
\item si $\Delta<0$ : \textbf{aucune solution réelle}.
\end{itemize}
\end{propriete}

Cette démonstration est exigible au Probatoire : elle est courte, élégante, et repose entièrement sur la forme canonique. Suivons-la pas à pas.

\begin{experience}[Démonstration du théorème à partir de la forme canonique]
\textbf{Étape 1 : forme canonique.} On a vu que pour tout réel $x$ :
\[
ax^2+bx+c=a\left[\left(x+\dfrac{b}{2a}\right)^2-\dfrac{\Delta}{4a^2}\right].
\]
Comme $a\neq 0$, l'équation $ax^2+bx+c=0$ équivaut à :
\[
\left(x+\dfrac{b}{2a}\right)^2=\dfrac{\Delta}{4a^2}.
\]
\textbf{Étape 2 : discussion selon le signe de $\Delta$.} Le membre de gauche est un carré, donc toujours positif ou nul. Le membre de droite a le signe de $\Delta$ (car $4a^2>0$).
\begin{itemize}
\item \textbf{Si $\Delta<0$} : un carré ne peut pas être égal à un nombre strictement négatif. Il n'y a \textbf{aucune solution réelle}.
\item \textbf{Si $\Delta=0$} : l'équation devient $\left(x+\dfrac{b}{2a}\right)^2=0$, donc $x+\dfrac{b}{2a}=0$, d'où l'unique solution $x_0=-\dfrac{b}{2a}$.
\item \textbf{Si $\Delta>0$} : un carré est égal à un nombre strictement positif dans deux cas :
\[
x+\dfrac{b}{2a}=\dfrac{\sqrt{\Delta}}{2a}\quad\text{ou}\quad x+\dfrac{b}{2a}=-\dfrac{\sqrt{\Delta}}{2a},
\]
ce qui donne exactement $x_1=\dfrac{-b-\sqrt{\Delta}}{2a}$ et $x_2=\dfrac{-b+\sqrt{\Delta}}{2a}$.
\end{itemize}
\textbf{Étape 3 : conclusion.} Le signe de $\Delta$ détermine donc entièrement le nombre de solutions : deux si $\Delta>0$, une si $\Delta=0$, aucune si $\Delta<0$. $\square$
\end{experience}

\begin{attention}[Les deux erreurs classiques dans les formules]
\begin{itemize}
\item Le dénominateur est $2a$, pas $2$ ni $a$ : écrire $x=\dfrac{-b\pm\sqrt{\Delta}}{2}$ est faux dès que $a\neq 1$.
\item Le signe de $a$ compte : si $a<0$, les formules restent valables telles quelles, à condition de garder $2a$ (négatif) au dénominateur. Ne « corrige » rien à la main.
\end{itemize}
\end{attention}

\begin{exoresolu}[Résoudre $2x^2-5x+3=0$ et $x^2+x+1=0$]
Résoudre dans $\mathbb{R}$ les équations : (1) $2x^2-5x+3=0$ ; (2) $x^2+x+1=0$.
\tcblower
\textbf{Équation (1).} Ici $a=2$, $b=-5$, $c=3$.
\[
\Delta=(-5)^2-4\times 2\times 3=25-24=1>0.
\]
L'équation admet deux solutions :
\begin{align*}
x_1&=\dfrac{-(-5)-\sqrt{1}}{2\times 2}=\dfrac{5-1}{4}=\dfrac{4}{4}=1,\\
x_2&=\dfrac{-(-5)+\sqrt{1}}{2\times 2}=\dfrac{5+1}{4}=\dfrac{6}{4}=\dfrac{3}{2}.
\end{align*}
Donc $S=\left\{1\,;\,\dfrac{3}{2}\right\}$. Vérification : $2(1)^2-5(1)+3=0$ et $2\left(\tfrac32\right)^2-5\left(\tfrac32\right)+3=\tfrac92-\tfrac{15}{2}+3=0$. \checkmark

\textbf{Équation (2).} Ici $a=1$, $b=1$, $c=1$.
\[
\Delta=1^2-4\times 1\times 1=-3<0.
\]
L'équation n'admet \textbf{aucune solution réelle} : $S=\emptyset$.
\end{exoresolu}

\begin{exemplebox}[Cas $\Delta=0$ : une racine double]
Résolvons $x^2-6x+9=0$. On a $\Delta=36-36=0$, donc une unique solution $x_0=-\dfrac{-6}{2\times 1}=3$. On reconnaît d'ailleurs l'identité remarquable : $x^2-6x+9=(x-3)^2$. Quand $\Delta=0$, le trinôme est toujours un carré parfait (au facteur $a$ près) : $ax^2+bx+c=a\left(x+\dfrac{b}{2a}\right)^2$.
\end{exemplebox}

\courssub{Interprétation graphique du signe de $\Delta$}

Géométriquement, résoudre $ax^2+bx+c=0$ revient à chercher les abscisses des points d'intersection de la parabole d'équation $y=ax^2+bx+c$ avec l'axe des abscisses. Le signe de $\Delta$ se lit alors directement sur le graphique : $\Delta>0$ correspond à deux points d'intersection, $\Delta=0$ à une parabole tangente à l'axe, $\Delta<0$ à une parabole qui ne rencontre jamais l'axe.

\begin{popfigure}
\begin{center}
\begin{tikzpicture}[scale=0.85]
\begin{axis}[width=0.36\linewidth,axis lines=middle,xlabel={$x$},ylabel={$y$},xtick=\empty,ytick=\empty,title={$\Delta>0$},title style={color=popblue},domain=-1:3,samples=100]
\addplot[popcurve]{x^2-2*x-0.5};
\addplot[popink,mark=*,only marks,mark size=1.5pt] coordinates {(-0.225,0) (2.225,0)};
\end{axis}
\end{tikzpicture}\hfill
\begin{tikzpicture}[scale=0.85]
\begin{axis}[width=0.36\linewidth,axis lines=middle,xlabel={$x$},ylabel={$y$},xtick=\empty,ytick=\empty,title={$\Delta=0$},title style={color=poporange},domain=-1:3,samples=100]
\addplot[popcurve]{(x-1)^2};
\addplot[popink,mark=*,only marks,mark size=1.5pt] coordinates {(1,0)};
\end{axis}
\end{tikzpicture}\hfill
\begin{tikzpicture}[scale=0.85]
\begin{axis}[width=0.36\linewidth,axis lines=middle,xlabel={$x$},ylabel={$y$},xtick=\empty,ytick=\empty,title={$\Delta<0$},title style={color=poppurple},domain=-1:3,samples=100]
\addplot[popcurve]{(x-1)^2+0.8};
\end{axis}
\end{tikzpicture}
\end{center}
\legende{Les trois positions possibles de la parabole par rapport à l'axe des abscisses : deux intersections ($\Delta>0$), un point de tangence ($\Delta=0$), aucune intersection ($\Delta<0$).}
\end{popfigure}

\begin{aretenir}[Stratégie de résolution : l'arbre de décision]
Face à une équation du second degré, la démarche est toujours la même : écrire l'équation sous la forme $ax^2+bx+c=0$, calculer $\Delta=b^2-4ac$, puis conclure selon son signe.
\end{aretenir}

\begin{popfigure}
\begin{center}
\begin{tikzpicture}[
node distance=0.9cm and 1.1cm,
etape/.style={rectangle,rounded corners,draw=popblue,fill=popblueL,thick,text width=4.6cm,align=center,font=\small},
branche/.style={rectangle,rounded corners,draw=popdark,thick,text width=3.6cm,align=center,font=\small},
fleche/.style={-{Stealth[length=2.5mm]},thick,popdark}]
\node[etape] (e1) {Écrire l'équation sous la forme $ax^2+bx+c=0$};
\node[etape,below=of e1] (e2) {Calculer $\Delta=b^2-4ac$};
\node[branche,below left=of e2,fill=popgreenL,draw=popgreen, align=center] (pos) {$\Delta>0$\\ deux solutions\\ $x_{1,2}=\dfrac{-b\pm\sqrt{\Delta}}{2a}$};
\node[branche,below=of e2,fill=poporangeL,draw=poporange, align=center] (nul) {$\Delta=0$\\ une solution double\\ $x_0=-\dfrac{b}{2a}$};
\node[branche,below right=of e2,fill=poppinkL,draw=poppink, align=center] (neg) {$\Delta<0$\\ aucune solution réelle\\ $S=\emptyset$};
\draw[fleche] (e1) -- (e2);
\draw[fleche] (e2) -- (pos);
\draw[fleche] (e2) -- (nul);
\draw[fleche] (e2) -- (neg);
\end{tikzpicture}
\end{center}
\legende{Arbre de décision pour la résolution de $ax^2+bx+c=0$ dans $\mathbb{R}$.}
\end{popfigure}

\courssub{Somme et produit des racines}

Lorsque $\Delta\geq 0$, les racines vérifient deux relations remarquables, très utiles pour vérifier un calcul, factoriser rapidement, ou résoudre les systèmes de la section 5.

\begin{propriete}[Somme et produit des racines]
Si l'équation $ax^2+bx+c=0$ ($a\neq 0$) admet deux racines réelles $x_1$ et $x_2$ (éventuellement égales), alors :
\[
S=x_1+x_2=-\dfrac{b}{a}\qquad\text{et}\qquad P=x_1\times x_2=\dfrac{c}{a}.
\]
En particulier, deux réels dont on connaît la somme $S$ et le produit $P$ sont les solutions de l'équation $x^2-Sx+P=0$.
\end{propriete}

\begin{experience}[Démonstration guidée]
Supposons $\Delta>0$ et posons $x_1=\dfrac{-b-\sqrt{\Delta}}{2a}$ et $x_2=\dfrac{-b+\sqrt{\Delta}}{2a}$.
\begin{enumerate}
\item \textbf{Somme.} En additionnant, les termes en $\sqrt{\Delta}$ s'éliminent :
\[
x_1+x_2=\dfrac{-b-\sqrt{\Delta}}{2a}+\dfrac{-b+\sqrt{\Delta}}{2a}=\dfrac{-2b}{2a}=-\dfrac{b}{a}.
\]
\item \textbf{Produit.} On utilise l'identité $(u-v)(u+v)=u^2-v^2$ avec $u=-b$ et $v=\sqrt{\Delta}$ :
\[
x_1\,x_2=\dfrac{(-b)^2-(\sqrt{\Delta})^2}{(2a)^2}=\dfrac{b^2-\Delta}{4a^2}=\dfrac{b^2-(b^2-4ac)}{4a^2}=\dfrac{4ac}{4a^2}=\dfrac{c}{a}.
\]
\item \textbf{Cas $\Delta=0$.} La racine double $x_0=-\dfrac{b}{2a}$ vérifie bien $x_0+x_0=-\dfrac{b}{a}$ et $x_0^2=\dfrac{b^2}{4a^2}=\dfrac{4ac}{4a^2}=\dfrac{c}{a}$ puisque $b^2=4ac$.
\end{enumerate}
$\square$
\end{experience}

\begin{exemplebox}[Vérification rapide]
Pour l'équation $2x^2-5x+3=0$ résolue plus haut : $S=-\dfrac{-5}{2}=\dfrac52$ et $P=\dfrac32$. Or $1+\dfrac32=\dfrac52$ et $1\times\dfrac32=\dfrac32$ : nos solutions $1$ et $\dfrac32$ sont confirmées en deux lignes, sans refaire tout le calcul.
\end{exemplebox}

\courssub{Équations se ramenant au second degré}

Beaucoup d'équations ne se présentent pas directement sous la forme $ax^2+bx+c=0$, mais s'y ramènent. Trois situations classiques sont au programme.

\textbf{1. Équations produits.} Un produit de facteurs est nul si et seulement si l'un au moins des facteurs est nul. On factorise, puis on résout chaque facteur.

\begin{exoresolu}[Équation produit]
Résoudre dans $\mathbb{R}$ : $(x^2-3x+2)(2x+1)=0$.
\tcblower
Un produit est nul si l'un des facteurs est nul.
\begin{itemize}
\item $2x+1=0$ donne $x=-\dfrac12$.
\item $x^2-3x+2=0$ : $\Delta=9-8=1>0$, d'où $x_1=\dfrac{3-1}{2}=1$ et $x_2=\dfrac{3+1}{2}=2$.
\end{itemize}
Conclusion : $S=\left\{-\dfrac12\,;\,1\,;\,2\right\}$.
\end{exoresolu}

\textbf{2. Équations bicarrées.} Une équation \cle{bicarrée} est une équation de la forme $ax^4+bx^2+c=0$ : elle ne contient que des puissances paires de $x$.

\begin{methode}[Résoudre une équation bicarrée $ax^4+bx^2+c=0$]
\begin{enumerate}
\item Poser le changement de variable $X=x^2$, avec la condition essentielle $X\geq 0$.
\item Résoudre l'équation du second degré $aX^2+bX+c=0$ en calculant son discriminant.
\item \textbf{Écarter toute solution $X<0$}, car un carré $x^2$ ne peut pas être négatif.
\item Pour chaque solution $X\geq 0$ retenue, résoudre $x^2=X$ : si $X>0$, on obtient $x=\sqrt{X}$ et $x=-\sqrt{X}$ ; si $X=0$, on obtient $x=0$.
\item Conclure en donnant l'ensemble $S$ de toutes les solutions réelles.
\end{enumerate}
\end{methode}

\begin{exoresolu}[Équation bicarrée]
Résoudre dans $\mathbb{R}$ : $x^4-5x^2+4=0$.
\tcblower
Posons $X=x^2$ ($X\geq 0$). L'équation devient $X^2-5X+4=0$.
\[
\Delta=25-16=9>0,\qquad X_1=\dfrac{5-3}{2}=1,\qquad X_2=\dfrac{5+3}{2}=4.
\]
Les deux valeurs sont positives, on les garde toutes les deux.
\begin{itemize}
\item $x^2=1$ donne $x=1$ ou $x=-1$ ;
\item $x^2=4$ donne $x=2$ ou $x=-2$.
\end{itemize}
Conclusion : $S=\{-2\,;\,-1\,;\,1\,;\,2\}$. Vérification pour $x=2$ : $16-20+4=0$. \checkmark
\end{exoresolu}

\begin{attention}[Le piège du changement de variable]
Si le second degré en $X$ donne une solution négative, elle ne fournit \textbf{aucune} solution en $x$. Par exemple, $x^4+3x^2+2=0$ donne $X^2+3X+2=0$, soit $X=-1$ ou $X=-2$ : toutes deux négatives, donc $S=\emptyset$. Écrire « $x^2=-1$ donc $x=\pm\sqrt{-1}$ » est une erreur grave dans $\mathbb{R}$.
\end{attention}

\textbf{3. Équations avec dénominateur.} On commence toujours par déterminer les \cle{valeurs interdites} (celles qui annulent un dénominateur), puis on se ramène à une équation du second degré en réduisant au même dénominateur, et enfin on \textbf{élimine} les solutions interdites.

\begin{exoresolu}[Équation avec dénominateur]
Résoudre dans $\mathbb{R}$ : $\dfrac{x+1}{x-1}=x-2$.
\tcblower
\textbf{Valeur interdite :} $x-1=0\Leftrightarrow x=1$. On cherche donc les solutions dans $\mathbb{R}\setminus\{1\}$.

Pour $x\neq 1$, on multiplie les deux membres par $x-1$ :
\[
x+1=(x-2)(x-1)\iff x+1=x^2-3x+2\iff x^2-4x+1=0.
\]
Discriminant : $\Delta=16-4=12>0$, d'où
\[
x_1=\dfrac{4-\sqrt{12}}{2}=\dfrac{4-2\sqrt{3}}{2}=2-\sqrt{3},\qquad x_2=2+\sqrt{3}.
\]
Aucune de ces deux valeurs n'est égale à $1$ : elles sont toutes deux acceptables.
\[
S=\{2-\sqrt{3}\,;\,2+\sqrt{3}\}.
\]
\end{exoresolu}

\begin{attention}[Ne jamais oublier les valeurs interdites]
Avant toute manipulation d'une équation contenant $x$ au dénominateur, écris explicitement l'ensemble de résolution (par exemple $\mathbb{R}\setminus\{1\}$). À la fin, compare chaque solution trouvée aux valeurs interdites : une solution qui annule le dénominateur doit être rejetée, même si elle est correctement calculée. C'est une source fréquente de perte de points au Probatoire.
\end{attention}

\courssub{Ce qui est hors programme — et ce qu'il faut retenir}

\begin{attention}[Limites officielles du programme de Première]
Conformément au programme officiel :
\begin{itemize}
\item le \textbf{discriminant réduit} $\Delta'=b'^2-ac$ (avec $b=2b'$) est \textbf{hors programme} : utilise toujours $\Delta=b^2-4ac$ ;
\item la résolution algébrique des équations du second degré \textbf{avec paramètres} (par exemple discuter selon $m$ le nombre de solutions de $mx^2+x+1=0$) est \textbf{hors programme} ;
\item les systèmes d'équations avec paramètres sont également hors programme.
\end{itemize}
Tout exercice au Probatoire portera donc sur des coefficients numériques explicites.
\end{attention}

\begin{aretenir}[L'essentiel de la section]
\begin{itemize}
\item Le discriminant de $ax^2+bx+c$ ($a\neq 0$) est $\Delta=b^2-4ac$.
\item $\Delta>0$ : deux solutions $x_{1,2}=\dfrac{-b\pm\sqrt{\Delta}}{2a}$ ; $\Delta=0$ : une solution double $x_0=-\dfrac{b}{2a}$ ; $\Delta<0$ : aucune solution réelle.
\item Graphiquement, $\Delta$ compte les intersections de la parabole avec l'axe des abscisses.
\item Si $x_1$ et $x_2$ sont les racines : $x_1+x_2=-\dfrac{b}{a}$ et $x_1x_2=\dfrac{c}{a}$.
\item Équations bicarrées : poser $X=x^2$ avec $X\geq 0$ et écarter les solutions négatives en $X$.
\item Équations à dénominateur : valeurs interdites d'abord, vérification à la fin.
\end{itemize}
\end{aretenir}

\begin{savaistu}[Pourquoi « discriminant » ?]
Le mot vient du latin \emph{discriminare}, « distinguer, séparer ». Introduit dans ce rôle au XIX\textsuperscript{e} siècle, le discriminant « distingue » effectivement les cas : il dit, avant tout calcul, combien de racines réelles possède l'équation. Les Babyloniens, il y a près de 4000 ans, savaient déjà résoudre des équations du second degré par des méthodes géométriques de « complétion du carré » — exactement l'idée de la forme canonique que tu as utilisée dans la démonstration. Au Cameroun, ces équations apparaissent dès qu'on modélise : trajectoire d'un ballon, rentabilité d'un champ de maïs en fonction de sa surface, dimensionnement d'une antenne CAMTEL... Tu les retrouveras dans la section 6 consacrée aux problèmes concrets.
\end{savaistu}

\begin{exercice}[Pour t'entraîner]
\begin{enumerate}
\item Résoudre dans $\mathbb{R}$ : $3x^2+7x-6=0$ ; $4x^2-4x+1=0$ ; $x^2+2x+5=0$.
\item Sans résoudre, déterminer le nombre de solutions réelles de $5x^2-3x+2=0$.
\item Résoudre l'équation bicarrée $x^4-13x^2+36=0$.
\item Résoudre dans $\mathbb{R}$ : $\dfrac{2x}{x+2}=x-1$ (penser aux valeurs interdites).
\item Trouver deux réels dont la somme vaut $7$ et le produit $12$.
\end{enumerate}
\end{exercice}

\begin{corrige}[Exercice : pour t'entraîner]
\begin{enumerate}
\item $3x^2+7x-6=0$ : $\Delta=49+72=121=11^2$, $x_1=\dfrac{-7-11}{6}=-3$, $x_2=\dfrac{-7+11}{6}=\dfrac23$. $S=\left\{-3\,;\,\dfrac23\right\}$.\\
$4x^2-4x+1=0$ : $\Delta=16-16=0$, $x_0=\dfrac{4}{8}=\dfrac12$. $S=\left\{\dfrac12\right\}$ (on reconnaît $(2x-1)^2$).\\
$x^2+2x+5=0$ : $\Delta=4-20=-16<0$, $S=\emptyset$.
\item $\Delta=9-40=-31<0$ : aucune solution réelle.
\item $X=x^2$ : $X^2-13X+36=0$, $\Delta=169-144=25$, $X_1=4$, $X_2=9$, toutes deux positives. $S=\{-3\,;\,-2\,;\,2\,;\,3\}$.
\item Valeur interdite : $x=-2$. Pour $x\neq -2$ : $2x=(x-1)(x+2)=x^2+x-2$, soit $x^2-x-2=0$, $\Delta=9$, $x_1=-1$, $x_2=2$. Aucune n'est interdite : $S=\{-1\,;\,2\}$.
\item Ce sont les solutions de $x^2-7x+12=0$ : $\Delta=49-48=1$, $x_1=3$, $x_2=4$. Les réels cherchés sont $3$ et $4$.
\end{enumerate}
\end{corrige}

\coursec{Factorisation et signe d'un polynôme du second degré}

Dans la section précédente, nous avons appris à résoudre l'équation $ax^2+bx+c=0$ grâce au discriminant $\Delta=b^2-4ac$ : nous savons désormais dire si un trinôme possède des racines dans $\mathbb{R}$ et les calculer. Mais une question naturelle se pose : une fois les racines connues, peut-on \emph{réécrire} le polynôme autrement ? Et surtout, peut-on savoir \emph{où} ce polynôme est positif ou négatif, sans calculer sa valeur en chaque point ? C'est précisément l'objet de cette section : la \cle{factorisation} du trinôme et l'étude de son \cle{signe}, deux outils indispensables pour résoudre les inéquations du second degré qui t'attendent au Probatoire.

\courssub{Factorisation d'un polynôme du second degré}

\textbf{Intuition.}\par
Quand tu connais les solutions d'une équation, tu connais en réalité les points où le polynôme « s'annule ». Or, en classe de Seconde, tu as vu qu'un polynôme qui s'annule en $x_1$ est divisible par $(x-x_1)$. Pour un trinôme qui s'annule deux fois, on peut donc espérer l'écrire comme un produit de deux facteurs du premier degré. C'est exactement ce que garantit le théorème suivant.

\begin{propriete}[Factorisation du trinôme selon le signe de $\Delta$]
Soit $P(x)=ax^2+bx+c$ un polynôme du second degré ($a\neq 0$) de discriminant $\Delta=b^2-4ac$.
\begin{itemize}
\item Si $\Delta>0$, $P$ admet deux racines réelles $x_1$ et $x_2$, et pour tout réel $x$ :
\[ P(x)=a(x-x_1)(x-x_2). \]
\item Si $\Delta=0$, $P$ admet une racine double $x_0=-\dfrac{b}{2a}$, et pour tout réel $x$ :
\[ P(x)=a(x-x_0)^2. \]
\item Si $\Delta<0$, $P$ \textbf{ne peut pas} se factoriser en produit de deux polynômes du premier degré dans $\mathbb{R}$.
\end{itemize}
Cette démonstration est exigible : elle repose uniquement sur la forme canonique.
\end{propriete}

\begin{experience}[Démonstration de la factorisation à partir de la forme canonique]
Rappelons la forme canonique obtenue à la section 2 :
\[ P(x)=a\left[\left(x+\frac{b}{2a}\right)^2-\frac{\Delta}{4a^2}\right]. \]
\textbf{Étape 1 : cas $\Delta>0$.} Comme $\Delta>0$, on peut écrire $\Delta=(\sqrt{\Delta})^2$, donc :
\[ P(x)=a\left[\left(x+\frac{b}{2a}\right)^2-\left(\frac{\sqrt{\Delta}}{2a}\right)^2\right]. \]
On reconnaît l'identité remarquable $A^2-B^2=(A-B)(A+B)$ avec $A=x+\dfrac{b}{2a}$ et $B=\dfrac{\sqrt{\Delta}}{2a}$ :
\[ P(x)=a\left(x+\frac{b}{2a}-\frac{\sqrt{\Delta}}{2a}\right)\left(x+\frac{b}{2a}+\frac{\sqrt{\Delta}}{2a}\right). \]
Or $x_1=\dfrac{-b-\sqrt{\Delta}}{2a}$ et $x_2=\dfrac{-b+\sqrt{\Delta}}{2a}$, donc $x+\dfrac{b}{2a}-\dfrac{\sqrt{\Delta}}{2a}=x-x_1$ et $x+\dfrac{b}{2a}+\dfrac{\sqrt{\Delta}}{2a}=x-x_2$. Ainsi :
\[ P(x)=a(x-x_1)(x-x_2). \]
\textbf{Étape 2 : cas $\Delta=0$.} La forme canonique devient immédiatement $P(x)=a\left(x+\dfrac{b}{2a}\right)^2=a(x-x_0)^2$.
\textbf{Étape 3 : cas $\Delta<0$.} Alors $-\dfrac{\Delta}{4a^2}>0$, donc $\left(x+\dfrac{b}{2a}\right)^2-\dfrac{\Delta}{4a^2}>0$ pour tout $x$ : le trinôme ne s'annule jamais, il ne peut donc contenir aucun facteur du premier degré réel.
\end{experience}

\begin{exemplebox}[Factoriser avec les racines]
Factorisons $P(x)=2x^2-5x+3$. On calcule $\Delta=25-24=1>0$, d'où $x_1=\dfrac{5-1}{4}=1$ et $x_2=\dfrac{5+1}{4}=\dfrac{3}{2}$. Donc :
\[ P(x)=2(x-1)\left(x-\frac{3}{2}\right)=(x-1)(2x-3). \]
Vérifions en développant : $(x-1)(2x-3)=2x^2-3x-2x+3=2x^2-5x+3$. Parfait.
\end{exemplebox}

\begin{attention}[N'oublie pas le coefficient $a$ !]
L'erreur la plus fréquente consiste à écrire $P(x)=(x-x_1)(x-x_2)$ en oubliant le facteur $a$. Par exemple, $2x^2-5x+3$ n'est \textbf{pas} égal à $(x-1)\left(x-\frac{3}{2}\right)$ : il manque le facteur $2$. Réflexe de survie : développe mentalement ta factorisation et vérifie que le terme en $x^2$ redonne bien $ax^2$.
\end{attention}

\courssub{Racine évidente : trouver l'autre racine sans calculer $\Delta$}

Parfois, une racine « saute aux yeux » : $1$, $-1$, $2$... Plutôt que de calculer le discriminant, on exploite les relations entre les coefficients et les racines.

\begin{propriete}[Somme et produit des racines]
Si $\Delta\geq 0$, les racines $x_1$ et $x_2$ de $ax^2+bx+c$ vérifient :
\[ S=x_1+x_2=-\frac{b}{a} \qquad\text{et}\qquad P=x_1x_2=\frac{c}{a}. \]
En effet, en développant $a(x-x_1)(x-x_2)=ax^2-a(x_1+x_2)x+ax_1x_2$ et en identifiant avec $ax^2+bx+c$, on obtient $-a(x_1+x_2)=b$ et $ax_1x_2=c$.
\end{propriete}

\begin{methode}[Utiliser une racine évidente]
\begin{enumerate}
\item \textbf{Tester} les valeurs simples : $P(1)=a+b+c$, $P(-1)=a-b+c$, $P(2)=4a+2b+c$. Si l'une est nulle, tu as une racine $x_1$.
\item \textbf{En déduire l'autre racine} par le produit : $x_2=\dfrac{c}{a\,x_1}$, ou par la somme : $x_2=-\dfrac{b}{a}-x_1$.
\item \textbf{Écrire la factorisation} $P(x)=a(x-x_1)(x-x_2)$.
\end{enumerate}
Cas particulier précieux : si $a+b+c=0$, alors $x_1=1$ et $x_2=\dfrac{c}{a}$. Si $a-b+c=0$, alors $x_1=-1$ et $x_2=-\dfrac{c}{a}$.
\end{methode}

\begin{exemplebox}[Racine évidente en action]
Soit $P(x)=3x^2+2x-5$. On remarque que $3+2-5=0$ : donc $x_1=1$ est racine évidente. Le produit des racines vaut $\dfrac{c}{a}=-\dfrac{5}{3}$, donc $x_2=-\dfrac{5}{3}$. D'où :
\[ P(x)=3(x-1)\left(x+\frac{5}{3}\right)=(x-1)(3x+5). \]
Gain de temps considérable par rapport au calcul de $\Delta=64$ !
\end{exemplebox}

\courssub{Signe d'un polynôme du second degré}

\textbf{Intuition.}\par
La courbe d'un trinôme est une parabole. Si elle coupe l'axe des abscisses en deux points (cas $\Delta>0$), elle se trouve d'un côté de l'axe entre les deux points de croisement, et de l'autre côté à l'extérieur. Quand $a>0$, les branches montent : la parabole est \emph{sous} l'axe entre les racines et \emph{au-dessus} à l'extérieur. Le théorème du signe du trinôme n'est que la traduction algébrique de cette image.

\begin{propriete}[Signe du trinôme $ax^2+bx+c$]
Soit $P(x)=ax^2+bx+c$ avec $a\neq 0$.
\begin{itemize}
\item Si $\Delta>0$ (racines $x_1<x_2$) : $P(x)$ est du \cle{signe de $a$} à l'extérieur de l'intervalle $[x_1\,;\,x_2]$, et du \cle{signe de $-a$} entre les racines. Il s'annule en $x_1$ et $x_2$.
\item Si $\Delta=0$ (racine double $x_0$) : $P(x)$ est du signe de $a$ pour tout $x\neq x_0$, et s'annule en $x_0$.
\item Si $\Delta<0$ : $P(x)$ est \textbf{strictement du signe de $a$} pour tout réel $x$.
\end{itemize}
\end{propriete}

\begin{experience}[Démonstration guidée du signe du trinôme (cas $\Delta>0$)]
\textbf{Étape 1.} Grâce à la factorisation, $P(x)=a(x-x_1)(x-x_2)$ avec $x_1<x_2$.
\textbf{Étape 2.} Étudions le signe de chaque facteur : $x-x_1$ est négatif avant $x_1$, positif après ; $x-x_2$ est négatif avant $x_2$, positif après.
\textbf{Étape 3.} Dressons le tableau de signes du produit $(x-x_1)(x-x_2)$ :
\begin{center}
\begin{tabular}{|c|ccccccc|}\hline
$x$ & $-\infty$ & & $x_1$ & & $x_2$ & & $+\infty$ \\ \hline
$x-x_1$ & & $-$ & $0$ & $+$ & & $+$ & \\ \hline
$x-x_2$ & & $-$ & & $-$ & $0$ & $+$ & \\ \hline
$(x-x_1)(x-x_2)$ & & $+$ & $0$ & $-$ & $0$ & $+$ & \\ \hline
\end{tabular}
\end{center}
\textbf{Étape 4.} Multiplier par $a$ : si $a>0$, le signe est conservé ; si $a<0$, il est renversé. Dans les deux cas, $P(x)$ est du signe de $a$ à l'extérieur des racines et du signe contraire entre les racines. $\blacksquare$
\par
Pour $\Delta=0$ : $P(x)=a(x-x_0)^2$ et un carré est toujours positif, donc $P$ est du signe de $a$ sauf en $x_0$. Pour $\Delta<0$ : la forme canonique montre que $P(x)=a\times(\text{quantité strictement positive})$, donc $P$ garde le signe de $a$ partout.
\end{experience}

\textbf{Les trois tableaux de signes à connaître par cœur.}\par
Cas $\Delta>0$, avec $a>0$ (si $a<0$, renverser tous les signes) :
\begin{center}
\begin{tabular}{|c|ccccccc|}\hline
$x$ & $-\infty$ & & $x_1$ & & $x_2$ & & $+\infty$ \\ \hline
$P(x)$ & & $+$ & $0$ & $-$ & $0$ & $+$ & \\ \hline
\end{tabular}
\end{center}
Cas $\Delta=0$ (avec $a>0$) :
\begin{center}
\begin{tabular}{|c|ccccc|}\hline
$x$ & $-\infty$ & & $x_0$ & & $+\infty$ \\ \hline
$P(x)$ & & $+$ & $0$ & $+$ & \\ \hline
\end{tabular}
\end{center}
Cas $\Delta<0$ (avec $a>0$) : $P(x)>0$ pour tout réel $x$.

\begin{popfigure}
\begin{center}
\begin{tikzpicture}
\begin{axis}[
  width=0.85\linewidth, height=6.5cm,
  axis lines=middle, xlabel={$x$}, ylabel={$y$},
  xtick={-1,4}, ytick=\empty,
  domain=-3:6, samples=200,
  xmin=-3, xmax=6, ymin=-8, ymax=8,
  every axis x label/.style={at={(ticklabel* cs:1)}, anchor=west},
]
\addplot[popcurve]{-x^2+3*x+4};
\addplot[draw=none, fill=popgreenL, domain=-1:4] {-x^2+3*x+4} \closedcycle;
\addplot[draw=none, fill=poppinkL, domain=-3:-1] {-x^2+3*x+4} \closedcycle;
\addplot[draw=none, fill=poppinkL, domain=4:6] {-x^2+3*x+4} \closedcycle;
\node[popgreen, font=\small] at (axis cs:1.5,4.5) {$P(x)>0$};
\node[poppink, font=\small] at (axis cs:-2.2,-4) {$P(x)<0$};
\node[poppink, font=\small] at (axis cs:5.2,-4) {$P(x)<0$};
\end{axis}
\end{tikzpicture}
\end{center}
\legende{Parabole de $P(x)=-x^2+3x+4$ (ici $a=-1<0$) : le trinôme est du signe de $a$ (négatif, zones roses) à l'extérieur des racines $-1$ et $4$, et du signe contraire (positif, zone verte) entre les racines.}
\end{popfigure}

\begin{exemplebox}[Tableau de signes complet de $-x^2+3x+4$]
Le trinôme $P(x)=-x^2+3x+4$ vérifie $a-b+c=0$ (car $-1-3+4=0$), donc $x_1=-1$ est racine évidente, et $x_2=\dfrac{c}{a\,x_1}=\dfrac{4}{(-1)(-1)}=4$. Comme $a=-1<0$, le trinôme est négatif à l'extérieur de $[-1\,;\,4]$ et positif entre les racines :
\begin{center}
\begin{tabular}{|c|ccccccc|}\hline
$x$ & $-\infty$ & & $-1$ & & $4$ & & $+\infty$ \\ \hline
$-x^2+3x+4$ & & $-$ & $0$ & $+$ & $0$ & $-$ & \\ \hline
\end{tabular}
\end{center}
\end{exemplebox}

\courssub{Résolution des inéquations du second degré}

\begin{methode}[Résoudre une inéquation du second degré en 4 étapes]
Pour résoudre $ax^2+bx+c<0$ (ou $\leq$, $>$, $\geq$) :
\begin{enumerate}
\item \textbf{Se ramener à une comparaison à $0$} : regrouper tous les termes dans un seul membre.
\item \textbf{Calculer le discriminant} $\Delta=b^2-4ac$ du trinôme obtenu.
\item \textbf{Déterminer les racines} éventuelles (par $\Delta$ ou par racine évidente).
\item \textbf{Dresser le tableau de signes} et lire l'ensemble des solutions : les intervalles où le signe convient, en incluant les bornes si l'inégalité est large ($\leq$ ou $\geq$) et en les excluant si elle est stricte ($<$ ou $>$).
\end{enumerate}
\end{methode}

\begin{exoresolu}[Résoudre $-2x^2+7x-3\geq 0$]
Résoudre dans $\mathbb{R}$ l'inéquation $-2x^2+7x-3\geq 0$.
\tcblower
\textbf{Étape 1.} L'inéquation est déjà sous la forme $P(x)\geq 0$ avec $P(x)=-2x^2+7x-3$.
\textbf{Étape 2.} $\Delta=7^2-4\times(-2)\times(-3)=49-24=25>0$.
\textbf{Étape 3.} $\sqrt{\Delta}=5$, d'où $x_1=\dfrac{-7+5}{-4}=\dfrac{1}{2}$ et $x_2=\dfrac{-7-5}{-4}=3$.
\textbf{Étape 4.} Comme $a=-2<0$, le trinôme est négatif à l'extérieur des racines et positif entre elles :
\begin{center}
\begin{tabular}{|c|ccccccc|}\hline
$x$ & $-\infty$ & & $\frac{1}{2}$ & & $3$ & & $+\infty$ \\ \hline
$-2x^2+7x-3$ & & $-$ & $0$ & $+$ & $0$ & $-$ & \\ \hline
\end{tabular}
\end{center}
On cherche $P(x)\geq 0$ : on garde l'intervalle entre les racines, bornes incluses car l'inégalité est large.
\[ S=\left[\frac{1}{2}\,;\,3\right]. \]
\end{exoresolu}

\begin{attention}[$\Delta<0$ ne signifie pas « pas de solution » !]
C'est le piège classique du Probatoire. Si $\Delta<0$, le trinôme \textbf{garde un signe constant} : celui de $a$. Donc :
\begin{itemize}
\item pour $2x^2+x+3>0$ : $\Delta=-23<0$ et $a=2>0$, le trinôme est toujours positif, donc $S=\mathbb{R}$ ;
\item pour $2x^2+x+3\leq 0$ : le trinôme n'est jamais négatif ni nul, donc $S=\emptyset$.
\end{itemize}
Conclure automatiquement « $S=\emptyset$ » dès que $\Delta<0$ est une erreur : tout dépend du sens de l'inégalité et du signe de $a$.
\end{attention}

\courssub{Inéquations se ramenant au second degré}

Beaucoup d'inéquations ne se présentent pas directement sous la forme $ax^2+bx+c\geq 0$, mais s'y ramènent.

\textbf{Inéquations produits.} Pour résoudre $(x-1)(x^2-3x+2)\leq 0$, on étudie le signe de chaque facteur (le trinôme se factorise : $x^2-3x+2=(x-1)(x-2)$) puis on dresse un tableau de signes global.

\textbf{Inéquations quotients avec trinômes.} Pour résoudre $\dfrac{x^2-4}{x^2-2x-3}\geq 0$, on factorise numérateur et dénominateur, on repère les \cle{valeurs interdites} (racines du dénominateur, exclues même si l'inégalité est large !) et on conclut par un tableau de signes.

\begin{exoresolu}[Inéquation quotient avec trinômes]
Résoudre dans $\mathbb{R}$ : $\dfrac{x^2-4}{x^2-2x-3}\geq 0$.
\tcblower
\textbf{Factorisations.} Numérateur : $x^2-4=(x-2)(x+2)$, racines $-2$ et $2$. Dénominateur : $x^2-2x-3$ ; racine évidente $-1$ (car $1+2-3=0$), l'autre racine est $\dfrac{c}{a\,x_1}=\dfrac{-3}{-1}=3$, donc $x^2-2x-3=(x+1)(x-3)$.
\textbf{Valeurs interdites :} $x=-1$ et $x=3$ (elles annulent le dénominateur).
\textbf{Tableau de signes :}
\begin{center}
\begin{tabular}{|c|ccccccccccc|}\hline
$x$ & $-\infty$ & & $-2$ & & $-1$ & & $2$ & & $3$ & & $+\infty$ \\ \hline
$x^2-4$ & & $+$ & $0$ & $-$ & & $-$ & $0$ & $+$ & & $+$ & \\ \hline
$x^2-2x-3$ & & $+$ & & $+$ & $0$ & $-$ & & $-$ & $0$ & $+$ & \\ \hline
Quotient & & $+$ & $0$ & $-$ & & $+$ & $0$ & $-$ & & $+$ & \\ \hline
\end{tabular}
\end{center}
Le quotient est positif ou nul sur $]-\infty\,;\,-2]$, sur $]-1\,;\,2]$ et sur $]3\,;\,+\infty[$. Les valeurs $-1$ et $3$ sont exclues (double barre), $-2$ et $2$ sont incluses (inégalité large) :
\[ S=]-\infty\,;\,-2]\;\cup\;]-1\,;\,2]\;\cup\;]3\,;\,+\infty[. \]
\end{exoresolu}

\begin{attention}[Inégalité large et dénominateur]
Même pour une inéquation du type $\geq 0$ ou $\leq 0$, les valeurs qui annulent le \textbf{dénominateur} sont toujours exclues de l'ensemble solution : on ne divise jamais par zéro. Dans le tableau, pense à la double barre ; dans la rédaction, utilise des crochets ouverts en ces bornes.
\end{attention}

\begin{exercice}[Pour t'entraîner]
\begin{enumerate}
\item Factoriser, si possible, les trinômes : $P(x)=x^2-6x+9$ ; $Q(x)=3x^2-7x+2$ ; $R(x)=x^2+x+1$.
\item Résoudre dans $\mathbb{R}$ : $x^2-5x+6<0$ ; $-3x^2+2x+8\leq 0$ ; $x^2+4x+5>0$.
\item Résoudre dans $\mathbb{R}$ : $\dfrac{2x^2-5x+2}{x-1}\leq 0$.
\end{enumerate}
\end{exercice}

\begin{corrige}[Exercices d'entraînement]
\begin{enumerate}
\item $P$ : $\Delta=36-36=0$, racine double $x_0=3$, donc $P(x)=(x-3)^2$. $Q$ : $\Delta=49-24=25$, $x_1=\dfrac{7-5}{6}=\dfrac{1}{3}$, $x_2=2$, donc $Q(x)=3\left(x-\dfrac{1}{3}\right)(x-2)=(3x-1)(x-2)$. $R$ : $\Delta=1-4=-3<0$, pas de factorisation dans $\mathbb{R}$.
\item $x^2-5x+6=(x-2)(x-3)$, $a>0$ : négatif entre les racines, $S=]2\,;\,3[$. Pour $-3x^2+2x+8$ : $\Delta=4+96=100$, $x_1=-\dfrac{4}{3}$, $x_2=2$, $a<0$ donc négatif à l'extérieur : $S=\left]-\infty\,;\,-\dfrac{4}{3}\right]\cup[2\,;\,+\infty[$. Pour $x^2+4x+5$ : $\Delta=16-20=-4<0$ et $a>0$, toujours positif : $S=\mathbb{R}$.
\item $2x^2-5x+2=(2x-1)(x-2)$, racines $\dfrac{1}{2}$ et $2$ ; valeur interdite $x=1$. Le tableau de signes donne un quotient négatif ou nul sur $\left[\dfrac{1}{2}\,;\,1\right[$ et sur $[2\,;\,+\infty[$ : $S=\left[\dfrac{1}{2}\,;\,1\right[\cup[2\,;\,+\infty[$.
\end{enumerate}
\end{corrige}

\begin{aretenir}[L'essentiel de la section]
\begin{itemize}
\item $\Delta>0$ : $P(x)=a(x-x_1)(x-x_2)$ ; $\Delta=0$ : $P(x)=a(x-x_0)^2$ ; $\Delta<0$ : pas de factorisation dans $\mathbb{R}$.
\item Somme et produit des racines : $S=-\dfrac{b}{a}$, $P=\dfrac{c}{a}$ — idéal avec une racine évidente.
\item Signe du trinôme : du signe de $a$ \textbf{à l'extérieur} des racines, du signe de $-a$ \textbf{entre} les racines ; si $\Delta<0$, signe de $a$ partout.
\item Résoudre une inéquation du second degré : se ramener à $0$, calculer $\Delta$, trouver les racines, lire le tableau de signes — et ne jamais oublier les valeurs interdites dans un quotient.
\end{itemize}
\end{aretenir}

\coursec{Systèmes d'équations se ramenant au second degré dans $\mathbb{R}^2$}

\noindent
Dans les sections précédentes, nous avons maîtrisé la factorisation et l'étude du signe des trinômes du second degré. Ces outils deviennent maintenant indispensables pour résoudre des problèmes où \emph{deux quantités inconnues} sont liées par \emph{deux relations} simultanées. Nous quittons ainsi l'équation à une inconnue pour entrer dans l'univers des \textbf{systèmes d'équations à deux inconnues}, dont la résolution se ramène, in fine, à une équation du second degré. C'est une compétence clé pour le Probatoire : modéliser une situation concrète (dimensions d'un champ, partage d'une somme, trajectoire d'un projectile) par un système, le résoudre rigoureusement et interpréter les résultats.

\courssub{Rappel : vérifier une solution et les systèmes linéaires $2 \times 2$}

\noindent
Avant d'aborder les systèmes du second degré, rappelons les bases. Un \textbf{système de deux équations à deux inconnues} s'écrit généralement :
\[
\begin{cases}
E_1(x,y) = 0 \\
E_2(x,y) = 0
\end{cases}
\]
Un \textbf{couple $(x_0 ; y_0) \in \mathbb{R}^2$} est une \textbf{solution} de ce système si et seulement si, en remplaçant $x$ par $x_0$ et $y$ par $y_0$, les \emph{deux} équations sont vérifiées simultanément.

\begin{methode}[Vérifier qu'un couple est solution]
\begin{enumerate}
    \item Remplacer $x$ et $y$ par les valeurs proposées dans la \textbf{première} équation. Vérifier l'égalité.
    \item Remplacer $x$ et $y$ par les mêmes valeurs dans la \textbf{deuxième} équation. Vérifier l'égalité.
    \item Conclure : « Le couple $(\dots ; \dots)$ est solution du système » ou « n'est pas solution ».
\end{enumerate}
\end{methode}

\begin{exemplebox}[Vérification]
Vérifier si $(2 ; 3)$ et $(1 ; -2)$ sont solutions du système $\begin{cases} x + y = 5 \\ x^2 + y^2 = 13 \end{cases}$.
\begin{itemize}
    \item Pour $(2 ; 3)$ : $2+3=5$ (ok) et $2^2+3^2=4+9=13$ (ok). \textbf{C'est une solution.}
    \item Pour $(1 ; -2)$ : $1+(-2)=-1 \neq 5$. \textbf{Ce n'est pas une solution.} (Inutile de tester la seconde).
\end{itemize}
\end{exemplebox}

\noindent
\emph{Rappel (Seconde) :} Un \textbf{système linéaire} $\begin{cases} ax+by=c \\ dx+ey=f \end{cases}$ se résout par \textbf{substitution} (exprimer $y$ en fonction de $x$ dans l'une, substituer dans l'autre) ou \textbf{combinaison linéaire} (multiplier les équations pour éliminer une inconnue). Géométriquement, c'est l'intersection de deux droites : une solution unique (droites sécantes), aucune (droites parallèles distinctes) ou une infinité (droites confondues).

\courssub{Les systèmes symétriques : la somme $S$ et le produit $P$}

\noindent
Considérons un système de la forme :
\[
(S) : \begin{cases} x + y = S \\ xy = P \end{cases}
\]
où $S$ et $P$ sont des nombres réels donnés. Ce système est dit \textbf{symétrique} car les équations ne changent pas si l'on échange $x$ et $y$. Si $(x_0 ; y_0)$ est solution, alors $(y_0 ; x_0)$ l'est aussi.

\begin{propriete}[Théorème fondamental : Somme et Produit]
\label{th:sommeproduit}
Deux nombres réels $x$ et $y$ sont de somme $S$ et de produit $P$ \textbf{si et seulement si} ils sont les deux solutions (éventuellement confondues) de l'équation du second degré \textbf{inconnue $X$} :
\[
X^2 - S X + P = 0
\]
\end{propriete}

\begin{experience}[Démonstration du théorème (exigible au Probatoire)]
\textbf{1. Sens direct ($\Rightarrow$).} Supposons $x+y=S$ et $xy=P$.
Calculons $(X-x)(X-y)$ :
\[
(X-x)(X-y) = X^2 - (x+y)X + xy = X^2 - S X + P.
\]
Ainsi, $x$ et $y$ annulent le trinôme $X^2 - SX + P$. Ils sont donc solutions de l'équation $X^2 - SX + P = 0$.

\textbf{2. Sens réciproque ($\Leftarrow$).} Soient $x$ et $y$ les deux solutions de $X^2 - SX + P = 0$ (dans $\mathbb{R}$).
D'après les relations coefficients-racines (Vieta) : la somme des solutions est $-\frac{-S}{1} = S$ et leur produit est $\frac{P}{1} = P$.
Donc $x+y=S$ et $xy=P$. \hfill $\square$
\end{experience}

\begin{aretenir}[Condition d'existence de solutions réelles]
L'équation $X^2 - SX + P = 0$ admet des solutions réelles \textbf{si et seulement si} son discriminant est positif ou nul :
\[
\Delta = S^2 - 4P \geq 0
\]
Si $\Delta < 0$, le système $(S)$ n'a \textbf{aucune solution dans $\mathbb{R}^2$}.
\cle{Condition d'existence} \cle{Discriminant} \cle{Système symétrique}
\end{aretenir}

\begin{methode}[Résoudre un système symétrique $\begin{cases} x+y=S \\ xy=P \end{cases}$]
\begin{enumerate}
    \item Écrire l'équation associée : $X^2 - SX + P = 0$.
    \item Calculer $\Delta = S^2 - 4P$.
    \item \textbf{Si $\Delta < 0$ :} Aucune solution réelle. Stop.
    \item \textbf{Si $\Delta \geq 0$ :} Calculer les deux racines $X_1 = \frac{S - \sqrt{\Delta}}{2}$ et $X_2 = \frac{S + \sqrt{\Delta}}{2}$.
    \item Écrire les couples solutions : $(X_1 ; X_2)$ et $(X_2 ; X_1)$.
    \item \emph{(Si $\Delta = 0$, les deux couples se confondent en un unique couple $(\frac{S}{2} ; \frac{S}{2})$).}
\end{enumerate}
\cle{Méthode} \cle{Somme} \cle{Produit}
\end{methode}

\begin{exoresolu}[Système symétrique classique]
\textbf{Énoncé :} Résoudre dans $\mathbb{R}^2$ le système $\begin{cases} x + y = 7 \\ xy = 12 \end{cases}$.
\tcblower
\textbf{Solution détaillée :}
\begin{enumerate}
    \item $S = 7$, $P = 12$. Équation : $X^2 - 7X + 12 = 0$.
    \item $\Delta = (-7)^2 - 4 \times 12 = 49 - 48 = 1 \geq 0$. Solutions réelles existantes.
    \item $\sqrt{\Delta} = 1$.
    \item $X_1 = \frac{7 - 1}{2} = 3$ ; $X_2 = \frac{7 + 1}{2} = 4$.
    \item Les couples solutions sont $\boxed{(3 ; 4)}$ et $\boxed{(4 ; 3)}$.
\end{enumerate}
\textbf{Vérification :} $3+4=7$, $3\times4=12$ ; $4+3=7$, $4\times3=12$. C'est bon.
\end{exoresolu}

\begin{attention}[Pièges fréquents sur les systèmes symétriques]
\begin{itemize}
    \item \textbf{Oublier le second couple :} Si $x \neq y$, il y a \emph{deux} couples solutions $(x_1;y_1)$ et $(y_1;x_1)$. Ne donnez pas seulement « $x=3, y=4$ ».
    \item \textbf{Ne pas vérifier $\Delta \geq 0$ :} Si $\Delta < 0$, écrire « Pas de solution » sans calculer de racines complexes (hors programme $\mathbb{R}$).
    \item \textbf{Confondre $S$ et $P$ :} L'équation est $X^2 - \mathbf{S}X + \mathbf{P} = 0$. Le signe moins devant $S$ est crucial.
\end{itemize}
\cle{Pièges} \cle{Recommandations}
\end{attention}

\courssub{Les systèmes mixtes : une droite et une courbe du second degré}

\noindent
Un \textbf{système mixte} comporte une équation du premier degré (droite) et une équation du second degré (parabole, hyperbole, cercle...). La méthode universelle est la \textbf{substitution}.

\begin{methode}[Résoudre un système mixte par substitution]
\begin{enumerate}
    \item Dans l'équation \textbf{linéaire}, exprimer une inconnue en fonction de l'autre (généralement $y = ax+b$ ou $x = cy+d$). Choisir l'expression la plus simple (éviter les fractions si possible).
    \item Substituer cette expression dans l'équation du \textbf{second degré}. On obtient une équation à une inconnue (souvent du second degré).
    \item Résoudre cette équation dans $\mathbb{R}$ (discriminant, factorisation, etc.).
    \item Pour chaque valeur trouvée, retrouver la valeur de l'autre inconnue grâce à l'expression de l'étape 1.
    \item \textbf{Vérifier} chaque couple dans les \textbf{deux} équations initiales (surtout si on a multiplié par une expression pouvant s'annuler, ou élevé au carré).
\end{enumerate}
\cle{Substitution} \cle{Système mixte}
\end{methode}

\begin{exoresolu}[Système mixte : Droite et Parabole]
\textbf{Énoncé :} Résoudre $\begin{cases} y = 2x + 1 \\ y = x^2 \end{cases}$.
\tcblower
\textbf{Solution détaillée :}
\begin{enumerate}
    \item La première équation donne directement $y = 2x+1$.
    \item Substitution dans la seconde : $x^2 = 2x + 1 \iff x^2 - 2x - 1 = 0$.
    \item $\Delta = (-2)^2 - 4(1)(-1) = 4 + 4 = 8 > 0$.
    \item $x_1 = \frac{2 - \sqrt{8}}{2} = 1 - \sqrt{2}$ ; $x_2 = \frac{2 + \sqrt{8}}{2} = 1 + \sqrt{2}$.
    \item $y_1 = 2(1-\sqrt{2})+1 = 3 - 2\sqrt{2}$ ; $y_2 = 2(1+\sqrt{2})+1 = 3 + 2\sqrt{2}$.
    \item Solutions : $\boxed{(1-\sqrt{2} ; 3-2\sqrt{2})}$ et $\boxed{(1+\sqrt{2} ; 3+2\sqrt{2})}$.
\end{enumerate}
\end{exoresolu}

\begin{popfigure}
\begin{tikzpicture}
\begin{axis}[
    width=\linewidth,
    height=8cm,
    axis lines=middle,
    xlabel=$x$, ylabel=$y$,
    xmin=-2.5, xmax=3.5,
    ymin=-1, ymax=6,
    xtick={-2,-1,1,2,3},
    ytick={1,2,3,4,5},
    grid=major, grid style={popline},
    clip=false,
    legend style={at={(0.02,0.98)}, anchor=north west, fill=popcream, draw=popmuted}
]
\addplot[popcurve, domain=-2.2:3.2, samples=200, thick] {x^2};
\addlegendentry{$y = x^2$ (Parabole)}
\addplot[poporange, domain=-2.5:3.5, samples=2, thick] {2*x+1};
\addlegendentry{$y = 2x+1$ (Droite)}
% Points d'intersection : x = 1 +/- sqrt(2) approx -0.4142 et 2.4142
% y = x^2 approx 0.1716 et 5.8284
\addplot[only marks, mark=*, mark size=2.5pt, popblue] coordinates {(-0.4142, 0.1716) (2.4142, 5.8284)};
\node[popblue, above left] at (axis cs:-0.4142, 0.1716) {$A$};
\node[popblue, above right] at (axis cs:2.4142, 5.8284) {$B$};
\end{axis}
\end{tikzpicture}
\legende{Intersection de la droite $y=2x+1$ et de la parabole $y=x^2$. Les abscisses des points $A$ et $B$ sont les solutions de $x^2-2x-1=0$.}
\end{popfigure}

\courssub{Interprétation géométrique : intersection droite / courbe}

\noindent
Résoudre un système mixte, c'est trouver les points d'intersection d'une \textbf{droite} et d'une \textbf{courbe du second degré} (parabole, hyperbole, cercle, ellipse).
\begin{itemize}
    \item $\Delta > 0$ : \textbf{Deux points d'intersection} distincts (la droite est sécante).
    \item $\Delta = 0$ : \textbf{Un point d'intersection} (la droite est \textbf{tangente} à la courbe).
    \item $\Delta < 0$ : \textbf{Aucun point d'intersection} (la droite ne rencontre pas la courbe).
\end{itemize}

\begin{exemplebox}[Cas tangent]
Système $\begin{cases} y = 2x - 1 \\ y = x^2 \end{cases}$.
Substitution : $x^2 = 2x - 1 \iff x^2 - 2x + 1 = 0 \iff (x-1)^2 = 0$.
$\Delta = 0$, solution unique $x=1$, $y=1$. La droite $y=2x-1$ est tangente à la parabole $y=x^2$ au point $(1;1)$.
\cle{Tangente} \cle{Intersection}
\end{exemplebox}

\courssub{Problèmes concrets : modélisation et résolution}

\noindent
C'est ici que les mathématiques prennent tout leur sens. Au Probatoire, l'énoncé ne donne pas le système ; \textbf{vous devez le construire}.

\begin{methode}[Résoudre un problème par un système du second degré]
\begin{enumerate}
    \item \textbf{Variables :} Choisir deux inconnues pertinentes (ex: $x$ = longueur, $y$ = largeur). Préciser leur ensemble de définition (souvent $>0$).
    \item \textbf{Traduction :} Écrire les deux relations données par l'énoncé sous forme d'équations (système).
    \item \textbf{Résolution :} Identifier le type (symétrique ou mixte) et appliquer la méthode adaptée.
    \item \textbf{Sélection :} Éliminer les solutions inadaptées au contexte (longueur négative, nombre non entier de personnes, etc.).
    \item \textbf{Réponse :} Formuler une phrase complète répondant à la question.
\end{enumerate}
\cle{Modélisation} \cle{Problème concret}
\end{methode}

\begin{exoresolu}[Le champ rectangulaire (Classique Probatoire)]
\textbf{Énoncé :} Un agriculteur à Bafia souhaite clôturer un champ rectangulaire. Il dispose de \SI{60}{m} de grillage pour le périmètre et souhaite une surface de \SI{200}{m^2}. Quelles sont les dimensions du champ ?
\tcblower
\textbf{Solution détaillée :}
\begin{enumerate}
    \item \textbf{Variables :} Soit $x$ la longueur (en m) et $y$ la largeur (en m). $x > 0, y > 0$.
    \item \textbf{Système :}
    \[
    \begin{cases}
    \text{Périmètre : } 2(x+y) = 60 \iff x+y = 30 \\
    \text{Aire : } xy = 200
    \end{cases}
    \]
    C'est un système \textbf{symétrique} avec $S=30$, $P=200$.
    \item \textbf{Résolution :} Équation $X^2 - 30X + 200 = 0$.
    $\Delta = 900 - 800 = 100$. $\sqrt{\Delta} = 10$.
    $X_1 = \frac{30-10}{2} = 10$ ; $X_2 = \frac{30+10}{2} = 20$.
    Couples : $(10 ; 20)$ et $(20 ; 10)$.
    \item \textbf{Sélection :} Les deux couples donnent des dimensions positives. Par convention, la longueur $\geq$ largeur.
    \item \textbf{Réponse :} Le champ mesure \textbf{\SI{20}{m} de long et \SI{10}{m} de large}.
\end{enumerate}
\end{exoresolu}

\begin{popfigure}
\begin{tikzpicture}[scale=0.6, every node/.style={font=\small}]
% Rectangle
\draw[popdark, thick] (0,0) rectangle (10,5);
% Dimensions
\draw[|<->|, popblue, thick] (0,-1) -- (10,-1) node[midway, below=2pt] {$x$ (longueur)};
\draw[|<->|, poporange, thick] (11,0) -- (11,5) node[midway, right=2pt] {$y$ (largeur)};
% Annotations
\node[popblue] at (5, -2.5) {Périmètre $= 2(x+y) = 60$};
\node[poporange] at (5, -3.5) {Aire $= xy = 200$};
\node[align=center, popmuted] at (5, 2.5) {Système symétrique :\\ $x+y=30$ ; $xy=200$};
\end{tikzpicture}
\legende{Modélisation géométrique du problème du champ : le système symétrique $x+y=S$, $xy=P$.}
\end{popfigure}

\begin{exoresolu}[Partage d'une somme (Contexte MTN Mobile Money)]
\textbf{Énoncé :} Deux commerçants, Amadou et Bintou, mettent en commun leur argent pour acheter un stock de cartes de recharge MTN. Le produit de leurs apports (en milliers de FCFA) est $120$ et la somme de leurs apports est $23$. Combien a apporté chacun ?
\tcblower
\textbf{Solution :}
\begin{enumerate}
    \item $x$ = apport d'Amadou (en kFCFA), $y$ = apport de Bintou (en kFCFA). $x>0, y>0$.
    \item Système : $\begin{cases} x+y = 23 \\ xy = 120 \end{cases}$ ($S=23, P=120$).
    \item $X^2 - 23X + 120 = 0$. $\Delta = 529 - 480 = 49$. $\sqrt{\Delta}=7$.
    \item $X_1 = \frac{23-7}{2} = 8$ ; $X_2 = \frac{23+7}{2} = 15$.
    \item Apports : \textbf{8 000 FCFA et 15 000 FCFA}. (On ne sait pas qui a donné quoi sans info supplémentaire).
\end{enumerate}
\end{exoresolu}

\begin{exoresolu}[Système mixte : Problème de mouvement / Trajectoire]
\textbf{Énoncé :} Un drone de surveillance (CAMTEL) suit une trajectoire parabolique $y = -x^2 + 4x + 5$ (hauteur $y$ en mètres, distance horizontale $x$ en mètres). Un câble tendu relie le point $(0;0)$ au point $(5;10)$. Le câble suit la droite $y=2x$. Le drone coupe-t-il le câble ? Si oui, où ?
\tcblower
\textbf{Solution :}
\begin{enumerate}
    \item Système mixte : $\begin{cases} y = 2x \\ y = -x^2 + 4x + 5 \end{cases}$.
    \item Substitution : $2x = -x^2 + 4x + 5 \iff x^2 - 2x - 5 = 0$.
    \item $\Delta = 4 + 20 = 24 > 0$. Deux intersections.
    \item $x_1 = 1 - \sqrt{6} \approx -1{,}45$ (rejeté, $x<0$ hors trajectoire utile) ; $x_2 = 1 + \sqrt{6} \approx 3{,}45$.
    \item $y_2 = 2(1+\sqrt{6}) = 2 + 2\sqrt{6} \approx 6{,}90$.
    \item \textbf{Réponse :} Oui, le drone coupe le câble au point de coordonnées approximatives \textbf{$(3{,}45 ; 6{,}90)$} (en mètres).
\end{enumerate}
\end{exoresolu}

\courssub{Inéquations se ramenant au second degré à deux inconnues (Rappel)}

\noindent
Le programme demande aussi de « Résoudre dans $\mathbb{R}^2$, des inéquations... ». Il s'agit de déterminer l'ensemble des couples $(x;y)$ vérifiant une inéquation comme $x^2 + y^2 < 9$ (disque ouvert) ou $y > x^2$ (région au-dessus de la parabole). La méthode graphique (hachurage) est privilégiée.

\begin{exemplebox}[Résolution graphique dans $\mathbb{R}^2$]
Résoudre $y \leq x^2 - 4$.
\begin{enumerate}
    \item Tracer la parabole $y = x^2 - 4$ (trait \textbf{plein} car $\leq$).
    \item Prendre un point test facile, ex: $O(0;0)$. $0 \leq -4$ ? \textbf{Faux}.
    \item Hachurer la région \textbf{ne contenant pas} $O$ (l'extérieur de la parabole, vers le haut).
\end{enumerate}
L'ensemble solution est $\{(x,y) \in \mathbb{R}^2 \mid y \leq x^2 - 4\}$.
\cle{Inéquation à deux inconnues} \cle{Résolution graphique}
\end{exemplebox}

\begin{attention}[Rappel officiel : Hors programme]
\begin{itemize}
    \item \textbf{Les systèmes d'équations \emph{avec paramètres} sont hors programme.} Vous ne rencontrerez pas de $m, k, a$ dans les coefficients du système au Probatoire. Les coefficients sont des nombres fixés.
    \item \textbf{La résolution algébrique des équations du second degré \emph{avec paramètres} est hors programme.}
    \item \textbf{Le discriminant réduit ($\Delta'$) est hors programme.} Utilisez toujours $\Delta = b^2 - 4ac$.
\end{itemize}
Concentrez votre énergie sur la maîtrise parfaite des systèmes à coefficients numériques et de la modélisation de problèmes concrets.
\cle{Hors programme} \cle{Probatoire}
\end{attention}

\begin{savaistu}[L'arpentage au Cameroun : des maths pour la terre]
Saviez-vous que les systèmes symétriques $x+y=S, xy=P$ sont au cœur de l'\textbf{arpentage} (bornage des terrains) depuis l'Antiquité ? Au Cameroun, lors de la délimitation d'une parcelle rectangulaire (certificat foncier, titre foncier), le géomètre-expert mesure souvent le \textbf{périmètre} (ou demi-périmètre $S$) et la \textbf{superficie} $P$. Retrouver les longueurs des côtés revient exactement à résoudre $X^2 - SX + P = 0$.
De même, dans les successions (partage d'héritage), si deux héritiers doivent recevoir des parts dont on connaît la somme totale et le produit (lié à une clause de partage proportionnel), on tombe sur ce même système. Ces problèmes « terre et argent » sont des classiques indémodables du Probatoire Série A ! Maîtrisez-les, ils rapportent des points faciles.
\cle{Arpentage} \cle{Foncier} \cle{Cameroun}
\end{savaistu}

\begin{exercice}[Entraînement Probatoire]
\begin{enumerate}
    \item Résoudre dans $\mathbb{R}^2$ : $\begin{cases} x+y=10 \\ xy=21 \end{cases}$.
    \item Résoudre : $\begin{cases} y = x - 2 \\ y = x^2 - 4 \end{cases}$. Donner une interprétation géométrique.
    \item Un rectangle a pour périmètre \SI{40}{cm} et pour diagonale \SI{15}{cm}. Trouver ses dimensions. (Indice : Théorème de Pythagore $\to x^2+y^2=225$).
    \item Deux nombres ont pour somme $14$ et pour somme des carrés $100$. Trouver ces nombres.
\end{enumerate}
\end{exercice}

\begin{corrige}[Exercice 1]
\begin{itemize}
\item $S=10, P=21$. $\Delta=100-84=16$. $X_1=3, X_2=7$. Solutions : $\boxed{(3;7)}$ et $\boxed{(7;3)}$.
\item Substitution : $x-2 = x^2-4 \iff x^2-x-2=0 \iff (x-2)(x+1)=0$. $x=2$ ou $x=-1$.
Couples : $(2;0)$ et $(-1;-3)$. Géométrie : Intersection de la droite $y=x-2$ et de la parabole $y=x^2-4$. Deux points d'intersection (droite sécante).
\item Soit $x, y$ les côtés. $\begin{cases} 2(x+y)=40 \Rightarrow x+y=20 \\ x^2+y^2=15^2=225 \end{cases}$.
On a $S=20$. On cherche $P=xy$. On sait que $(x+y)^2 = x^2+y^2+2xy \Rightarrow 400 = 225 + 2P \Rightarrow 2P=175 \Rightarrow P=87{,}5$.
Système symétrique : $X^2 - 20X + 87{,}5 = 0$. $\Delta = 400 - 350 = 50$.
$X = \frac{20 \pm \sqrt{50}}{2} = 10 \pm \frac{5\sqrt{2}}{2}$.
Dimensions : $\boxed{\left(10+\frac{5\sqrt{2}}{2} ; 10-\frac{5\sqrt{2}}{2}\right)}$ cm (environ $13{,}54$ cm et $6{,}46$ cm).
\item Soit $x, y$ les nombres. $\begin{cases} x+y=14 \\ x^2+y^2=100 \end{cases}$.
Même astuce : $196 = 100 + 2xy \Rightarrow xy = 48$.
Système symétrique $S=14, P=48$. $X^2-14X+48=0$. $\Delta=196-192=4$.
$X_1=6, X_2=8$. Les nombres sont \textbf{6 et 8}.
\end{itemize}
\end{corrige}

\coursec{Problèmes de synthèse et modélisation}

Dans les sections précédentes, tu as appris à résoudre des équations et inéquations-quotients, à écrire la forme canonique d'un trinôme, à utiliser le discriminant, à factoriser et à étudier le signe d'un polynôme du second degré, puis à résoudre des systèmes se ramenant au second degré. Tu disposes donc maintenant de toute la « boîte à outils » de cette leçon. Mais à quoi servent ces outils ? À résoudre des \cle{problèmes concrets} : calculer les dimensions d'un champ, fixer un prix de vente rentable, répartir une somme d'argent, déterminer une durée de trajet. C'est exactement ce que le programme appelle « résoudre des problèmes concrets se ramenant à la résolution d'équations, d'inéquations et de systèmes étudiés ». Cette dernière section t'apprend à passer du texte d'un problème aux mathématiques, puis des mathématiques à la réponse concrète.

\courssub{La méthodologie générale de modélisation}

\begin{definition}[Modéliser un problème]
\cle{Modéliser} un problème concret, c'est traduire sa situation en langage mathématique : on choisit une ou plusieurs \cle{inconnues}, on traduit les contraintes de l'énoncé en \cle{équations} ou \cle{inéquations}, on les résout, puis on \emph{interprète} les résultats dans le contexte de départ.
\end{definition}

L'erreur la plus fréquente des élèves n'est pas de mal résoudre l'équation : c'est de mal la construire, ou d'oublier de vérifier que la solution trouvée a un sens concret. Voici la démarche complète, à suivre systématiquement.

\begin{methode}[Les 5 étapes de la résolution d'un problème par mise en équation]
\begin{enumerate}
\item \textbf{Comprendre et choisir l'inconnue.} Lis attentivement l'énoncé, repère ce qu'on cherche, et pose $x = $ « la quantité cherchée ». Précise immédiatement l'\cle{ensemble des valeurs possibles} de $x$ (par exemple $x > 0$ pour une longueur, $x \in \mathbb{N}$ pour un nombre d'objets).
\item \textbf{Mettre en équation.} Traduis chaque information de l'énoncé par une relation mathématique portant sur $x$. C'est l'étape la plus délicate : procède phrase par phrase.
\item \textbf{Résoudre.} Résous l'équation, l'inéquation ou le système obtenu avec les techniques du chapitre (discriminant, tableau de signes, substitution\ldots).
\item \textbf{Vérifier et critiquer.} Confronte chaque solution mathématique aux contraintes du contexte : une longueur négative, un prix négatif, un nombre d'élèves non entier doivent être \cle{rejetés}.
\item \textbf{Conclure en français.} Rédige une phrase de conclusion qui répond exactement à la question posée, avec les unités.
\end{enumerate}
\end{methode}

\begin{attention}[Solution mathématiquement correcte mais impossible dans le contexte]
Une équation du second degré donne souvent \emph{deux} solutions. Toutes ne sont pas acceptables ! Si $x$ désigne une longueur, un prix ou une quantité, toute solution négative ou nulle doit être rejetée. De même, si $x$ est un nombre de personnes ou d'articles, seules les valeurs entières conviennent. \textbf{Au Probatoire, ne pas rejeter une solution hors contexte coûte des points}, même si la résolution est parfaite. Prends l'habitude d'écrire explicitement : « cette solution est à rejeter car une longueur est positive ».
\end{attention}

\courssub{Problèmes se ramenant à une équation du second degré}

Les problèmes d'\cle{aire} et de \cle{dimensions} conduisent naturellement au second degré, car une aire fait intervenir un produit de deux longueurs.

\begin{exoresolu}[Le jardin de Monsieur Fotso]
Monsieur Fotso possède à Bafoussam un jardin rectangulaire de \SI{20}{m} de long sur \SI{12}{m} de large. Il veut aménager une allée de largeur constante $x$ (en mètres) tout autour du jardin, à l'intérieur, de sorte que la partie cultivable restante ait une aire de \SI{162}{m^2}. Déterminer la largeur $x$ de l'allée.
\tcblower
\textbf{Étape 1 — Inconnue.} Soit $x$ la largeur de l'allée, en mètres. Comme l'allée est aménagée des deux côtés de chaque dimension, on doit avoir $0 < x < 6$ (sinon la largeur $12 - 2x$ deviendrait négative).

\textbf{Étape 2 — Mise en équation.} La partie cultivable est un rectangle de longueur $20 - 2x$ et de largeur $12 - 2x$. Son aire vaut $162$, donc :
\[ (20 - 2x)(12 - 2x) = 162. \]

\textbf{Étape 3 — Résolution.} Développons :
\begin{align*}
240 - 40x - 24x + 4x^2 &= 162\\
4x^2 - 64x + 240 - 162 &= 0\\
4x^2 - 64x + 78 &= 0\\
2x^2 - 32x + 39 &= 0 \quad (\text{en divisant par } 2).
\end{align*}
Calculons le discriminant : $\Delta = (-32)^2 - 4 \times 2 \times 39 = 1024 - 312 = 712$. Or $712 = 4 \times 178$, donc $\sqrt{\Delta} = 2\sqrt{178}$. Les solutions sont :
\[ x_1 = \frac{32 - 2\sqrt{178}}{4} = \frac{16 - \sqrt{178}}{2} \approx 1{,}33 \qquad x_2 = \frac{16 + \sqrt{178}}{2} \approx 14{,}67. \]

\textbf{Étape 4 — Critique.} $x_2 \approx 14{,}67 > 6$ : cette solution est à rejeter, car l'allée ne peut pas être plus large que la moitié de la largeur du jardin. Seule $x_1 \approx 1{,}33$ convient.

\textbf{Étape 5 — Conclusion.} L'allée doit avoir une largeur d'environ \SI{1,33}{m}. Vérification : $(20 - 2 \times 1{,}33)(12 - 2 \times 1{,}33) \approx 17{,}34 \times 9{,}34 \approx 162$.
\end{exoresolu}

\begin{popfigure}
\begin{center}
\begin{tikzpicture}[scale=0.32]
% Allée (bande) : on dessine le rectangle extérieur en couleur allée puis l'intérieur en vert
\fill[popgoldL] (0,0) rectangle (20,12);
\fill[popgreen] (1.33,1.33) rectangle (18.67,10.67);
\draw[very thick,popdark] (0,0) rectangle (20,12);
\draw[very thick,popdark] (1.33,1.33) rectangle (18.67,10.67);
% Cotes
\draw[<->,>=stealth,popink] (0,-0.9) -- (20,-0.9) node[midway,below]{$20$ m};
\draw[<->,>=stealth,popink] (-0.9,0) -- (-0.9,12) node[midway,left]{$12$ m};
\draw[<->,>=stealth,poporange] (20.6,10.67) -- (20.6,12) node[midway,right]{$x$};
\draw[<->,>=stealth,poporange] (18.67,12.6) -- (20,12.6) node[midway,above]{$x$};
\node[popcream] at (10,6.4) {\textbf{Zone cultivable}};
\node[popcream] at (10,5.2) {$(20-2x)(12-2x)=162$};
\node[popdark] at (10,11.4) {\small allée};
\end{tikzpicture}
\end{center}
\legende{Le jardin de Monsieur Fotso : l'allée de largeur $x$ entoure la zone cultivable.}
\end{popfigure}

\courssub{Problèmes se ramenant à une inéquation : seuil de rentabilité}

Quand un énoncé contient les expressions « au moins », « au plus », « ne pas dépasser », « être rentable », c'est une \cle{inéquation} qu'il faut écrire, et non une équation. La résolution repose alors sur le \cle{tableau de signes} d'un trinôme ou d'un quotient.

\begin{exemplebox}[La recette de Mama Ngo Bell]
Mama Ngo Bell vend des beignets au marché de Douala. Si elle fixe le prix d'un sachet à $x$ francs CFA, elle estime en vendre $300 - 2x$ sachets par jour (plus le prix est élevé, moins elle en vend). Sa recette journalière est donc, en francs :
\[ R(x) = x(300 - 2x) = -2x^2 + 300x. \]
Elle souhaite une recette \textbf{d'au moins} \num{10000} FCFA par jour. Pour quels prix $x$ ce seuil est-il atteint ?

\textbf{Étape 1.} $x$ est un prix : $x > 0$, et il faut vendre un nombre positif de sachets : $300 - 2x > 0$, soit $x < 150$.

\textbf{Étape 2.} L'inéquation à résoudre est $-2x^2 + 300x \geq \num{10000}$, c'est-à-dire $-2x^2 + 300x - \num{10000} \geq 0$, ou encore, en divisant par $-2$ (attention au changement de sens !) : $x^2 - 150x + \num{5000} \leq 0$.

\textbf{Étape 3.} Discriminant : $\Delta = 150^2 - 4 \times \num{5000} = \num{22500} - \num{20000} = \num{2500} = 50^2$. Racines :
\[ x_1 = \frac{150 - 50}{2} = 50, \qquad x_2 = \frac{150 + 50}{2} = 100. \]
Le trinôme $x^2 - 150x + \num{5000}$ (coefficient de $x^2$ positif) est négatif \emph{entre} ses racines. Donc $x^2 - 150x + \num{5000} \leq 0 \iff x \in [50\,;\,100]$.

\textbf{Étape 4.} L'intervalle $[50\,;\,100]$ est bien inclus dans les contraintes $]0\,;\,150[$ : tout est acceptable.

\textbf{Étape 5.} Mama Ngo Bell doit fixer son prix entre \num{50} et \num{100} FCFA le sachet pour réaliser au moins \num{10000} FCFA de recette.
\end{exemplebox}

\begin{popfigure}
\begin{center}
\begin{tikzpicture}
\begin{axis}[
  width=\linewidth, height=6.2cm,
  axis lines=middle,
  xlabel={Prix $x$ (FCFA)}, ylabel={Recette $R(x)$ (FCFA)},
  xmin=0, xmax=160, ymin=0, ymax=12500,
  xtick={0,50,75,100,150},
  ytick={2500,5000,7500,10000,11250},
  tick label style={font=\small},
  label style={font=\small},
  clip=false
]
% Zone colorée où R(x) >= 10000
\addplot[draw=none, fill=popgreenL, domain=50:100, samples=60] {-2*x^2+300*x} \closedcycle;
\addplot[draw=none, fill=white, domain=50:100, samples=2] {10000} \closedcycle;
% Courbe de la recette
\addplot[popcurve, domain=0:150, samples=200] {-2*x^2+300*x};
% Seuil 10000
\addplot[dashed, poporange, thick, domain=0:160, samples=2] {10000};
% Points remarquables
\addplot[only marks, mark=*, popink, mark size=1.6pt] coordinates {(50,10000) (100,10000) (75,11250)};
\node[popink, font=\small, anchor=south] at (axis cs:75,11350) {Maximum : $R(75)=\num{11250}$};
\node[poporange, font=\small, anchor=south west] at (axis cs:105,10000) {seuil \num{10000}};
\node[popink, font=\small, anchor=north] at (axis cs:50,9600) {$50$};
\node[popink, font=\small, anchor=north] at (axis cs:100,9600) {$100$};
\end{axis}
\end{tikzpicture}
\end{center}
\legende{Recette $R(x) = -2x^2 + 300x$. La zone colorée correspond à $R(x) \geq \num{10000}$, c'est-à-dire $x \in [50\,;\,100]$.}
\end{popfigure}

\begin{savaistu}[La forme canonique donne le maximum immédiatement]
Regarde la recette de Mama Ngo Bell : $R(x) = -2x^2 + 300x$. Sa forme canonique est
\[ R(x) = -2(x - 75)^2 + \num{11250}. \]
Comme $-2(x-75)^2 \leq 0$ pour tout $x$, la recette est \emph{maximale} lorsque $x = 75$, et ce maximum vaut \num{11250} FCFA : c'est le sommet de la parabole, visible sur le graphique. Aucune équation à résoudre ! En \textbf{Terminale}, tu retrouveras cette idée avec un outil plus puissant : la \emph{dérivation}. La dérivée de $R$ est $R'(x) = -4x + 300$, qui s'annule précisément en $x = 75$ — l'abscisse du sommet. La forme canonique de Première et la dérivation de Terminale sont deux visages du même phénomène. \emph{(Ce lien est un complément culturel, hors programme strict de Première A.)}
\end{savaistu}

\courssub{Problèmes se ramenant à un système}

Les problèmes de \cle{partages}, de \cle{mélanges} ou de recherche de deux quantités liées par deux conditions se traduisent par un système à deux inconnues, qui se ramène souvent au second degré (méthode somme-produit ou substitution).

\begin{exemplebox}[Trouver deux nombres]
Trouver deux nombres réels dont la différence est $5$ et le produit $84$.

\textbf{Mise en système.} Soient $x$ et $y$ ces deux nombres, avec $x \geq y$. L'énoncé se traduit par :
\[ \begin{cases} x - y = 5 \\ xy = 84 \end{cases} \]

\textbf{Résolution par substitution.} De la première équation, $x = y + 5$. En reportant dans la seconde :
\[ (y+5)y = 84 \iff y^2 + 5y - 84 = 0. \]
Discriminant : $\Delta = 25 + 336 = 361 = 19^2$. D'où :
\[ y_1 = \frac{-5 - 19}{2} = -12, \qquad y_2 = \frac{-5 + 19}{2} = 7. \]
Si $y = -12$, alors $x = -12 + 5 = -7$. Si $y = 7$, alors $x = 7 + 5 = 12$.

\textbf{Vérification.} $-7 - (-12) = 5$ et $(-7) \times (-12) = 84$ ; $12 - 7 = 5$ et $12 \times 7 = 84$.

\textbf{Conclusion.} Il y a deux couples solutions : $(12\,;\,7)$ et $(-7\,;\,-12)$.
\end{exemplebox}

\begin{methode}[Vérifier qu'un couple est solution d'un système]
Pour vérifier qu'un couple $(x\,;\,y)$ est solution d'un système, on \textbf{remplace} $x$ et $y$ par leurs valeurs dans \emph{chacune} des équations et on contrôle que \emph{toutes} les égalités sont vraies. Un couple qui vérifie une seule équation n'est pas solution. Exemple : $(9\,;\,4)$ vérifie $x - y = 5$ mais $9 \times 4 = 36 \neq 84$ : ce n'est pas solution du système ci-dessus.
\end{methode}

\courssub{Problèmes à quotient : moyennes et partages inverses}

Certaines situations font intervenir des rapports : vitesses moyennes, prix unitaires, répartitions. Elles conduisent à des équations du type $\dfrac{ax+b}{cx+d} = k$ ou à des inéquations-quotients, résolues par tableau de signes (section 1).

\begin{exoresolu}[Le transporteur de la Sawa]
Un transporteur effectue le trajet Douala--Bafoussam de \SI{270}{km}. À l'aller, il roule à une vitesse moyenne de $v$ km/h. Au retour, alourdi par des marchandises, sa vitesse chute de \SI{15}{km/h} et le trajet dure $1$ h de plus qu'à l'aller. Déterminer $v$.
\tcblower
\textbf{Mise en équation.} Durée de l'aller : $\dfrac{270}{v}$ ; durée du retour : $\dfrac{270}{v-15}$. La contrainte s'écrit :
\[ \frac{270}{v-15} - \frac{270}{v} = 1, \qquad \text{avec } v > 15. \]
\textbf{Résolution.} Réduisons au même dénominateur :
\[ \frac{270v - 270(v-15)}{v(v-15)} = 1 \iff \frac{270 \times 15}{v(v-15)} = 1 \iff v(v-15) = \num{4050}. \]
D'où $v^2 - 15v - \num{4050} = 0$. Discriminant : $\Delta = 225 + \num{16200} = \num{16425}$. Or $\num{16425} = 225 \times 73$, donc $\sqrt{\Delta} = 15\sqrt{73}$ et
\[ v = \frac{15 \pm 15\sqrt{73}}{2}. \]
La racine $\dfrac{15 - 15\sqrt{73}}{2} < 0$ est rejetée (une vitesse est positive). Il reste $v = \dfrac{15 + 15\sqrt{73}}{2} \approx \dfrac{15 + 128{,}2}{2} \approx 71{,}6$ km/h.
\textbf{Conclusion.} Le transporteur roule à environ \SI{71,6}{km/h} à l'aller et \SI{56,6}{km/h} au retour. Vérification : $\dfrac{270}{56{,}6} - \dfrac{270}{71{,}6} \approx 4{,}77 - 3{,}77 = 1$ h.
\end{exoresolu}

\courssub{Complément pour la Terminale : position relative de deux courbes}

\begin{savaistu}[Étudier le signe d'une différence — hors programme strict]
Pour savoir si la courbe d'une fonction $f$ est au-dessus ou au-dessous de celle d'une fonction $g$, on étudie le \cle{signe de la différence} $f(x) - g(x)$ :
\begin{itemize}
\item si $f(x) - g(x) > 0$, la courbe de $f$ est \emph{au-dessus} de celle de $g$ ;
\item si $f(x) - g(x) < 0$, elle est \emph{en dessous}.
\end{itemize}
Exemple : comparer $f(x) = x^2$ et $g(x) = 3x - 2$. On étudie $f(x) - g(x) = x^2 - 3x + 2 = (x-1)(x-2)$, positif à l'extérieur de $[1\,;\,2]$, négatif à l'intérieur. Cette technique, qui n'utilise que le signe d'un trinôme, deviendra centrale en Terminale pour l'étude des fonctions. \emph{(Complément culturel, hors programme strict de Première A.)}
\end{savaistu}

\courssub{Exercices de synthèse}

\begin{exercice}[Le champ rectangulaire]
Un champ rectangulaire a un périmètre de \SI{100}{m} et une aire de \SI{600}{m^2}. Déterminer ses dimensions.
\end{exercice}

\begin{corrige}[Exercice 1]
Soient $L$ et $l$ les dimensions. On a $2(L+l) = 100$ donc $L + l = 50$, et $Ll = 600$. $L$ et $l$ sont les racines de $X^2 - 50X + 600 = 0$. $\Delta = \num{2500} - \num{2400} = 100$, d'où $X_1 = \dfrac{50 - 10}{2} = 20$ et $X_2 = \dfrac{50 + 10}{2} = 30$. Les dimensions sont \SI{20}{m} et \SI{30}{m}. Vérification : $2(20+30) = 100$ et $20 \times 30 = 600$.
\end{corrige}

\begin{exercice}[Budget d'un cybercafé]
Un cybercafé de Yaoundé a des charges fixes de \num{45000} FCFA par mois et dépense $100$ FCFA par heure de connexion vendue. Il facture l'heure $250$ FCFA. Combien d'heures $h$ doit-il vendre au minimum par mois pour être bénéficiaire ?
\end{exercice}

\begin{corrige}[Exercice 2]
Le bénéfice est $B(h) = 250h - 100h - \num{45000} = 150h - \num{45000}$. On veut $B(h) > 0$, soit $h > 300$. Le cybercafé doit vendre \textbf{plus de $300$ heures} de connexion par mois pour être bénéficiaire (le seuil de rentabilité est exactement $300$ h).
\end{corrige}

\begin{exercice}[Partage d'une prime]
Deux employés se partagent une prime de \num{120000} FCFA. Le produit des deux parts (en dizaines de milliers de francs) vaut $32$ et leur somme vaut $12$ (toujours en dizaines de milliers). Quelles sont les deux parts ?
\end{exercice}

\begin{corrige}[Exercice 3]
Les parts $x$ et $y$ (en dizaines de milliers) vérifient $x + y = 12$ et $xy = 32$ : ce sont les racines de $X^2 - 12X + 32 = 0$. $\Delta = 144 - 128 = 16$, d'où $X_1 = \dfrac{12 - 4}{2} = 4$ et $X_2 = \dfrac{12 + 4}{2} = 8$. Les parts sont \num{40000} FCFA et \num{80000} FCFA. Vérification : $\num{40000} + \num{80000} = \num{120000}$ FCFA.
\end{corrige}

\begin{aretenir}[L'essentiel de la section]
\begin{itemize}
\item Tout problème concret se résout en \textbf{5 étapes} : choix de l'inconnue (avec son ensemble de valeurs), mise en équation, résolution, critique des solutions, conclusion rédigée.
\item « Au moins », « au plus », « rentable » $\Rightarrow$ \cle{inéquation} et tableau de signes.
\item Somme et produit de deux nombres $\Rightarrow$ équation $X^2 - SX + P = 0$.
\item Une solution négative pour une longueur, un prix ou une vitesse est \textbf{toujours rejetée}, en le justifiant explicitement.
\item La forme canonique donne directement le maximum ou le minimum d'une quantité du second degré — pont vers la dérivation de Terminale.
\end{itemize}
\end{aretenir}

\newpage

\coursec{Activité d'intégration}

\courssub{Objectifs de l'activité}

Cette activité d'intégration te place dans une situation authentique de la vie rurale camerounaise : l'aménagement d'un champ. En travaillant en groupe, tu devras mobiliser simultanément plusieurs savoir-faire de la leçon pour :

\begin{itemize}
\item modéliser une contrainte géométrique (un périmètre de clôture) par une expression algébrique ;
\item écrire un polynôme du second degré sous \cle{forme canonique} et en déduire un extremum ;
\item résoudre une \cle{inéquation du second degré} à l'aide du discriminant et d'un tableau de signes ;
\item résoudre un \cle{système somme-produit} en se ramenant à l'équation $X^2 - SX + P = 0$ ;
\item critiquer des solutions mathématiques au regard du contexte concret (valeurs impossibles, hors domaine).
\end{itemize}

\begin{attention}[Compétences évaluées]
L'objectif n'est pas seulement de trouver les bonnes réponses, mais de \emph{justifier chaque étape} : conditions d'existence, choix de la méthode, cohérence des résultats avec la réalité du problème. Au Probatoire, une réponse sans justification ni conclusion rédigée perd des points.
\end{attention}

\courssub{Situation}

\textbf{Le champ de Pa'a Essomba.}

Pa'a Essomba est un cultivateur d'Okola, dans la région du Centre. Il possède un terrain rectangulaire le long de la route, dont l'un des grands côtés est déjà bordé par un mur en parpaings appartenant à son voisin. Il souhaite clôturer les \emph{trois autres côtés} du champ avec du fil de fer, afin de protéger ses plants de maïs et de manioc des chèvres du quartier.

Chez le quincaillier d'Okola, il a acheté exactement $120$~mètres de fil de fer, et il ne veut pas en acheter davantage.

On note $x$ (en mètres) la \textbf{largeur} du champ, c'est-à-dire la longueur de chacun des deux côtés perpendiculaires au mur, et $L$ (en mètres) la \textbf{longueur} du champ, c'est-à-dire le côté parallèle au mur. Le mur remplace la clôture sur un des deux côtés de longueur $L$.

Pa'a Essomba se pose quatre questions :

\begin{enumerate}
\item Quelle est l'aire de son champ en fonction de $x$ ?
\item Quelles dimensions doit-il choisir pour que l'aire du champ soit \textbf{maximale} ?
\item Il estime qu'une aire d'au moins $\num{1600}$~m\textsuperscript{2} est nécessaire pour rentabiliser sa culture. Pour quelles valeurs de $x$ cette condition est-elle satisfaite ?
\item Pa'a Essomba veut partager le champ entre ses deux fils, en deux parcelles rectangulaires dont les aires (en m\textsuperscript{2}) sont deux nombres positifs dont la \textbf{somme} vaut $\num{1800}$ et le \textbf{produit} vaut $\num{720000}$. Quelles seront les aires des deux parcelles ?
\end{enumerate}

\begin{popfigure}
\begin{tikzpicture}[scale=0.09]
\fill[pattern=north east lines, pattern color=popmuted] (-4,-6) rectangle (64,0);
\draw[very thick, popdark] (-4,0) -- (64,0) node[midway, above, text=popdark]{\small Mur du voisin (non clôturé)};
\draw[very thick, poporange] (0,0) -- (0,40) node[midway, left, text=poporange]{\small $x$};
\draw[very thick, poporange] (60,0) -- (60,40) node[midway, right, text=poporange]{\small $x$};
\draw[very thick, poporange] (0,40) -- (60,40) node[midway, below, text=poporange]{\small $L$};
\node[text=popblue] at (30,20) {\small Champ de Pa'a Essomba};
\end{tikzpicture}
\legende{Le champ rectangulaire adossé au mur : seuls trois côtés sont clôturés.}
\end{popfigure}

\courssub{Consignes}

Travail en groupes de quatre élèves. Chaque groupe rédige ses réponses sur une feuille, avec toutes les justifications.

\begin{enumerate}
\item \textbf{Modélisation.} En utilisant la contrainte sur le fil de fer, exprimer $L$ en fonction de $x$. En déduire que l'aire du champ, en m\textsuperscript{2}, est donnée par $A(x) = -2x^2 + 120x$. Préciser l'ensemble des valeurs de $x$ qui ont un sens dans ce problème.
\item \textbf{Optimisation.} Écrire $A(x)$ sous forme canonique. En déduire la valeur de $x$ qui rend l'aire maximale, puis les dimensions du champ et l'aire maximale correspondante.
\item \textbf{Contrainte de rentabilité.} Montrer que la condition \og aire au moins égale à $\num{1600}$~m\textsuperscript{2} \fg{} se traduit par l'inéquation $-2x^2 + 120x - 1600 \geq 0$. Calculer le discriminant du trinôme, résoudre cette inéquation à l'aide d'un tableau de signes, puis conclure en tenant compte du domaine trouvé à la question 1.
\item \textbf{Partage entre les deux fils.} On cherche deux nombres réels positifs $u$ et $v$ tels que $u + v = \num{1800}$ et $u \times v = \num{720000}$. Justifier que $u$ et $v$ sont les solutions de l'équation $X^2 - 1800X + 720000 = 0$, puis résoudre cette équation. Conclure sur les aires des deux parcelles.
\item \textbf{Critique.} Un élève du groupe affirme : \og on pourrait aussi choisir $x = 70$~m, cela donnerait un champ plus grand \fg{}. Expliquer pourquoi cette proposition est absurde dans le contexte.
\end{enumerate}

\courssub{Ressources à mobiliser}

\begin{aretenir}[Boîte à outils de la leçon]
\begin{itemize}
\item \textbf{Forme canonique} : pour $a \neq 0$, $ax^2 + bx + c = a\left(x + \dfrac{b}{2a}\right)^2 - \dfrac{b^2 - 4ac}{4a}$. Si $a < 0$, le polynôme admet un \cle{maximum} en $x = -\dfrac{b}{2a}$.
\item \textbf{Discriminant} : $\Delta = b^2 - 4ac$. Si $\Delta > 0$, deux racines $x_1 = \dfrac{-b - \sqrt{\Delta}}{2a}$ et $x_2 = \dfrac{-b + \sqrt{\Delta}}{2a}$.
\item \textbf{Signe du trinôme} : si $\Delta > 0$, $ax^2+bx+c$ est du signe de $a$ à l'extérieur des racines, du signe contraire entre les racines.
\item \textbf{Système somme-produit} : deux nombres de somme $S$ et de produit $P$ sont les solutions de $X^2 - SX + P = 0$ (ils existent si et seulement si $S^2 - 4P \geq 0$).
\end{itemize}
\end{aretenir}

\courssub{Démarche et corrigé commenté}

\begin{exoresolu}[Le champ de Pa'a Essomba : corrigé complet]
Énoncé : voir la situation ci-dessus. On résout les cinq consignes dans l'ordre.
\tcblower
\textbf{Question 1 : modélisation.}

Le fil de fer clôture deux largeurs et une longueur (le mur remplace l'autre longueur). La contrainte s'écrit :
\[ 2x + L = 120 \quad \text{donc} \quad L = 120 - 2x. \]
L'aire du rectangle est alors :
\[ A(x) = x \times L = x(120 - 2x) = -2x^2 + 120x. \]
Pour que le champ existe, il faut $x > 0$ et $L > 0$, c'est-à-dire $120 - 2x > 0$, soit $x < 60$. Le domaine du problème est donc $x \in \,]0\,; 60[$.

\textbf{Question 2 : optimisation par la forme canonique.}

On factorise par $a = -2$ :
\[ A(x) = -2\left(x^2 - 60x\right). \]
Or $x^2 - 60x = (x - 30)^2 - 900$, donc :
\[ A(x) = -2\left[(x-30)^2 - 900\right] = -2(x-30)^2 + 1800. \]
Comme $-2 < 0$, la quantité $-2(x-30)^2$ est toujours négative ou nulle, et elle est maximale (nulle) lorsque $x = 30$. L'aire admet donc un \textbf{maximum} :
\[ A(30) = 1800 \ \text{m}^2, \quad \text{atteint pour } x = 30 \ \text{m}. \]
La longueur correspondante est $L = 120 - 2 \times 30 = 60$~m. \textbf{Conclusion :} Pa'a Essomba doit choisir un champ de $30$~m de largeur et $60$~m de longueur ; l'aire maximale est $\num{1800}$~m\textsuperscript{2}.

\textbf{Question 3 : inéquation de rentabilité.}

La condition $A(x) \geq \num{1600}$ s'écrit :
\[ -2x^2 + 120x \geq 1600 \iff -2x^2 + 120x - 1600 \geq 0. \]
On étudie le trinôme $f(x) = -2x^2 + 120x - 1600$ avec $a = -2$, $b = 120$, $c = -1600$ :
\[ \Delta = 120^2 - 4 \times (-2) \times (-1600) = 14400 - 12800 = 1600. \]
$\Delta > 0$, donc deux racines :
\[ x_1 = \frac{-120 - \sqrt{1600}}{2 \times (-2)} = \frac{-120 - 40}{-4} = 40, \qquad x_2 = \frac{-120 + 40}{-4} = 20. \]
On dresse le tableau de variations de $f$ sur $\mathbb{R}$ (les valeurs remarquables sont les racines $20$, $40$ et le sommet $30$) :

\begin{center}
\begin{tabular}{|c|ccccccccc|}\hline
$x$ & $-\infty$ & & $20$ & & $30$ & & $40$ & & $+\infty$ \\ \hline
$f'(x)$ & & $+$ & & $+$ & $0$ & $-$ & & $-$ & \\ \hline
$f(x)$ & $-\infty$ & $\nearrow$ & $0$ & $\nearrow$ & $200$ & $\searrow$ & $0$ & $\searrow$ & $-\infty$ \\ \hline
\end{tabular}
\end{center}

Le trinôme est positif ou nul ($f(x) \geq 0$) sur l'intervalle $[20\,; 40]$. En intersectant avec le domaine du problème $]0\,; 60[$, on obtient l'ensemble des solutions :
\[ \mathcal{S} = [20\,; 40]. \]
\textbf{Conclusion :} l'aire dépasse (ou égale) $\num{1600}$~m\textsuperscript{2} dès que la largeur $x$ est comprise entre $20$~m et $40$~m. Le choix optimal $x = 30$~m de la question 2 appartient bien à cet intervalle.

\textbf{Question 4 : partage somme-produit.}

On cherche $u$ et $v$ tels que $u + v = S = \num{1800}$ et $uv = P = \num{720000}$. D'après la propriété somme-produit, $u$ et $v$ sont les racines de :
\[ X^2 - SX + P = 0 \iff X^2 - 1800X + 720000 = 0. \]
Discriminant :
\[ \Delta = 1800^2 - 4 \times 720000 = \num{3240000} - \num{2880000} = \num{360000}, \qquad \sqrt{\Delta} = 600. \]
Les deux solutions sont :
\[ X_1 = \frac{1800 - 600}{2} = 600, \qquad X_2 = \frac{1800 + 600}{2} = 1200. \]
\textbf{Conclusion :} les deux parcelles auront des aires de $600$~m\textsuperscript{2} et $\num{1200}$~m\textsuperscript{2}. On vérifie : $600 + 1200 = \num{1800}$ et $600 \times 1200 = \num{720000}$. Remarquons que la somme $\num{1800}$~m\textsuperscript{2} correspond exactement à l'aire maximale du champ : le partage n'est possible que si Pa'a Essomba choisit les dimensions optimales !

\textbf{Question 5 : critique.}

Si $x = 70$~m, alors $L = 120 - 2 \times 70 = -20$~m : une longueur négative, ce qui est impossible. Le fil de $120$~m ne suffirait même pas à clôturer les deux largeurs ($2 \times 70 = 140 > 120$). Cette valeur est hors du domaine $]0\,; 60[$ : une solution mathématiquement calculable peut être \emph{hors contexte}, et il faut toujours confronter les résultats à la réalité du problème.
\end{exoresolu}

\courssub{Bilan de l'activité}

Cette activité a permis de mettre en évidence plusieurs idées essentielles de la leçon :

\begin{itemize}
\item La \cle{forme canonique} n'est pas un simple exercice de style : c'est l'outil naturel pour résoudre un problème d'\textbf{optimisation} (ici, maximiser une aire), car elle fait apparaître directement l'extremum et la valeur où il est atteint.
\item Le \cle{discriminant} et le \cle{tableau de signes} (ou de variations) du trinôme permettent de traiter des \textbf{contraintes} (aire minimale exigée) en résolvant des inéquations du second degré.
\item Les \cle{systèmes somme-produit} se ramènent toujours à une équation $X^2 - SX + P = 0$ : c'est la méthode standard pour les problèmes de partage, de dimensions, de nombres cherchés.
\item Enfin, la modélisation impose un \textbf{domaine de validité} : toute solution doit être vérifiée dans le contexte. Les mathématiques donnent des candidats ; le bon sens et les contraintes concrètes font le tri.
\end{itemize}

\begin{savaistu}[Et dans la vraie vie ?]
Les problèmes d'optimisation sous contrainte comme celui de Pa'a Essomba sont partout : dimensionnement des parcelles agricoles, emballages de volume maximal, tarification des forfaits MTN ou Orange pour maximiser les recettes. Le cas du rectangle adossé à un mur est un classique : à périmètre de clôture fixé, la longueur optimale est le double de la largeur. Retiens ce réflexe : face à \og maximal \fg{} ou \og minimal \fg{}, pense forme canonique.
\end{savaistu}

\newpage

\coursec{Méthodes et exercices résolus}

\coursec{Équations et inéquations se ramenant à un quotient $\frac{ax+b}{cx+d}$}

\begin{definition}[Fonction homographique, équation et inéquation associées]
    Soient $a,b,c,d \in \mathbb{R}$ avec $c \neq 0$ ou $d \neq 0$ (pour éviter le cas constant). On appelle \textbf{fonction homographique} la fonction $f$ définie sur $\mathbb{R} \setminus \{-\frac{d}{c}\}$ par $f(x) = \frac{ax+b}{cx+d}$.
    Une \textbf{équation homographique} est une équation de la forme $\frac{ax+b}{cx+d} = 0$ (ou $= k$, ramenée à $0$).
    Une \textbf{inéquation homographique} est une inéquation de la forme $\frac{ax+b}{cx+d} \lesseqgtr 0$.
\end{definition}

\begin{methode}[Résoudre une équation $\frac{ax+b}{cx+d}=0$]
    \textbf{Principe fondamental :} Un quotient de réels est nul \textbf{si et seulement si} son numérateur est nul \textbf{et} son dénominateur est non nul.
    \begin{enumerate}
        \item \textbf{Condition de définition (obligatoire) :} Écrire $cx+d \neq 0 \iff x \neq -\frac{d}{c}$ (si $c \neq 0$ ; si $c=0$, l'équation est du premier degré standard).
        \item \textbf{Résolution du numérateur :} Résoudre $ax+b=0 \iff x = -\frac{b}{a}$ (si $a \neq 0$ ; si $a=0$, l'équation n'a pas de solution sauf si $b=0$ mais alors dénominateur nul).
        \item \textbf{Vérification d'admissibilité :} La solution trouvée doit satisfaire la condition de définition.
        \item \textbf{Conclusion :} Écrire l'ensemble solution $\mathcal{S}$ (singleton ou $\varnothing$).
    \end{enumerate}
    \textbf{Réflexe Probatoire :} Toujours écrire explicitement la condition $cx+d \neq 0$ \textbf{avant} de résoudre. Si la solution du numérateur annule le dénominateur, l'équation n'a \textbf{aucune solution} ($\mathcal{S}=\varnothing$).
\end{methode}

\begin{exoresolu}[Équation homographique simple -- Type Probatoire]
    Résoudre dans $\mathbb{R}$ l'équation : $\frac{3x-5}{2x+1}=0$.
    \tcblower
    \textbf{Résolution.}
    \begin{enumerate}
        \item \textbf{Condition de définition :} Le dénominateur ne doit pas être nul.
        \[2x+1 \neq 0 \iff 2x \neq -1 \iff x \neq -\frac{1}{2}.\]
        \item \textbf{Résolution du numérateur :}
        \[3x-5=0 \iff 3x=5 \iff x=\frac{5}{3}.\]
        \item \textbf{Vérification :} $\frac{5}{3} \neq -\frac{1}{2}$. La valeur trouvée est admissible.
    \end{enumerate}
    \textbf{Conclusion :} L'ensemble solution est $\mathcal{S}=\left\{\frac{5}{3}\right\}$.
\end{exoresolu}

\begin{methode}[Résoudre une inéquation $\frac{ax+b}{cx+d} \lesseqgtr 0$ par tableau de signes]
    \textbf{Principe :} Le signe d'un quotient est le produit des signes du numérateur et du dénominateur (règle des signes). On dresse un tableau de signes rigoureux.
    \begin{enumerate}
        \item \textbf{Ramener à $0$ :} Mettre l'inéquation sous la forme $\frac{N(x)}{D(x)} \lesseqgtr 0$ (numérateur et dénominateur du premier degré).
        \item \textbf{Trouver les valeurs racines :}
        \begin{itemize}
            \item Racine du numérateur : $x_N = -\frac{b}{a}$ (zéro de la fonction).
            \item Racine du dénominateur : $x_D = -\frac{d}{c}$ (valeur \textbf{interdite}, asymptote verticale).
        \end{itemize}
        \item \textbf{Construire le tableau :}
        \begin{itemize}
            \item Placer $x_D$ et $x_N$ dans l'ordre croissant sur la première ligne.
            \item \textbf{Marquer la valeur interdite $x_D$ par une double barre $||$ (ou trait de suspension).}
            \item Déterminer le signe de $ax+b$ et $cx+d$ sur chaque intervalle (signe du coefficient directeur pour $x \to +\infty$).
            \item Déduire le signe du quotient par multiplication (ligne du bas).
        \end{itemize}
        \item \textbf{Lire la solution :}
        \begin{itemize}
            \item Inégalité \textbf{stricte} ($<$ ou $>$) : valeurs interdites exclues, racine du numérateur exclue.
            \item Inégalité \textbf{large} ($\le$ ou $\ge$) : racine du numérateur \textbf{incluse} (crochets), racine du dénominateur \textbf{toujours exclue}.
        \end{itemize}
    \end{enumerate}
    \textbf{Astuce rapide :} Pour $x \to +\infty$, le signe du quotient est $\frac{\text{signe}(a)}{\text{signe}(c)}$.
\end{methode}

\begin{exoresolu}[Inéquation homographique -- Modélisation tarification MTN/Orange]
    Un forfait internet coûte $1500$ FCFA pour les $2$ premiers Go, puis $300$ FCFA par Go supplémentaire. Soit $x$ le nombre de Go supplémentaires ($x \ge 0$). Le coût moyen par Go total est $m(x)=\frac{1500+300x}{2+x}$.
    Déterminer les valeurs de $x$ pour lesquelles le coût moyen est \textbf{strictement inférieur} à $400$ FCFA/Go.
    \tcblower
    \textbf{Résolution.}
    On cherche $x \ge 0$ tel que $\frac{1500+300x}{2+x} < 400$.
    \begin{enumerate}
        \item \textbf{Ramener à $0$ et simplifier :}
        \[\frac{1500+300x}{2+x} - 400 < 0 \iff \frac{1500+300x - 400(2+x)}{2+x} < 0\]
        \[\iff \frac{1500+300x - 800 - 400x}{2+x} < 0 \iff \frac{700 - 100x}{2+x} < 0.\]
        Factoriser par $100>0$ (ne change pas le signe) : $\frac{100(7 - x)}{x+2} < 0 \iff \frac{7 - x}{x+2} < 0$.
        \item \textbf{Valeurs remarquables :}
        \begin{itemize}
            \item Numérateur $7-x=0 \Rightarrow x_N = 7$.
            \item Dénominateur $x+2=0 \Rightarrow x_D = -2$ (\textbf{valeur interdite}).
        \end{itemize}
        \item \textbf{Tableau de signes :} (2 valeurs $\to$ 5 colonnes pour les bornes + 3 intervalles = 8 colonnes au total)
        \begin{center}
        \begin{tabular}{|c|ccccccc|}\hline
        $x$ & $-\infty$ & & $-2$ & & $7$ & & $+\infty$ \\ \hline
        $7-x$ & $+$ & & $+$ & & $0$ & & $-$ \\
        $x+2$ & $-$ & & $||$ & & $+$ & & $+$ \\ \hline
        $\frac{7-x}{x+2}$ & $-$ & & $||$ & & $0$ & & $+$ \\ \hline
        \end{tabular}
        \end{center}
        \item \textbf{Lecture :} L'inéquation est stricte ($<0$). Le quotient est négatif sur $]-\infty; -2[$ et $]7; +\infty[$.
        \item \textbf{Contrainte du réel :} $x \ge 0$ (Go supplémentaires).
        Intersection : $[0; +\infty[ \cap \left(]-\infty; -2[ \cup ]7; +\infty[\right) = ]7; +\infty[$.
    \end{enumerate}
    \textbf{Conclusion :} Le coût moyen est inférieur à $400$ FCFA/Go pour \textbf{plus de $7$ Go supplémentaires} (soit un total de plus de $9$ Go).
\end{exoresolu}

\begin{popfigure}
\begin{tikzpicture}[scale=0.8, domain=-5:10]
    \draw[->] (-5.5,0) -- (10.5,0) node[right] {$x$ (Go suppl.)};
    \draw[->] (0,-5.5) -- (0,5.5) node[above] {$m(x)$};
    \draw[dashed, popmuted] (-2,-5) -- (-2,5) node[above left] {$x=-2$};
    \begin{scope}
    \clip (-5.5,-5.5) rectangle (10.5,5.5);
    \draw[popcurve, thick, samples=150, domain=-5:-2.25] plot (\x, {(7-\x)/(\x+2)});
    \draw[popcurve, thick, samples=150, domain=-1.75:10] plot (\x, {(7-\x)/(\x+2)});
    \end{scope}
    \draw[popgreen, thick] (7,0) circle (3pt) node[below right] {$(7,0)$};
    \draw[poporange, dashed] (0,3.5) -- (10,3.5) node[right] {$y=3.5$};
    \node[popblue] at (8, -1) {$<0$};
    \node[popblue] at (3, 2) {$>0$};
    \node[popblue] at (-4, -2) {$<0$};
\end{tikzpicture}
\legende{Représentation graphique de $m(x)-400 \propto \frac{7-x}{x+2}$. La zone $x>7$ (contexte $x\ge 0$) donne un coût moyen $<400$.}
\end{popfigure}

\coursec{Forme canonique d'un polynôme du second degré}

\begin{definition}[Forme canonique]
    Soit $P(x) = ax^2+bx+c$ un polynôme du second degré ($a \neq 0$). Sa \textbf{forme canonique} est :
    \[P(x) = a\left[\left(x + \frac{b}{2a}\right)^2 - \frac{\Delta}{4a^2}\right] = a\left(x - \alpha\right)^2 + \beta\]
    où $\alpha = -\frac{b}{2a}$ (abscisse du sommet), $\beta = -\frac{\Delta}{4a}$ (ordonnée du sommet) et $\Delta = b^2-4ac$ (discriminant).
\end{definition}

\begin{experience}[Découverte de la forme canonique -- Méthode du carré parfait]
    On veut écrire $P(x) = 2x^2 - 8x + 5$ sous la forme $a(x-\alpha)^2 + \beta$.
    \begin{enumerate}
        \item Factoriser $a=2$ sur les termes en $x$ : $P(x) = 2(x^2 - 4x) + 5$.
        \item Compléter le carré : $x^2 - 4x = (x-2)^2 - 4$. (On ajoute et soustrait $(\frac{-4}{2})^2 = 4$).
        \item $P(x) = 2\left[(x-2)^2 - 4\right] + 5 = 2(x-2)^2 - 8 + 5 = 2(x-2)^2 - 3$.
        \item \textbf{Forme canonique :} $P(x) = 2(x-2)^2 - 3$. Sommet $S(2; -3)$. Comme $a=2>0$, $P$ a un minimum $-3$ en $x=2$.
    \end{enumerate}
\end{experience}

\begin{methode}[Mettre $ax^2+bx+c$ sous forme canonique (Algorithme)]
    \begin{enumerate}
        \item Factoriser $a$ sur les deux premiers termes : $a\left(x^2+\frac{b}{a}x\right)+c$.
        \item Compléter le carré : ajouter et soustraire $\left(\frac{b}{2a}\right)^2$ à l'intérieur.
        \item Écrire le carré remarquable : $a\left[\left(x+\frac{b}{2a}\right)^2 - \frac{b^2}{4a^2}\right]+c$.
        \item Simplifier la constante : $-\frac{ab^2}{4a^2}+c = \frac{-b^2+4ac}{4a} = -\frac{\Delta}{4a}$.
        \item Écrire le résultat final : $a\left(x+\frac{b}{2a}\right)^2 - \frac{\Delta}{4a}$.
    \end{enumerate}
    \textbf{Utilité :} Lecture immédiate du sommet de la parabole ($-\frac{b}{2a}; -\frac{\Delta}{4a}$), du sens de variation, du signe de $\Delta$ (via $\beta$), et résolution aisée de $P(x)=k$.
\end{methode}

\begin{exoresolu}[Forme canonique et résolution -- Application trajectoire projectile]
    Un projectile est lancé depuis le sol. Sa hauteur (en mètres) en fonction du temps $t$ (en secondes) est modélisée par $h(t)=-5t^2+20t$.
    \begin{enumerate}
        \item Mettre $h(t)$ sous forme canonique. En déduire la hauteur maximale et l'instant où elle est atteinte.
        \item Résoudre $h(t)=15$. Interpréter.
    \end{enumerate}
    \tcblower
    \textbf{Solution.}
    \begin{enumerate}
        \item $h(t)=-5t^2+20t = -5(t^2-4t)$.
        Compléter le carré : $t^2-4t = (t-2)^2 - 4$.
        Donc $h(t) = -5\left[(t-2)^2 - 4\right] = -5(t-2)^2 + 20$.
        \textbf{Forme canonique :} $h(t) = -5(t-2)^2 + 20$.
        Comme $a=-5<0$, la parabole est tournée vers le bas : \textbf{maximum} $20$ atteint pour $t=2$.
        \textit{Hauteur max $= 20$ m à $t=2$ s.}
        \item $h(t)=15 \iff -5t^2+20t=15 \iff -5t^2+20t-15=0$.
        Diviser par $-5$ : $t^2-4t+3=0$. ($a=1, b=-4, c=3$).
        $\Delta = (-4)^2 - 4\times 1 \times 3 = 16-12=4 > 0$.
        $\sqrt{\Delta}=2$.
        $t_1 = \frac{4-2}{2}=1$ ; $t_2 = \frac{4+2}{2}=3$.
        \textbf{Interprétation :} Le projectile est à $15$ m de hauteur à deux instants : \textbf{$t=1$ s (montée)} et \textbf{$t=3$ s (descente)}.
    \end{enumerate}
\end{exoresolu}

\begin{popfigure}
\begin{tikzpicture}[scale=0.6, domain=0:4]
    \draw[->] (-0.5,0) -- (4.5,0) node[right] {$t$ (s)};
    \draw[->] (0,-0.5) -- (0,22) node[above] {$h$ (m)};
    \draw[popcurve, thick, samples=200] plot (\x, {-5*\x*\x+20*\x});
    \draw[dashed, popmuted] (2,0) -- (2,20) node[above] {$t=2$};
    \draw[dashed, popmuted] (0,20) -- (2,20) node[left] {$h=20$};
    \draw[popgreen, thick] (2,20) circle (4pt) node[above right] {Sommet $S(2;20)$};
    \draw[poporange, thick] (1,15) circle (4pt) node[below left] {$(1;15)$};
    \draw[poporange, thick] (3,15) circle (4pt) node[below right] {$(3;15)$};
    \draw[dashed, poporange] (1,0) -- (1,15) (3,0) -- (3,15) (0,15) -- (1,15) (0,15) -- (3,15);
\end{tikzpicture}
\legende{Trajectoire du projectile $h(t)=-5t^2+20t$. Hauteur max 20m à 2s. Passage à 15m aux instants 1s et 3s.}
\end{popfigure}

\coursec{Discriminant et résolution des équations du second degré}

\begin{propriete}[Discriminant $\Delta$ et nombre de solutions réelles]
    Soit $P(x) = ax^2+bx+c$ avec $a \neq 0$ et $\Delta = b^2 - 4ac$.
    \begin{itemize}
        \item Si $\Delta > 0$ : L'équation $P(x)=0$ admet \textbf{deux solutions distinctes} :
        \[x_1 = \frac{-b - \sqrt{\Delta}}{2a}, \quad x_2 = \frac{-b + \sqrt{\Delta}}{2a} \quad (x_1 < x_2).\]
        \item Si $\Delta = 0$ : L'équation $P(x)=0$ admet \textbf{une solution double} (racine double) :
        \[x_0 = -\frac{b}{2a}.\]
        \item Si $\Delta < 0$ : L'équation $P(x)=0$ n'admet \textbf{aucune solution réelle} ($\mathcal{S} = \varnothing$).
    \end{itemize}
    \textbf{Exigible au Probatoire :} La formule des racines, le lien $\Delta = b^2-4ac$, et l'interprétation graphique (intersections parabole / axe $Ox$).
\end{propriete}

\begin{attention}[Pièges classiques sur le discriminant et les racines]
    \begin{enumerate}
        \item \textbf{Ordre des racines :} Par convention $x_1 \le x_2$. La formule $\frac{-b \pm \sqrt{\Delta}}{2a}$ donne $x_1$ avec $-$ et $x_2$ avec $+$ \textbf{seulement si $a>0$}. Si $a<0$, l'ordre s'inverse ! \textbf{Astuce :} Calculer les deux valeurs numériques, puis les ranger dans l'ordre croissant pour le tableau de signes.
        \item \textbf{Racine carrée :} $\sqrt{\Delta}$ désigne la racine carrée \textbf{positive} de $\Delta$. Ne pas écrire $\sqrt{\Delta} = \pm \dots$
        \item \textbf{Simplification :} Si $\Delta$ est un carré parfait ($49, 64, 100\dots$), simplifier $\sqrt{\Delta}$ avant de calculer $x_1, x_2$.
        \item \textbf{Équation réduite :} Diviser par $a$ si $a \neq 1$ pour simplifier les calculs ($x^2 + \frac{b}{a}x + \frac{c}{a} = 0$), mais attention : $\Delta$ change ($\Delta' = (\frac{b}{a})^2 - 4\frac{c}{a} = \frac{\Delta}{a^2}$). Préférer $\Delta$ initial.
    \end{enumerate}
\end{attention}

\coursec{Factorisation et signe d'un polynôme du second degré}

\begin{propriete}[Factorisation d'un trinôme du second degré]
    Soit $P(x) = ax^2+bx+c$ ($a \neq 0$) et $\Delta = b^2-4ac$.
    \begin{itemize}
        \item Si $\Delta > 0$ : $P(x) = a(x-x_1)(x-x_2)$ où $x_1, x_2$ sont les deux racines réelles distinctes.
        \item Si $\Delta = 0$ : $P(x) = a(x-x_0)^2$ où $x_0$ est la racine double.
        \item Si $\Delta < 0$ : $P(x)$ \textbf{ne se factorise pas} dans $\mathbb{R}$ (pas de facteur du premier degré à coefficients réels).
    \end{itemize}
    \textbf{Remarque :} La factorisation est l'outil principal pour dresser le tableau de signes et résoudre les inéquations.
\end{propriete}

\begin{methode}[Dresser le tableau de signes de $ax^2+bx+c$]
    \begin{enumerate}
        \item Calculer $\Delta = b^2-4ac$.
        \item \textbf{Cas $\Delta < 0$ :} Le trinôme a le \textbf{signe constant de $a$} sur $\mathbb{R}$.
        \item \textbf{Cas $\Delta = 0$ :} Le trinôme a le signe de $a$ sur $\mathbb{R}$, \textbf{nul en $x_0 = -\frac{b}{2a}$}.
        \item \textbf{Cas $\Delta > 0$ :} Calculer $x_1 < x_2$. Le trinôme a le signe de $a$ \textbf{en dehors} des racines ($]-\infty; x_1[ \cup ]x_2; +\infty[$) et le signe \textbf{contraire de $a$} \textbf{entre} les racines ($]x_1; x_2[$).
        \item \textbf{Tableau type ($\Delta > 0$, $x_1 < x_2$) :}
        \begin{center}
        \begin{tabular}{|c|ccccccc|}\hline
        $x$ & $-\infty$ & & $x_1$ & & $x_2$ & & $+\infty$ \\ \hline
        $ax^2+bx+c$ & signe($a$) & & $0$ & signe($-a$) & $0$ & & signe($a$) \\ \hline
        \end{tabular}
        \end{center}
    \end{enumerate}
    \textbf{Mnémotechnique :} « Le trinôme a le signe de $a$ \textbf{en dehors} des racines, et le signe \textbf{contraire} de $a$ \textbf{entre} les racines. »
\end{methode}

\begin{exoresolu}[Factorisation, signes et inéquation du second degré]
    Soit $f(x)=2x^2-5x-3$.
    \begin{enumerate}
        \item Calculer $\Delta$ et factoriser $f(x)$.
        \item Dresser le tableau de signes de $f$.
        \item Résoudre $f(x) \le 0$.
    \end{enumerate}
    \tcblower
    \textbf{Solution.}
    \begin{enumerate}
        \item $a=2, b=-5, c=-3$.
        $\Delta = (-5)^2 - 4 \times 2 \times (-3) = 25 + 24 = 49$.
        $\sqrt{\Delta}=7$.
        $x_1 = \frac{5-7}{4} = -\frac{1}{2}$ ; $x_2 = \frac{5+7}{4} = 3$. (Ordre croissant vérifié : $-\frac{1}{2} < 3$).
        \textbf{Factorisation :} $f(x) = 2\left(x+\frac{1}{2}\right)(x-3) = (2x+1)(x-3)$.
        \item $a=2>0$, racines $-\frac{1}{2}$ et $3$. Signe $+$ en dehors, $-$ entre.
        \begin{center}
        \begin{tabular}{|c|ccccccc|}\hline
        $x$ & $-\infty$ & & $-\frac{1}{2}$ & & $3$ & & $+\infty$ \\ \hline
        $f(x)$ & $+$ & & $0$ & $-$ & $0$ & & $+$ \\ \hline
        \end{tabular}
        \end{center}
        \item $f(x) \le 0 \iff$ signes $-$ ou $0$ (inégalité large).
        Solution : $\left[-\frac{1}{2};\; 3\right]$.
    \end{enumerate}
\end{exoresolu}

\begin{popfigure}
\begin{tikzpicture}[scale=0.7, domain=-2:4]
    \draw[->] (-2.5,0) -- (4.5,0) node[right] {$x$};
    \draw[->] (0,-5.5) -- (0,5.5) node[above] {$y$};
    \draw[popcurve, thick, samples=200] plot (\x, {2*\x*\x-5*\x-3});
    \draw[popgreen, thick] (-0.5,0) circle (3pt) node[below] {$-\frac{1}{2}$};
    \draw[popgreen, thick] (3,0) circle (3pt) node[below] {$3$};
    \draw[poppurple, thick] (1.25,-4.125) circle (3pt) node[below right] {Sommet};
    \node[popblue] at (-1.5, 3) {$+$};
    \node[popred] at (1.25, -2) {$-$};
    \node[popblue] at (4, 3) {$+$};
\end{tikzpicture}
\legende{Graphique de $f(x)=2x^2-5x-3$. Signe $+$ en dehors des racines, $-$ entre. Solution de $f(x)\le 0$ : segment $[-\frac{1}{2}; 3]$.}
\end{popfigure}

\coursec{Systèmes d'équations se ramenant au second degré dans $\mathbb{R}^2$}

\begin{methode}[Résoudre un système $\begin{cases} E_1(x,y)=0 \\ E_2(x,y)=0 \end{cases}$ se ramenant au second degré]
    \textbf{Stratégie par substitution (la plus générale) :}
    \begin{enumerate}
        \item \textbf{Isoler une inconnue} (souvent $y$) dans l'équation la plus simple (souvent linéaire ou symétrique).
        \item \textbf{Substituer} cette expression dans l'autre équation.
        \item Obtenir une équation \textbf{à une inconnue} (généralement du second degré en $x$ ou $y$).
        \item \textbf{Résoudre} cette équation (Discriminant, factorisation, forme canonique).
        \item Pour chaque valeur trouvée, \textbf{calculer la valeur de l'autre inconnue} via l'expression de l'étape 1.
        \item \textbf{Vérifier} les couples $(x;y)$ dans le système initial (indispensable si on a élevé au carré, multiplié par une variable, ou simplifié par une expression pouvant être nulle).
        \item Écrire l'ensemble solution $\mathcal{S} \subset \mathbb{R}^2$ comme ensemble de couples : $\mathcal{S} = \{(x_1;y_1); (x_2;y_2); \dots\}$.
    \end{enumerate}
    \textbf{Cas particulier -- Système symétrique ($x+y=S$, $xy=P$) :}
    $x$ et $y$ sont les racines du trinôme $X^2 - SX + P = 0$. $\Delta = S^2 - 4P$.
    \begin{itemize}
        \item $\Delta < 0 \implies \mathcal{S} = \varnothing$.
        \item $\Delta \ge 0 \implies$ les solutions sont les couples $(x_1; x_2)$ et $(x_2; x_1)$.
    \end{itemize}
\end{methode}

\begin{exoresolu}[Système symétrique -- Problème « Somme et Produit »]
    Trouver deux nombres réels dont la somme est $10$ et le produit est $21$.
    \tcblower
    \textbf{Résolution.}
    Soient $x$ et $y$ ces deux nombres.
    \begin{enumerate}
        \item \textbf{Modélisation :}
        $\begin{cases} x+y=10 & (1) \\ xy=21 & (2) \end{cases}$
        \item \textbf{Méthode 1 (Substitution classique) :} D'après (1) : $y = 10-x$.
        Substitution dans (2) : $x(10-x)=21 \iff 10x-x^2=21 \iff x^2-10x+21=0$.
        \item \textbf{Méthode 2 (Trinôme associé) :} $x$ et $y$ sont racines de $X^2 - 10X + 21 = 0$. Même équation.
        \item $\Delta = (-10)^2 - 4\times 21 = 100-84=16$. $\sqrt{\Delta}=4$.
        $x_1 = \frac{10-4}{2}=3$ ; $x_2 = \frac{10+4}{2}=7$.
        \item Valeurs correspondantes de $y$ (par $y=10-x$) :
        Si $x=3 \Rightarrow y=7$.
        Si $x=7 \Rightarrow y=3$.
        \item Vérification : $3+7=10$, $3\times 7=21$. OK.
    \end{enumerate}
    \textbf{Conclusion :} Les deux nombres sont \textbf{$3$ et $7$}. L'ensemble solution du système est $\mathcal{S}=\{(3;7);(7;3)\}$.
\end{exoresolu}

\begin{exoresolu}[Système non symétrique -- Intersection droite/parabole]
    Résoudre dans $\mathbb{R}^2$ : $\begin{cases} y = x^2 - 2x - 3 \\ y = 2x - 5 \end{cases}$
    \tcblower
    \textbf{Résolution.}
    \begin{enumerate}
        \item Substitution : $x^2 - 2x - 3 = 2x - 5 \iff x^2 - 4x + 2 = 0$.
        \item $\Delta = 16 - 8 = 8 > 0$. $\sqrt{\Delta} = 2\sqrt{2}$.
        $x_1 = 2 - \sqrt{2}$, $x_2 = 2 + \sqrt{2}$.
        \item $y = 2x - 5$.
        $y_1 = 2(2-\sqrt{2})-5 = -1 - 2\sqrt{2}$.
        $y_2 = 2(2+\sqrt{2})-5 = -1 + 2\sqrt{2}$.
        \item Vérification rapide : $(2\pm\sqrt{2})^2 - 2(2\pm\sqrt{2}) - 3 = 4\pm4\sqrt{2}+2 -4\mp2\sqrt{2}-3 = -1\pm2\sqrt{2}$. OK.
    \end{enumerate}
    \textbf{Conclusion :} $\mathcal{S} = \{(2-\sqrt{2}; -1-2\sqrt{2}); (2+\sqrt{2}; -1+2\sqrt{2})\}$.
\end{exoresolu}

\coursec{Inéquations du second degré et problèmes de synthèse}

\begin{methode}[Résoudre une inéquation $ax^2+bx+c \lesseqgtr 0$]
    \textbf{Algorithme unifié :}
    \begin{enumerate}
        \item Calculer $\Delta = b^2-4ac$.
        \item \textbf{Selon $\Delta$ et $a$ :}
        \begin{itemize}
            \item $\Delta < 0$ : Signe constant = signe($a$).
            \item $\Delta = 0$ : Signe de $a$ partout, nul en $x_0$.
            \item $\Delta > 0$ : Racines $x_1 < x_2$. Signe de $a$ en dehors, signe opposé entre.
        \end{itemize}
        \item Lire la solution selon l'inégalité (crochets pour $\le/\ge$ aux racines du numérateur/trinôme ; parenthèses pour $</>$ et toujours pour valeurs interdites).
        \item \textbf{Écrire la solution sous forme d'intervalle(s) ou d'ensemble.}
    \end{enumerate}
    \textbf{Cas particuliers (gain de temps au Probatoire) :}
    \begin{itemize}
        \item $a>0, \Delta<0$ : $P(x) > 0 \ \forall x \in \mathbb{R}$ ; $P(x) \ge 0 \ \forall x \in \mathbb{R}$ ; $P(x) < 0$ ou $\le 0 \implies \varnothing$.
        \item $a<0, \Delta<0$ : $P(x) < 0 \ \forall x \in \mathbb{R}$ ; $P(x) \le 0 \ \forall x \in \mathbb{R}$ ; $P(x) > 0$ ou $\ge 0 \implies \varnothing$.
    \end{itemize}
\end{methode}

\begin{exoresolu}[Inéquation du second degré -- Optimisation recette vente]
    Un commerçant vend des ananas. S'il les vend à $x$ FCFA l'unité ($x>0$), il en vend $q(x)=100-2x$ par jour. Sa recette journalière est $R(x)=x(100-2x)$.
    Pour quels prix $x$ la recette est-elle \textbf{supérieure ou égale à $1200$ FCFA} ?
    \tcblower
    \textbf{Résolution.}
    \begin{enumerate}
        \item \textbf{Inéquation :} $R(x) \ge 1200 \iff x(100-2x) \ge 1200 \iff 100x-2x^2-1200 \ge 0$.
        \item \textbf{Simplifier :} Diviser par $-2$ (\textbf{changement de sens !}) : $x^2-50x+600 \le 0$.
        \item \textbf{Étude du trinôme :} $a=1, b=-50, c=600$. $\Delta = 2500 - 2400 = 100$. $\sqrt{\Delta}=10$.
        Racines : $x_1 = \frac{50-10}{2}=20$ ; $x_2 = \frac{50+10}{2}=30$.
        \item \textbf{Tableau de signes ($a=1>0$) :}
        \begin{center}
        \begin{tabular}{|c|ccccccc|}\hline
        $x$ & $-\infty$ & & $20$ & & $30$ & & $+\infty$ \\ \hline
        $x^2-50x+600$ & $+$ & & $0$ & $-$ & $0$ & & $+$ \\ \hline
        \end{tabular}
        \end{center}
        \item \textbf{Lecture :} $\le 0 \iff$ signes $-$ ou $0 \iff x \in [20; 30]$.
        \item \textbf{Contraintes du réel :} $x>0$ et $q(x)=100-2x \ge 0 \Rightarrow x \le 50$. $[20;30] \subset [0;50]$.
    \end{enumerate}
    \textbf{Conclusion :} Le prix doit être compris entre \textbf{$200$ et $300$ FCFA} (inclus) pour assurer une recette $\ge 1200$ FCFA.
    \textit{Note : L'énoncé donne $x$ en FCFA, le calcul donne 20 et 30. Il faut vérifier l'unité : si $x$ est en centaines de FCFA, alors 20 = 2000 FCFA. Ici, l'énoncé dit "$x$ FCFA", le résultat est 20 FCFA et 30 FCFA. Si le contexte implique des centaines, l'énoncé doit le préciser. On garde l'unité de l'énoncé.}
\end{exoresolu}

\coursec{Problèmes de synthèse et modélisation}

\begin{methode}[Résoudre un problème concret (Démarche en 5 étapes -- Standard Probatoire)]
    \begin{enumerate}
        \item \textbf{Variables :} Définir clairement les inconnues (lettres, unités, ensemble de définition). Ex: « Soit $x$ la longueur en mètres, $x>0$ ».
        \item \textbf{Modélisation :} Traduire les contraintes en équations ou inéquations. Justifier chaque traduction par une phrase.
        \item \textbf{Résolution mathématique :} Appliquer les méthodes vues (quotient, trinôme, système). Montrer les calculs (discriminant, factorisation, tableau de signes).
        \item \textbf{Interprétation / Filtrage :} Revenir au contexte. Éliminer les solutions mathématiques non physiques (longueur négative, temps négatif, nombre d'objets non entier si requis).
        \item \textbf{Conclusion :} Phrase réponse complète, avec unités, répondant exactement à la question posée.
    \end{enumerate}
\end{methode}

\begin{exoresolu}[Problème synthèse : Dimensions d'un terrain rectangulaire]
    Un agriculteur veut clôturer un terrain rectangulaire de $600\,\text{m}^2$ le long d'une rivière (donc 3 côtés à clôturer). Il dispose de $100\,\text{m}$ de grillage.
    Déterminer les dimensions possibles du terrain.
    \tcblower
    \textbf{Résolution.}
    \begin{enumerate}
        \item \textbf{Variables :} Soit $x$ la longueur (parallèle à la rivière) en mètres, $y$ la largeur (perpendiculaire) en mètres. $x>0, y>0$.
        \item \textbf{Contraintes :}
        \begin{itemize}
            \item Aire : $xy = 600$ \quad (1)
            \item Grillage (3 côtés) : $x + 2y = 100$ \quad (2)
        \end{itemize}
        \item \textbf{Système :} $\begin{cases} xy=600 \\ x+2y=100 \end{cases}$
        De (2) : $x = 100-2y$.
        Substitution dans (1) : $(100-2y)y = 600 \iff 100y-2y^2=600 \iff 2y^2-100y+600=0$.
        Diviser par $2$ : $y^2-50y+300=0$.
        \item \textbf{Résolution :} $\Delta = 2500 - 1200 = 1300 = 100 \times 13$. $\sqrt{\Delta}=10\sqrt{13} \approx 36,06$.
        $y_1 = \frac{50-10\sqrt{13}}{2} = 25-5\sqrt{13} \approx 6,97$.
        $y_2 = 25+5\sqrt{13} \approx 43,03$.
        \item \textbf{Dimensions correspondantes :}
        Si $y_1 \approx 6,97 \Rightarrow x_1 = 100-2y_1 \approx 86,06$.
        Si $y_2 \approx 43,03 \Rightarrow x_2 = 100-2y_2 \approx 13,94$.
        Les deux couples donnent des dimensions positives.
    \end{enumerate}
    \textbf{Conclusion :} Deux configurations possibles :
    \begin{itemize}
        \item Longueur $\approx 86,1$ m, Largeur $\approx 7,0$ m.
        \item Longueur $\approx 13,9$ m, Largeur $\approx 43,0$ m.
    \end{itemize}
    (Valeurs exactes : $(100-2(25\pm5\sqrt{13}); 25\pm5\sqrt{13})$).
\end{exoresolu}

\begin{savaistu}[Histoire et applications au Cameroun]
    \textbf{Al-Khwarizmi (vers 780-850) :} Son ouvrage \emph{Al-Kitab al-mukhtasar fi hisab al-jabr wa'l-muqabala} (« Le livre compendieux sur le calcul par la restauration et la comparaison ») a donné le mot « algèbre » (\emph{al-jabr}). Il classifiait les équations du second degré en 6 types (sans nombres négatifs ni zéro comme coefficient) et les résolvait géométriquement (complétion du carré).
    \textbf{Contexte camerounais :}
    \begin{itemize}
        \item \textbf{Agriculture :} Optimisation de la surface cultivée pour un périmètre de clôture donné (problème du terrain).
        \item \textbf{Télécoms (MTN/Orange) :} Modélisation des coûts marginaux et moyens par fonctions homographiques (forfaits data, appels).
        \item \textbf{Transport :} Calcul de rentabilité (point mort) pour un taxi ou un bus : Recette = Prix $\times$ Clients(Prix) $\to$ trinôme du second degré.
        \item \textbf{Électrification rurale :} Dimensionnement de câbles (section, chute de tension) $\to$ équations du second degré.
    \end{itemize}
\end{savaistu}

\begin{attention}[Les 5 erreurs qui coûtent des points au Probatoire]
    \begin{enumerate}
        \item \textbf{Oublier la condition de définition pour les quotients.} Ne pas écrire $cx+d \neq 0$ avant de résoudre $\frac{ax+b}{cx+d}=0$ ou l'inéquation. \textbf{Sanction :} Perte de points majeure, solution fausse si la racine annule le dénominateur.
        \item \textbf{Ne pas changer le sens de l'inégalité en multipliant/divisant par un négatif.} Ex: $-2x^2+... \ge 0 \implies x^2-... \le 0$. \textbf{Réflexe :} Écrire explicitement « Division par $-2$ (négatif), le sens change ».
        \item \textbf{Confondre « racine du numérateur » et « racine du dénominateur » dans le tableau de signes.} La racine du numérateur \textbf{est incluse} si $\le$ ou $\ge$ ; la racine du dénominateur est \textbf{JAMAIS incluse} (valeur interdite, double barre $||$).
        \item \textbf{Écrire $x_1 = \frac{-b-\sqrt{\Delta}}{2a}$ et $x_2 = \frac{-b+\sqrt{\Delta}}{2a}$ sans vérifier l'ordre.} Par convention $x_1 \le x_2$. Si $a<0$, la formule donne $x_1 > x_2$ ! \textbf{Astuce :} Calculer les deux valeurs, puis les ranger dans l'ordre croissant pour le tableau de signes.
        \item \textbf{Réponse finale incomplète ou mal rédigée.} Ensemble solution $\mathcal{S}$ non écrit, absence de phrase de conclusion, unités oubliées (m, FCFA, s), solutions « mathématiques » non filtrées par le réel (ex: longueur $-5$ m gardée).
    \end{enumerate}
\end{attention}

\begin{exercice}[Entraînement -- Équations et inéquations quotients]
    \begin{enumerate}
        \item Résoudre dans $\mathbb{R}$ : $\frac{2x-1}{x+3} = 2$.
        \item Résoudre dans $\mathbb{R}$ : $\frac{x-4}{2x+1} \le 0$.
        \item Un taxi facture une prise en charge de $500$ FCFA plus $100$ FCFA par km. Soit $x$ le nombre de km. Le coût moyen par km est $C(x) = \frac{500+100x}{x}$. Pour quelles distances le coût moyen est-il $\le 150$ FCFA/km ?
    \end{enumerate}
\end{exercice}

\begin{exercice}[Entraînement -- Trinôme du second degré]
    \begin{enumerate}
        \item Soit $P(x) = -3x^2 + 6x + 9$.
        \begin{enumerate}
            \item Mettre sous forme canonique. Donner le sommet.
            \item Calculer $\Delta$, factoriser $P(x)$.
            \item Dresser le tableau de signes.
            \item Résoudre $P(x) > 0$.
        \end{enumerate}
        \item Déterminer $m \in \mathbb{R}$ pour que l'équation $x^2 - 2mx + m+3 = 0$ ait une racine double. Donner cette racine.
    \end{enumerate}
\end{exercice}

\begin{exercice}[Entraînement -- Systèmes et Problèmes]
    \begin{enumerate}
        \item Résoudre dans $\mathbb{R}^2$ : $\begin{cases} x+y=5 \\ x^2+y^2=13 \end{cases}$.
        \item Un rectangle a un périmètre de $40$ cm et une aire de $96$ cm$^2$. Déterminer ses dimensions.
        \item Une entreprise produit $x$ unités d'un bien. Le coût total est $CT(x) = 2x^2 - 40x + 500$ (en milliers de FCFA). Le prix de vente unitaire est fixé à $120$ (milliers FCFA). Pour quelles quantités $x$ l'entreprise fait-elle un bénéfice (Recette $>$ Coût) ?
    \end{enumerate}
\end{exercice}

\begin{corrige}[Exercice 1]
    \begin{enumerate}
        \item $\frac{2x-1}{x+3}=2 \iff \frac{2x-1-2(x+3)}{x+3}=0 \iff \frac{-7}{x+3}=0$. Numérateur constant non nul $\implies \mathcal{S}=\varnothing$.
        \item $\frac{x-4}{2x+1} \le 0$. Racines : $x_N=4$, $x_D=-\frac{1}{2}$ (interdit). Tableau : $]-\infty; -\frac{1}{2}[$ : $-$ ; $]-\frac{1}{2}; 4]$ : $+$ puis $-$ ? Non. $x-4$ : $-$ puis $0$ puis $+$. $2x+1$ : $-$ ($||$) puis $+$. Quotient : $+$ ($||$) $-$ ($0$) $+$. Solution $\le 0$ : $]-\frac{1}{2}; 4]$.
        \item $C(x) \le 150 \iff \frac{500+100x}{x} \le 150 \iff \frac{500-50x}{x} \le 0 \iff \frac{10-x}{x} \le 0$ (divisé par $50>0$). Racines $10$ et $0$ (interdit). Tableau : $]-\infty;0[$ $+$ ; $]0;10]$ $-$ ; $]10;+\infty[$ $+$. Solution $\le 0$ : $]0;10]$. Contexte $x>0$ : $]0;10]$. Distance $\le 10$ km.
    \end{enumerate}
\end{corrige}

\begin{corrige}[Exercice 2]
    \begin{enumerate}
        \item $P(x)=-3x^2+6x+9 = -3(x^2-2x-3) = -3[(x-1)^2-4] = -3(x-1)^2+12$. Sommet $S(1;12)$. Max $12$.
        \item $\Delta = 36 - 4(-3)(9) = 36+108=144$. $\sqrt{\Delta}=12$. $x_1=\frac{-6-12}{-6}=3$, $x_2=\frac{-6+12}{-6}=-1$. Ordonné : $-1, 3$. $P(x)=-3(x+1)(x-3)$.
        \item $a=-3<0$. Signe $-$ en dehors, $+$ entre.
        \begin{tabular}{|c|ccccccc|}\hline $x$ & $-\infty$ & & $-1$ & & $3$ & & $+\infty$ \\ \hline $P(x)$ & $-$ & & $0$ & $+$ & $0$ & $-$ \\ \hline\end{tabular}
        \item $P(x)>0 \iff ]-1; 3[$.
        \item Racine double $\iff \Delta=0$. $\Delta' = (-m)^2 - (m+3) = m^2-m-3=0$ (ou $\Delta=4m^2-4(m+3)=0 \iff m^2-m-3=0$). $\Delta_m = 1+12=13$. $m = \frac{1\pm\sqrt{13}}{2}$. Racine double $x_0 = m = \frac{1\pm\sqrt{13}}{2}$.
    \end{enumerate}
\end{corrige}

\begin{corrige}[Exercice 3]
    \begin{enumerate}
        \item $y=5-x$. $x^2+(5-x)^2=13 \iff 2x^2-10x+12=0 \iff x^2-5x+6=0$. $\Delta=1$. $x=2$ ou $3$. Couples $(2;3)$ et $(3;2)$. $\mathcal{S}=\{(2;3);(3;2)\}$.
        \item $2(L+l)=40 \Rightarrow L+l=20$. $Ll=96$. $L,l$ racines de $X^2-20X+96=0$. $\Delta=400-384=16$. $L=12, l=8$ (ou inversement). Dimensions : $12$ cm et $8$ cm.
        \item Bénéfice $>0 \iff Rx - CT(x) > 0 \iff 120x - (2x^2-40x+500) > 0 \iff -2x^2+160x-500 > 0 \iff x^2-80x+250 < 0$ (divisé par $-2$). $\Delta=6400-1000=5400$. $\sqrt{\Delta}=30\sqrt{6}\approx 73,48$. Racines $40\pm15\sqrt{6}$. $a=1>0$, $<0$ entre les racines. $x \in ]40-15\sqrt{6}; 40+15\sqrt{6}[ \approx ]3,5; 76,5[$. Quantité entière : $x \in [4; 76]$.
    \end{enumerate}
\end{corrige}

\newpage

\coursec{Exercices}

\courssub{Je vérifie mes connaissances}

\begin{exercice}[QCM et Vrai/Faux sur les définitions]
Parmi les affirmations suivantes, une seule est correcte. Laquelle ? Justifier brièvement.
\begin{enumerate}
    \item L'équation $\dfrac{2x-1}{x+3}=0$ admet pour solution $x=-\dfrac{1}{2}$.
    \item Le trinôme $P(x)=x^2-4x+4$ s'écrit sous forme canonique $P(x)=(x-2)^2$.
    \item Le discriminant du trinôme $3x^2+5x-2$ est $\Delta = 1$.
    \item L'inéquation $\dfrac{x-2}{x+1} \ge 0$ a pour ensemble de solutions $]-\infty; -1[ \cup [2; +\infty[$.
\end{enumerate}
\end{exercice}

\begin{exercice}[Calculs directs : quotient homographique]
Résoudre dans $\mathbb{R}$ les équations et inéquations suivantes en précisant à chaque fois l'ensemble de définition.
\begin{enumerate}
    \item $\dfrac{3x+6}{x-2}=0$
    \item $\dfrac{3x+6}{x-2} \le 0$
    \item $\dfrac{5-2x}{x+4} > 0$
\end{enumerate}
\end{exercice}

\begin{exercice}[Forme canonique et discriminant]
Soit le trinôme $P(x) = 2x^2 - 8x + 5$.
\begin{enumerate}
    \item Écrire $P(x)$ sous forme canonique.
    \item Calculer son discriminant $\Delta$ et déduire le nombre de racines réelles.
    \item Donner les racines exactes si elles existent.
\end{enumerate}
\end{exercice}

\begin{exercice}[Factorisation et signe d'un trinôme]
Soit $Q(x) = -x^2 + 3x + 4$.
\begin{enumerate}
    \item Factoriser $Q(x)$ sous la forme $a(x-x_1)(x-x_2)$.
    \item Dresser le tableau de signes de $Q(x)$ sur $\mathbb{R}$.
    \item Résoudre l'inéquation $Q(x) < 0$.
\end{enumerate}
\end{exercice}

\courssub{Je m'entraîne}

\begin{exercice}[Équations et inéquations à quotient -- Cas général]
Résoudre dans $\mathbb{R}$ :
\begin{enumerate}
    \item $\dfrac{2x-7}{3x+1} = 0$
    \item $\dfrac{2x-7}{3x+1} \ge 0$
    \item $\dfrac{x^2-4}{x-1} = 0$ \quad (Attention à l'ensemble de définition)
    \item $\dfrac{x^2-4}{x-1} \le 0$
\end{enumerate}
\end{exercice}

\begin{exercice}[Étude complète d'un trinôme du second degré]
Soit la fonction $f$ définie sur $\mathbb{R}$ par $f(x) = -2x^2 + 4x + 6$.
\begin{enumerate}
    \item Écrire $f(x)$ sous forme canonique. En déduire le maximum de $f$ et la valeur de $x$ pour laquelle il est atteint.
    \item Calculer le discriminant de $f$ et factoriser $f(x)$.
    \item Dresser le tableau de variations de $f$ (ou tableau de signes de $f(x)$).
    \item Résoudre dans $\mathbb{R}$ l'inéquation $f(x) \ge 0$.
    \item Représenter graphiquement la courbe représentative de $f$ dans un repère orthonormé (indiquer le sommet, les intersections avec les axes).
\end{enumerate}
\end{exercice}

\begin{exercice}[Résolution d'équations du second degré -- Cas particuliers]
Résoudre dans $\mathbb{R}$ les équations suivantes. Factoriser le second membre lorsque c'est possible.
\begin{enumerate}
    \item $3x^2 - 5x - 2 = 0$
    \item $4x^2 - 12x + 9 = 0$
    \item $x^2 + x + 1 = 0$
    \item $2x^2 - 8 = 0$ \quad (Résoudre aussi par factorisation avec $a^2-b^2$)
\end{enumerate}
\end{exercice}

\begin{exercice}[Systèmes se ramenant au second degré -- Méthode de la somme et du produit]
Résoudre dans $\mathbb{R}^2$ les systèmes suivants.
\begin{enumerate}
    \item $\begin{cases} x+y=7 \\ xy=10 \end{cases}$
    \item $\begin{cases} 2x+y=9 \\ xy=4 \end{cases}$
    \item $\begin{cases} x-y=3 \\ x^2+y^2=17 \end{cases}$
\end{enumerate}
Vérifier que les couples trouvés sont bien solutions des systèmes.
\end{exercice}

\courssub{Je me prépare au Probatoire}

\begin{exercice}[Problème de géométrie : Le rectangle du marché Mokolo]
Un commerçant du marché Mokolo souhaite clôturer un terrain rectangulaire. Il dispose de $34$ mètres de grillage pour le périmètre et souhaite que la surface du terrain soit de $60\,\text{m}^2$.
\begin{enumerate}
    \item Soit $x$ la longueur et $y$ la largeur (en mètres). Traduire les conditions du problème par un système de deux équations à deux inconnues.
    \item Montrer que $x$ et $y$ sont les solutions d'une équation du second degré de la forme $X^2 - SX + P = 0$ où $S$ et $P$ sont des entiers à préciser.
    \item Résoudre cette équation et en déduire les dimensions du terrain (longueur et largeur). Préciser laquelle est la longueur et laquelle est la largeur.
    \item Vérifier que les dimensions trouvées conviennent bien aux conditions du problème.
\end{enumerate}
\end{exercice}

\begin{exercice}[Problème de rentabilité : Transport Douala--Yaoundé]
Un transporteur assure la liaison Douala--Yaoundé avec un car de $70$ places. Il a modélisé son bénéfice net (en milliers de francs CFA) en fonction du nombre $x$ de passagers transportés par la fonction $B(x) = -x^2 + 80x - 1200$, pour $x \in [0; 70]$.
\begin{enumerate}
    \item Calculer $B(0)$ et interpréter ce résultat dans le contexte.
    \item Étudier le signe du trinôme $B(x)$ sur $\mathbb{R}$ : forme canonique, discriminant, factorisation, tableau de signes.
    \item Déterminer l'intervalle de valeurs de $x$ pour lequel le transporteur est bénéficiaire (bénéfice strictement positif).
    \item Quel est le nombre de passagers qui maximise le bénéfice ? Calculer ce bénéfice maximal.
    \item Le transporteur envisage d'augmenter le prix du billet, ce qui modifierait la fonction de bénéfice en $B_{\text{new}}(x) = -x^2 + 80x - 1500$. Cette modification est-elle favorable ? Justifier en comparant les seuils de rentabilité.
\end{enumerate}
\end{exercice}

\begin{exercice}[Synthèse : Équation bicarrée et inéquation rationnelle]
\begin{enumerate}
    \item Résoudre dans $\mathbb{R}$ l'équation bicarrée $x^4 - 13x^2 + 36 = 0$ en effectuant le changement de variable $X = x^2$.
    \item Résoudre dans $\mathbb{R}$ l'inéquation $\dfrac{x^2 - 9}{x - 2} \ge 0$.
    \begin{itemize}
        \item Préciser l'ensemble de définition.
        \item Factoriser le numérateur.
        \item Dresser un tableau de signes à trois lignes (numérateur, dénominateur, quotient).
        \item Donner l'ensemble des solutions sous forme d'intervalles.
    \end{itemize}
    \item Soit $f(x) = \dfrac{x^2 - 9}{x - 2}$. On admet que la courbe représentative de $f$ admet une asymptote oblique d'équation $y = x + 2$. Construire le tableau de variations de $f$ sur $\mathbb{R} \setminus \{2\}$ (dérivée non demandée, utiliser le tableau de signes de $f'(x)$ si connu ou la forme canonique du numérateur de la dérivée). \textbf{Indication :} $f'(x) = \dfrac{x^2 - 4x + 9}{(x-2)^2}$.
\end{enumerate}
\end{exercice}

\newpage

\coursec{Corrigés des exercices}

\begin{corrige}[Exercice 1 — QCM et Vrai/Faux sur les définitions]
\begin{enumerate}
    \item \textbf{Faux.} $\dfrac{2x-1}{x+3}=0 \iff 2x-1=0$ et $x\neq -3$, soit $x=\dfrac{1}{2}$ et non $-\dfrac{1}{2}$.
    \item \textbf{Vrai.} $P(x)=x^2-4x+4=(x-2)^2$ (identité remarquable). La forme canonique est bien $(x-2)^2+0$.
    \item \textbf{Faux.} $\Delta = b^2-4ac = 25-4\times 3\times(-2)=25+24=49\neq 1$.
    \item \textbf{Vrai.} Le numérateur s'annule en $2$, le dénominateur en $-1$ (valeur interdite). Le quotient est positif à l'extérieur des racines : $S=]-\infty;-1[\cup[2;+\infty[$, avec $-1$ exclu (barre verticale double) et $2$ inclus.
\end{enumerate}
\textbf{Conclusion :} les affirmations 2 et 4 sont correctes ; si une seule réponse est attendue, c'est la \textbf{2} (la 4 est également exacte — vérifier l'énoncé).
\end{corrige}

\begin{corrige}[Exercice 2 — Calculs directs : quotient homographique]
\begin{enumerate}
    \item \textbf{Ensemble de définition :} $x-2\neq 0 \iff x\neq 2$, donc $D=\mathbb{R}\setminus\{2\}$.\\
    $\dfrac{3x+6}{x-2}=0 \iff 3x+6=0 \iff x=-2$. Comme $-2\in D$ : $S=\{-2\}$.
    \item \textbf{Tableau de signes} sur $\mathbb{R}\setminus\{2\}$ : $3x+6=0 \iff x=-2$ ; $x-2=0 \iff x=2$.
    \begin{center}
    \begin{tabular}{|c|ccccccc|}\hline
    $x$ & $-\infty$ & & $-2$ & & $2$ & & $+\infty$ \\ \hline
    $3x+6$ & & $-$ & $0$ & $+$ & & $+$ & \\ \hline
    $x-2$ & & $-$ & & $-$ & $0$ & $+$ & \\ \hline
    $\dfrac{3x+6}{x-2}$ & & $+$ & $0$ & $-$ & $\|$ & $+$ & \\ \hline
    \end{tabular}
    \end{center}
    Le quotient est $\le 0$ sur $[-2;2[$ ($2$ exclu, valeur interdite). $S=[-2;2[$.
    \item $D=\mathbb{R}\setminus\{-4\}$. $5-2x=0 \iff x=\dfrac{5}{2}$ ; $x+4=0 \iff x=-4$. Le quotient est strictement positif entre les racines : $S=\left]-4;\dfrac{5}{2}\right[$ (bornes exclues car inégalité stricte et $-4$ interdit).
\end{enumerate}
\textbf{Point de vigilance :} une valeur qui annule le dénominateur est toujours exclue, même avec une inégalité large.
\end{corrige}

\begin{corrige}[Exercice 3 — Forme canonique et discriminant]
$P(x)=2x^2-8x+5$.
\begin{enumerate}
    \item $P(x)=2\left(x^2-4x\right)+5=2\left[(x-2)^2-4\right]+5=2(x-2)^2-8+5=2(x-2)^2-3$.\\
    Forme canonique : $P(x)=2(x-2)^2-3$.
    \item $\Delta=b^2-4ac=(-8)^2-4\times 2\times 5=64-40=24$. Comme $\Delta>0$, $P$ admet \textbf{deux racines réelles distinctes}.
    \item $x_1=\dfrac{8-\sqrt{24}}{4}=\dfrac{8-2\sqrt{6}}{4}=\dfrac{4-\sqrt{6}}{2}$ et $x_2=\dfrac{4+\sqrt{6}}{2}$.
\end{enumerate}
\textbf{Vérification :} $x_1+x_2=4=-\dfrac{b}{a}$ et $x_1x_2=\dfrac{16-6}{4}=\dfrac{10}{4}=\dfrac{5}{2}=\dfrac{c}{a}$. Correct.
\end{corrige}

\begin{corrige}[Exercice 4 — Factorisation et signe d'un trinôme]
$Q(x)=-x^2+3x+4$.
\begin{enumerate}
    \item $\Delta=9-4\times(-1)\times 4=9+16=25>0$. Racines : $x_1=\dfrac{-3-5}{-2}=4$ et $x_2=\dfrac{-3+5}{-2}=-1$.\\
    Factorisation : $Q(x)=-(x-4)(x+1)$. \textbf{Vérification :} $-(x^2-3x-4)=-x^2+3x+4$. Correct.
    \item Tableau de signes ($a=-1<0$, le trinôme est positif entre les racines) :
    \begin{center}
    \begin{tabular}{|c|ccccccc|}\hline
    $x$ & $-\infty$ & & $-1$ & & $4$ & & $+\infty$ \\ \hline
    $Q(x)$ & & $-$ & $0$ & $+$ & $0$ & $-$ & \\ \hline
    \end{tabular}
    \end{center}
    \item $Q(x)<0 \iff x\in\,]-\infty;-1[\cup]4;+\infty[$.
\end{enumerate}
\textbf{Point de vigilance :} ne pas oublier le facteur $a=-1$ dans la factorisation ; écrire $(x-4)(x+1)$ seul serait faux.
\end{corrige}

\begin{corrige}[Exercice 5 — Équations et inéquations à quotient : cas général]
\begin{enumerate}
    \item $D=\mathbb{R}\setminus\left\{-\dfrac{1}{3}\right\}$. $\dfrac{2x-7}{3x+1}=0 \iff x=\dfrac{7}{2}$. $S=\left\{\dfrac{7}{2}\right\}$.
    \item Zéros : $x=\dfrac{7}{2}$ (numérateur) et $x=-\dfrac{1}{3}$ (dénominateur, interdit). Le quotient est du signe de $(2x-7)(3x+1)$, positif à l'extérieur des racines :
    $S=\left]-\infty;-\dfrac{1}{3}\right[\cup\left[\dfrac{7}{2};+\infty\right[$.
    \item $D=\mathbb{R}\setminus\{1\}$. $x^2-4=0 \iff x=2$ ou $x=-2$, toutes deux dans $D$. $S=\{-2;2\}$.
    \item Signes de $x^2-4=(x-2)(x+2)$ et de $x-1$ :
    \begin{center}
    \begin{tabular}{|c|ccccccccc|}\hline
    $x$ & $-\infty$ & & $-2$ & & $1$ & & $2$ & & $+\infty$ \\ \hline
    $x^2-4$ & & $+$ & $0$ & $-$ & & $-$ & $0$ & $+$ & \\ \hline
    $x-1$ & & $-$ & & $-$ & $0$ & $+$ & & $+$ & \\ \hline
    Quotient & & $-$ & $0$ & $+$ & $\|$ & $-$ & $0$ & $+$ & \\ \hline
    \end{tabular}
    \end{center}
    $S=]-\infty;-2]\cup]1;2]$ ($1$ exclu, $-2$ et $2$ inclus).
\end{enumerate}
\end{corrige}

\begin{corrige}[Exercice 6 — Étude complète d'un trinôme du second degré]
$f(x)=-2x^2+4x+6$.
\begin{enumerate}
    \item $f(x)=-2\left(x^2-2x\right)+6=-2\left[(x-1)^2-1\right]+6=-2(x-1)^2+8$.\\
    Comme $a=-2<0$, $f$ admet un \textbf{maximum} égal à $8$, atteint en $x=1$.
    \item $\Delta=16-4\times(-2)\times 6=16+48=64>0$. Racines : $x_1=\dfrac{-4-8}{-4}=3$, $x_2=\dfrac{-4+8}{-4}=-1$.\\
    Factorisation : $f(x)=-2(x-3)(x+1)$.
    \item Variations : $f$ est croissante sur $]-\infty;1]$, décroissante sur $[1;+\infty[$, maximum $f(1)=8$. Signes : $f(x)<0$ sur $]-\infty;-1[\cup]3;+\infty[$, $f(x)>0$ sur $]-1;3[$.
    \item $f(x)\ge 0 \iff x\in[-1;3]$.
    \item Éléments du graphique : sommet $S(1;8)$ ; intersections avec l'axe des abscisses : $A(-1;0)$ et $B(3;0)$ ; intersection avec l'axe des ordonnées : $C(0;6)$ ; axe de symétrie $x=1$. La parabole est tournée vers le bas.
    \begin{popfigure}
    \begin{tikzpicture}
    \begin{axis}[width=0.7\linewidth,axis lines=middle,xlabel=$x$,ylabel=$y$,xmin=-3,xmax=5,ymin=-6,ymax=10,xtick={-2,-1,1,2,3,4},ytick={-4,2,4,6,8},grid=both,grid style={popline!40}]
    \addplot[popcurve,domain=-2.2:4.2,samples=200]{-2*x^2+4*x+6};
    \addplot[mark=*,only marks,poporange] coordinates {(1,8) (-1,0) (3,0) (0,6)};
    \end{axis}
    \end{tikzpicture}
    \legende{Courbe de $f(x)=-2x^2+4x+6$ : sommet $(1;8)$, zéros en $-1$ et $3$.}
    \end{popfigure}
\end{enumerate}
\end{corrige}

\begin{corrige}[Exercice 7 — Résolution d'équations du second degré : cas particuliers]
\begin{enumerate}
    \item $3x^2-5x-2=0$ : $\Delta=25+24=49$. $x_1=\dfrac{5-7}{6}=-\dfrac{1}{3}$, $x_2=\dfrac{5+7}{6}=2$. $S=\left\{-\dfrac{1}{3};2\right\}$. Factorisation : $3x^2-5x-2=3\left(x+\dfrac{1}{3}\right)(x-2)=(3x+1)(x-2)$.
    \item $4x^2-12x+9=0$ : $\Delta=144-144=0$, racine double $x_0=\dfrac{12}{8}=\dfrac{3}{2}$. On reconnaît $(2x-3)^2=0$. $S=\left\{\dfrac{3}{2}\right\}$.
    \item $x^2+x+1=0$ : $\Delta=1-4=-3<0$. Pas de racine réelle : $S=\varnothing$.
    \item $2x^2-8=0 \iff x^2=4 \iff x=\pm 2$. Par factorisation : $2(x^2-4)=2(x-2)(x+2)=0$. $S=\{-2;2\}$.
\end{enumerate}
\end{corrige}

\begin{corrige}[Exercice 8 — Systèmes se ramenant au second degré]
\begin{enumerate}
    \item $x$ et $y$ ont pour somme $S=7$ et produit $P=10$ : ce sont les racines de $X^2-7X+10=0$. $\Delta=49-40=9$ ; $X_1=2$, $X_2=5$.\\
    $S=\{(2;5);(5;2)\}$. \textbf{Vérification :} $2+5=7$, $2\times 5=10$. Correct.
    \item De la première équation : $y=9-2x$. Substitution : $x(9-2x)=4 \iff -2x^2+9x-4=0 \iff 2x^2-9x+4=0$.\\
    $\Delta=81-32=49$ ; $x_1=\dfrac{9-7}{4}=\dfrac{1}{2}$, $x_2=\dfrac{9+7}{4}=4$.\\
    Si $x=\dfrac{1}{2}$ : $y=9-1=8$. Si $x=4$ : $y=9-8=1$.\\
    $S=\left\{\left(\dfrac{1}{2};8\right);(4;1)\right\}$. \textbf{Vérification :} $\dfrac{1}{2}\times 8=4$ et $2\times\dfrac{1}{2}+8=9$ ; $4\times 1=4$ et $8+1=9$. Correct.
    \item $y=x-3$. Alors $x^2+(x-3)^2=17 \iff 2x^2-6x+9=17 \iff 2x^2-6x-8=0 \iff x^2-3x-4=0$.\\
    $\Delta=9+16=25$ ; $x_1=-1$, $x_2=4$.\\
    Si $x=-1$ : $y=-4$. Si $x=4$ : $y=1$.\\
    $S=\{(-1;-4);(4;1)\}$. \textbf{Vérification :} $(-1)-(-4)=3$ et $1+16=17$ ; $4-1=3$ et $16+1=17$. Correct.
\end{enumerate}
\textbf{Point de vigilance :} chaque valeur de $x$ donne une seule valeur de $y$ ; ne pas mélanger les couples.
\end{corrige}

\begin{corrige}[Exercice 9 — Problème de géométrie : le rectangle du marché Mokolo]
\begin{enumerate}
    \item Périmètre : $2(x+y)=34$, soit $x+y=17$. Aire : $xy=60$. Système : $\begin{cases} x+y=17 \\ xy=60 \end{cases}$ avec $x>0$, $y>0$.
    \item $x$ et $y$ ont pour somme $S=17$ et produit $P=60$ : ce sont les racines de $X^2-17X+60=0$.
    \item $\Delta=289-240=49$ ; $X_1=\dfrac{17-7}{2}=5$, $X_2=\dfrac{17+7}{2}=12$.\\
    La longueur étant la plus grande dimension : longueur $=12$ m, largeur $=5$ m.
    \item \textbf{Vérification :} périmètre $=2(12+5)=34$ m (le grillage suffit exactement) ; aire $=12\times 5=60\,\text{m}^2$. Les conditions sont satisfaites.
\end{enumerate}
\textbf{Barème indicatif (4 pts) :} système correct (1 pt) ; équation $X^2-SX+P=0$ justifiée (1 pt) ; résolution avec $\Delta$ (1 pt) ; conclusion avec distinction longueur/largeur et vérification (1 pt).\\
\textbf{Points de vigilance :} exiger $x,y>0$ ; citer le théorème somme-produit ; la vérification finale est attendue dans un problème concret.
\end{corrige}

\begin{corrige}[Exercice 10 — Problème de rentabilité : transport Douala--Yaoundé]
$B(x)=-x^2+80x-1200$ sur $[0;70]$.
\begin{enumerate}
    \item $B(0)=-1200$ : sans passager, le transporteur perd \num{1200000} FCFA (charges fixes : carburant, péages, salaires).
    \item Forme canonique : $B(x)=-\left(x^2-80x\right)-1200=-\left[(x-40)^2-1600\right]-1200=-(x-40)^2+400$.\\
    $\Delta=6400-4800=1600>0$ ; racines : $x_1=\dfrac{-80-40}{-2}=60$, $x_2=\dfrac{-80+40}{-2}=20$.\\
    Factorisation : $B(x)=-(x-20)(x-60)$. Comme $a<0$, $B(x)>0$ entre les racines :
    \begin{center}
    \begin{tabular}{|c|ccccccc|}\hline
    $x$ & $-\infty$ & & $20$ & & $60$ & & $+\infty$ \\ \hline
    $B(x)$ & & $-$ & $0$ & $+$ & $0$ & $-$ & \\ \hline
    \end{tabular}
    \end{center}
    \item Bénéfice strictement positif : $x\in\,]20;60[$. Le transporteur est bénéficiaire pour un nombre de passagers compris entre $21$ et $59$ (entiers).
    \item Maximum atteint en $x=40$ (forme canonique) : $B(40)=400$, soit un bénéfice maximal de \num{400000} FCFA pour $40$ passagers.
    \item $B_{\text{new}}(x)=-x^2+80x-1500$ : $\Delta=6400-6000=400$ ; racines $x_1=\dfrac{-80-20}{-2}=50$, $x_2=\dfrac{-80+20}{-2}=30$. La zone de rentabilité devient $]30;50[$, plus étroite que $]20;60[$, et le maximum devient $B_{\text{new}}(40)=-(40)^2+3200-1500=100$ (soit \num{100000} FCFA au lieu de \num{400000}). La modification est \textbf{défavorable} : les seuils de rentabilité se resserrent et le bénéfice maximal diminue.
\end{enumerate}
\textbf{Barème indicatif (5 pts) :} $B(0)$ interprété (0,5 pt) ; forme canonique (1 pt) ; discriminant, racines, factorisation (1 pt) ; tableau de signes et intervalle de rentabilité (1 pt) ; maximum (0,5 pt) ; comparaison argumentée (1 pt).\\
\textbf{Points de vigilance :} interpréter les résultats dans le contexte (unités : milliers de FCFA) ; pour la question 5, comparer à la fois les racines et les maxima.
\end{corrige}

\begin{corrige}[Exercice 11 — Synthèse : équation bicarrée et inéquation rationnelle]
\begin{enumerate}
    \item Posons $X=x^2$ ($X\ge 0$). L'équation devient $X^2-13X+36=0$. $\Delta=169-144=25$ ; $X_1=\dfrac{13-5}{2}=4$, $X_2=\dfrac{13+5}{2}=9$, toutes deux positives.\\
    $x^2=4 \iff x=\pm 2$ ; $x^2=9 \iff x=\pm 3$. $S=\{-3;-2;2;3\}$.\\
    \textbf{Vérification :} $3^4-13\times 9+36=81-117+36=0$. Correct.
    \item \textbf{Ensemble de définition :} $x\neq 2$, $D=\mathbb{R}\setminus\{2\}$.\\
    Numérateur : $x^2-9=(x-3)(x+3)$, zéros en $-3$ et $3$. Dénominateur nul en $2$.
    \begin{center}
    \begin{tabular}{|c|ccccccccc|}\hline
    $x$ & $-\infty$ & & $-3$ & & $2$ & & $3$ & & $+\infty$ \\ \hline
    $x^2-9$ & & $+$ & $0$ & $-$ & & $-$ & $0$ & $+$ & \\ \hline
    $x-2$ & & $-$ & & $-$ & $0$ & $+$ & & $+$ & \\ \hline
    $\dfrac{x^2-9}{x-2}$ & & $-$ & $0$ & $+$ & $\|$ & $-$ & $0$ & $+$ & \\ \hline
    \end{tabular}
    \end{center}
    $S=[-3;2[\cup[3;+\infty[$ ($-3$ et $3$ inclus, $2$ exclu).
    \item On admet $f'(x)=\dfrac{x^2-4x+9}{(x-2)^2}$. Étudions le signe du numérateur : $\Delta'=16-36=-20<0$ et coefficient dominant positif, donc $x^2-4x+9>0$ pour tout $x$ réel. Comme $(x-2)^2>0$ sur $D$, on a $f'(x)>0$ sur $]-\infty;2[$ et sur $]2;+\infty[$ : $f$ est strictement croissante sur chacun de ces intervalles.\\
    Limites : $\lim_{x\to 2^-}f(x)=+\infty$ (numérateur $\to -5<0$, dénominateur $\to 0^-$... en fait $x^2-9\to -5$ et $x-2\to 0^-$ donc $f\to +\infty$) et $\lim_{x\to 2^+}f(x)=-\infty$. De plus $f(x)-(x+2)=\dfrac{x^2-9-(x+2)(x-2)}{x-2}=\dfrac{-5}{x-2}\to 0$ en $\pm\infty$ : asymptote oblique $y=x+2$ confirmée.
    \begin{center}
    \begin{tabular}{|c|ccccc|}\hline
    $x$ & $-\infty$ & & $2$ & & $+\infty$ \\ \hline
    $f'(x)$ & & $+$ & $\|$ & $+$ & \\ \hline
    $f$ & $-\infty$ & $\nearrow$ & $\|\ +\infty \quad -\infty\ \|$ & $\nearrow$ & $+\infty$ \\ \hline
    \end{tabular}
    \end{center}
    (Précisément : $f$ croît de $-\infty$ vers $+\infty$ sur $]-\infty;2[$, puis de $-\infty$ vers $+\infty$ sur $]2;+\infty[$.)
\end{enumerate}
\textbf{Barème indicatif (5 pts) :} changement de variable et résolution (1,5 pt) ; ensemble de définition et factorisation (1 pt) ; tableau de signes à trois lignes (1,5 pt) ; tableau de variations avec justification du signe de $f'$ (1 pt).\\
\textbf{Points de vigilance :} pour la bicarrée, ne retenir que les valeurs $X\ge 0$ ; dans le tableau de signes, double barre en $x=2$ ; pour la question 3, justifier le signe du numérateur de $f'$ par son discriminant négatif.
\end{corrige}

\newpage

\coursec{Fiche bilan}

\begin{popfigure}
\begin{tikzpicture}[mindmap, grow cyclic, every node/.style={concept, font=\small},
  root concept/.append style={concept color=popblue, font=\bfseries\small},
  level 1/.append style={level distance=3.4cm, sibling angle=60},
  level 2/.append style={level distance=2.1cm, sibling angle=45, font=\scriptsize}]
\node[root concept]{Équations, inéquations, systèmes}
  child[concept color=poporange]{ node {Quotient $\frac{ax+b}{cx+d}$}
    child { node {Valeur interdite $x=-\frac{d}{c}$} }
    child { node {Tableau de signes} }
  }
  child[concept color=popgreen]{ node {Forme canonique}
    child { node {$a(x-\alpha)^2+\beta$} }
    child { node {$\alpha=-\frac{b}{2a}$ ; $\beta=f(\alpha)$} }
  }
  child[concept color=poppurple]{ node {Discriminant}
    child { node {$\Delta=b^2-4ac$} }
    child { node {Signe de $\Delta$ $\Rightarrow$ nombre de racines} }
  }
  child[concept color=poppink]{ node {Factorisation et signe}
    child { node {$a(x-x_1)(x-x_2)$} }
    child { node {Signe de $a$ hors racines} }
  }
  child[concept color=popgold]{ node {Systèmes dans $\mathbb{R}^2$}
    child { node {Substitution} }
    child { node {Somme et produit} }
  }
  child[concept color=popblue!70]{ node {Problèmes concrets}
    child { node {Mise en équation} }
    child { node {Vérification des solutions} }
  };
\end{tikzpicture}
\legende{Carte mentale de la leçon : les six compétences clés autour des équations, inéquations et systèmes.}
\end{popfigure}

\begin{bilan}
\textbf{1. Quotient $\dfrac{ax+b}{cx+d}$ (fonction homographique)}
\begin{itemize}
  \item \cle{Valeur interdite} : $x=-\dfrac{d}{c}$ (celle qui annule le dénominateur). On la signale par une double barre dans le tableau de signes.
  \item Équation $\dfrac{ax+b}{cx+d}=0$ : un quotient est nul si et seulement si son \cle{numérateur est nul} (et le dénominateur non nul) : $x=-\dfrac{b}{a}$, à condition que ce ne soit pas la valeur interdite.
  \item Inéquation $\dfrac{ax+b}{cx+d}<0$ (ou $\leq 0$, $>0$, $\geq 0$) : on dresse le \cle{tableau de signes} du numérateur et du dénominateur, puis on lit le signe du quotient. Jamais de produit en croix avec une inéquation !
\end{itemize}

\textbf{2. Forme canonique de $f(x)=ax^2+bx+c$ ($a\neq 0$)}
\[ f(x)=a(x-\alpha)^2+\beta \quad \text{avec} \quad \alpha=-\frac{b}{2a} \quad \text{et} \quad \beta=f(\alpha)=-\frac{\Delta}{4a}. \]
Méthode express : factoriser $a$ devant les termes en $x$, puis reconnaître le début d'une identité remarquable.

\textbf{3. Discriminant et équations du second degré}
\begin{center}
\begin{tabularx}{\linewidth}{|l|X|X|X|}\hline
\rowcolor{popblueL} Signe de $\Delta=b^2-4ac$ & Racines de $ax^2+bx+c=0$ & Factorisation & Signe du trinôme \\ \hline
$\Delta>0$ & $x_1=\dfrac{-b-\sqrt{\Delta}}{2a}$, $x_2=\dfrac{-b+\sqrt{\Delta}}{2a}$ & $a(x-x_1)(x-x_2)$ & signe de $a$ à l'extérieur des racines, signe de $-a$ entre \\ \hline
$\Delta=0$ & racine double $x_0=-\dfrac{b}{2a}$ & $a(x-x_0)^2$ & signe de $a$, nul en $x_0$ \\ \hline
$\Delta<0$ & aucune racine réelle & pas de factorisation dans $\mathbb{R}$ & signe de $a$ partout \\ \hline
\end{tabularx}
\end{center}
Relations utiles : $x_1+x_2=-\dfrac{b}{a}$ et $x_1 x_2=\dfrac{c}{a}$.

\textbf{4. Inéquations du second degré}
\begin{enumerate}
  \item Tout ramener à la forme $ax^2+bx+c<0$ (ou $\leq$, $>$, $\geq$).
  \item Calculer $\Delta$, trouver les racines éventuelles.
  \item Dresser le tableau de signes du trinôme.
  \item Lire l'intervalle solution selon l'inégalité demandée (bornes incluses si l'inégalité est large).
\end{enumerate}

\textbf{5. Systèmes se ramenant au second degré}
\begin{itemize}
  \item \cle{Substitution} : exprimer une inconnue grâce à l'équation du premier degré, reporter dans l'autre $\Rightarrow$ équation du second degré.
  \item \cle{Somme et produit} : si $x+y=S$ et $xy=P$, alors $x$ et $y$ sont les racines de $X^2-SX+P=0$ (solutions réelles seulement si $S^2-4P\geq 0$).
  \item Toujours \cle{vérifier} chaque couple $(x\,;y)$ dans les deux équations de départ.
\end{itemize}

\textbf{6. Résolution de problèmes}
Méthode en 4 temps : choisir les inconnues $\to$ traduire les données en équations/inéquations $\to$ résoudre $\to$ confronter les solutions au contexte (longueur positive, prix en FCFA, effectif entier\ldots) et conclure par une phrase.
\end{bilan}

\begin{aretenir}[Checklist avant le Probatoire]
\begin{itemize}
  \item[$\square$] Je sais repérer et exclure la valeur interdite d'un quotient $\dfrac{ax+b}{cx+d}$.
  \item[$\square$] Je sais résoudre $\dfrac{ax+b}{cx+d}=0$ en annulant uniquement le numérateur.
  \item[$\square$] Je sais dresser le tableau de signes d'un quotient et résoudre une inéquation quotient, sans jamais multiplier en croix.
  \item[$\square$] Je sais mettre $ax^2+bx+c$ sous forme canonique $a(x-\alpha)^2+\beta$ avec $\alpha=-\dfrac{b}{2a}$.
  \item[$\square$] Je connais par c\oe ur $\Delta=b^2-4ac$ et les formules $x_1=\dfrac{-b-\sqrt{\Delta}}{2a}$, $x_2=\dfrac{-b+\sqrt{\Delta}}{2a}$.
  \item[$\square$] Je sais factoriser un trinôme lorsque $\Delta\geq 0$ et je sais qu'il est factorisable seulement dans ce cas.
  \item[$\square$] Je sais donner le signe d'un trinôme : signe de $a$ à l'extérieur des racines, signe contraire entre les racines.
  \item[$\square$] Je sais résoudre une inéquation du second degré grâce au tableau de signes, avec les bonnes bornes (ouvertes ou fermées).
  \item[$\square$] Je sais résoudre un système par substitution ou par somme et produit ($X^2-SX+P=0$).
  \item[$\square$] Je pense à vérifier mes solutions (couples, racines) dans les équations de départ.
  \item[$\square$] Je sais mettre un problème concret en équation et rejeter les solutions incompatibles avec le contexte.
  \item[$\square$] Je rédige proprement : ensemble de définition, calcul de $\Delta$, conclusion sous forme d'intervalle ou d'ensemble $S=\{\ldots\}$.
\end{itemize}
\end{aretenir}

\findecours

\end{document}
